% Simple and Damped Harmonic motion — Further Mathematics Upper Sixth — Cours signé OnBuch+
% Official MINESEC syllabus. Compiler avec : tectonic FMA09-simple-and-damped-harmonic-motion.tex
\newcommand{\DOCMATIERE}{Further Mathematics}
\newcommand{\DOCNIVEAU}{Upper Sixth}
\newcommand{\DOCMODULE}{Upper Sixth — Further mathematics}
\newcommand{\DOCLECON}{9}
\newcommand{\DOCTITRE}{Simple and Damped Harmonic motion}
% OnBuch+ — préambule « cours signés » ENRICHI (compilé avec tectonic / XeLaTeX)
% Système visuel « Pop » d'OnBuch+ : fond crème (jamais blanc), bordures encre
% épaisses, ombres « dures » décalées, titres Archivo Black en capitales.
% Version MATHÉMATIQUES : graphes (pgfplots/tikz), boîtes pédagogiques
% (définition, propriété, méthode, attention, le savais-tu, exercice résolu,
% exercice, TP, bilan), page de garde et signature OnBuch+ ancrée en bas.
%
% Macros attendues dans main.tex AVANT \input{preamble} :
%   \DOCMATIERE  \DOCNIVEAU  \DOCMODULE  \DOCLECON (numéro)  \DOCTITRE
\documentclass[11pt]{article}

\usepackage{fontspec}
\usepackage[english]{babel}
\usepackage[a4paper,margin=2cm,top=3.4cm,bottom=2.8cm,headheight=1.4cm,footskip=1.3cm]{geometry}
\usepackage{amsmath,amssymb,amsfonts}
\usepackage{mathtools} % \xmapsto, \coloneqq... souvent utilisés par les modèles
\usepackage{cancel} % simplifications barrées (bilans de forces)
% \abs{z} (module, valeur absolue) et \norm{u} (norme), utilisés par les modèles
\providecommand{\abs}{}\let\abs\relax\DeclarePairedDelimiter{\abs}{\lvert}{\rvert}
\providecommand{\norm}{}\let\norm\relax\DeclarePairedDelimiter{\norm}{\lVert}{\rVert}
\usepackage[table]{xcolor} % [table] : fournit aussi \rowcolors (alternance de lignes)
\usepackage{pagecolor}
\usepackage{tikz}
\usetikzlibrary{arrows.meta,positioning,calc,shapes.geometric,shapes.misc,shapes.arrows,decorations.pathmorphing,decorations.pathreplacing,decorations.markings,patterns,shadows,mindmap,trees,fit,backgrounds,matrix,intersections,angles,quotes}
\usepackage{pgfplots}
\usepackage[version=4]{mhchem} % équations nucléaires \ce{^{235}_{92}U}
\usepackage[european,siunitx]{circuitikz} % schémas électriques
\pgfplotsset{compat=1.18}
\usepgfplotslibrary{fillbetween,groupplots} % aires entre courbes (calcul intégral)
\usepackage{enumitem}
\usepackage{fancyhdr}
% [most] charge toutes les bibliothèques tcolorbox (listings, minted, raster,
% documentation...) et gonfle inutilement la mémoire TeX sur un document long
% (~300 boîtes) : on ne charge que ce qui est réellement utilisé.
\usepackage[skins,breakable]{tcolorbox}
\usepackage{graphicx}
\usepackage{colortbl}
\usepackage{booktabs}
\usepackage{tabularx}
\usepackage{diagbox} % case d'en-tête diagonale des tableaux à double entrée
% Colonnes extensibles centrée / gauche / droite (C, L, R), souvent utilisées par les modèles
\newcolumntype{C}{>{\centering\arraybackslash}X}
\newcolumntype{L}{>{\raggedright\arraybackslash}X}
\newcolumntype{R}{>{\raggedleft\arraybackslash}X}
\usepackage{array}
\usepackage{multicol}
\usepackage{siunitx}
% parse-numbers=false : accepte aussi \SI{2,5 \times 10^{-3}}{...}
\sisetup{locale=UK,output-decimal-marker={.},group-minimum-digits=4,parse-numbers=false}
\DeclareSIUnit\ohm{\ensuremath{\symup{\Omega}}} % U+2126 absent de la police
\DeclareSIUnit\division{div} % réglages d'oscilloscope (ms/div, V/div)
\DeclareSIUnit\tr{tr} % tours (vitesse de rotation, ex. \SI{1500}{\tr\per\minute})
\usepackage{qrcode}
% \degree : commande fréquemment utilisée par les modèles pour un angle ou une
% température en dehors de \SI{}{} (non fournie par les paquets chargés).
\providecommand{\degree}{^{\circ}}
% \qed (fin de démonstration), utilisé par les modèles sans amsthm
\providecommand{\qed}{\hfill$\square$}
% \barre : double barre des valeurs interdites dans un tableau de variations (tkz-tab)
\providecommand{\barre}{\ensuremath{\|}}
% environnement proof (amsthm non chargé) utilisé par les modèles
\ifdefined\proof\else\newenvironment{proof}[1][Proof]{\par\noindent\textit{#1.}\ }{\qed\par}\fi
\usepackage{lastpage}
\usepackage{hyperref}
\hypersetup{hidelinks}

% ============================================================
% Palette « Pop » OnBuch+ — mêmes hex que la classe P (plus_ui.dart)
% ============================================================
\definecolor{poporange}{HTML}{F59321}
\definecolor{popdark}{HTML}{E07A0C}
\definecolor{popcream}{HTML}{FFF6E8}
\definecolor{popink}{HTML}{1C1714}
\definecolor{poppurple}{HTML}{7A5AE0}
\definecolor{popgreen}{HTML}{2E9E6A}
\definecolor{poppink}{HTML}{E0457B}
\definecolor{popblue}{HTML}{2F7DE1}
\definecolor{popgold}{HTML}{C98A12}
\definecolor{popmuted}{HTML}{8A7F76}
\definecolor{popline}{HTML}{EADFD2}
\definecolor{popgray}{HTML}{8A7F76}
\colorlet{popgrey}{popgray}
% Alias tolérants (le modèle nomme parfois les couleurs différemment)
\colorlet{popwhite}{white}
\colorlet{popblack}{popink}
\colorlet{popred}{poppink}
\colorlet{popyellow}{popgold}
% Teintes claires pour fonds de tableaux / zones
\colorlet{popblueL}{popblue!10!white}
\colorlet{poporangeL}{poporange!14!white}
\colorlet{popgreenL}{popgreen!12!white}
\colorlet{poppurpleL}{poppurple!12!white}
\colorlet{poppinkL}{poppink!12!white}
\colorlet{popgoldL}{popgold!14!white}

\pagecolor{popcream}

% ============================================================
% Typographie
% ============================================================
\defaultfontfeatures{Ligatures=TeX}
\setmainfont{PlusJakartaSans-Medium.ttf}[Path=../../../fonts/,BoldFont=PlusJakartaSans-Bold.ttf]
% Maths en sans-serif (Fira Math) avec chiffres et lettres droites de Plus
% Jakarta, pour que les formules se fondent dans le texte.
\usepackage[math-style=ISO,bold-style=ISO]{unicode-math}
\setmathfont{FiraMath-Regular.otf}[Path=../../../fonts/]
\setmathfont{PlusJakartaSans-Medium.ttf}[Path=../../../fonts/,range={up/num,up/{Latin,latin},\mathup}]
% (icomma retiré : décimales avec point en anglais)
% Caractères mathématiques Unicode absents des polices de texte (Plus Jakarta,
% Archivo Black) : rendus par leur équivalent mathématique, en texte comme en
% formule — sinon ils s'affichent en carré vide (ex. « Arithmétique dans ℤ »).
\usepackage{newunicodechar}
\newunicodechar{ℝ}{\ensuremath{\mathbb{R}}}
\newunicodechar{ℤ}{\ensuremath{\mathbb{Z}}}
\newunicodechar{ℂ}{\ensuremath{\mathbb{C}}}
\newunicodechar{ℕ}{\ensuremath{\mathbb{N}}}
\newunicodechar{ℚ}{\ensuremath{\mathbb{Q}}}
\newunicodechar{𝔻}{\ensuremath{\mathbb{D}}}
\newunicodechar{α}{\ensuremath{\alpha}}
\newunicodechar{β}{\ensuremath{\beta}}
\newunicodechar{γ}{\ensuremath{\gamma}}
\newunicodechar{δ}{\ensuremath{\delta}}
\newunicodechar{ε}{\ensuremath{\varepsilon}}
\newunicodechar{θ}{\ensuremath{\theta}}
\newunicodechar{λ}{\ensuremath{\lambda}}
\newunicodechar{μ}{\ensuremath{\mu}}
\newunicodechar{σ}{\ensuremath{\sigma}}
\newunicodechar{τ}{\ensuremath{\tau}}
\newunicodechar{φ}{\ensuremath{\varphi}}
\newunicodechar{ψ}{\ensuremath{\psi}}
\newunicodechar{ω}{\ensuremath{\omega}}
\newunicodechar{Σ}{\ensuremath{\Sigma}}
% majuscules produites par \MakeUppercase dans les titres (α-aminés -> Α)
\newunicodechar{Α}{\ensuremath{\alpha}}
\newunicodechar{Β}{\ensuremath{\beta}}
\newunicodechar{Γ}{\ensuremath{\Gamma}}
\newunicodechar{Δ}{\ensuremath{\Delta}}
\newunicodechar{Ε}{\ensuremath{\varepsilon}}
\newunicodechar{Θ}{\ensuremath{\Theta}}
\newunicodechar{Λ}{\ensuremath{\Lambda}}
\newunicodechar{Μ}{\ensuremath{\mu}}
\newunicodechar{Τ}{\ensuremath{\tau}}
\newunicodechar{Φ}{\ensuremath{\Phi}}
\newunicodechar{Ψ}{\ensuremath{\Psi}}
\newunicodechar{Ω}{\ensuremath{\Omega}}
\newunicodechar{↦}{\ensuremath{\mapsto}}
\newunicodechar{∈}{\ensuremath{\in}}
\newunicodechar{∉}{\ensuremath{\notin}}
\newunicodechar{∧}{\ensuremath{\wedge}}
\newunicodechar{⇔}{\ensuremath{\Leftrightarrow}}
\newunicodechar{⟺}{\ensuremath{\Longleftrightarrow}}
\newunicodechar{⇒}{\ensuremath{\Rightarrow}}
\newunicodechar{⟹}{\ensuremath{\Longrightarrow}}
\newunicodechar{∘}{\ensuremath{\circ}}
\newunicodechar{∪}{\ensuremath{\cup}}
\newunicodechar{∩}{\ensuremath{\cap}}
\newunicodechar{≡}{\ensuremath{\equiv}}
\newunicodechar{⊂}{\ensuremath{\subset}}
\newunicodechar{⊥}{\ensuremath{\perp}}
\newunicodechar{∃}{\ensuremath{\exists}}
\newunicodechar{∀}{\ensuremath{\forall}}
\newunicodechar{∛}{\ensuremath{\sqrt[3]{\,}}}
\newunicodechar{∜}{\ensuremath{\sqrt[4]{\,}}}
\newunicodechar{⃗}{\ensuremath{{}^{\to}}}
\newunicodechar{ⁿ}{\ifmmode^{n}\else\textsuperscript{\ensuremath{n}}\fi}
\newunicodechar{ˣ}{\ifmmode^{x}\else\textsuperscript{\ensuremath{x}}\fi}
\newunicodechar{ᵇ}{\ifmmode^{b}\else\textsuperscript{\ensuremath{b}}\fi}
\newunicodechar{ʳ}{\ifmmode^{r}\else\textsuperscript{\ensuremath{r}}\fi}
\newunicodechar{ᵘ}{\ifmmode^{u}\else\textsuperscript{\ensuremath{u}}\fi}
\newunicodechar{ʸ}{\ifmmode^{y}\else\textsuperscript{\ensuremath{y}}\fi}
\newunicodechar{ᵃ}{\ifmmode^{a}\else\textsuperscript{\ensuremath{a}}\fi}
\newunicodechar{ᵉ}{\ifmmode^{e}\else\textsuperscript{\ensuremath{e}}\fi}
\newunicodechar{ᵏ}{\ifmmode^{k}\else\textsuperscript{\ensuremath{k}}\fi}
\newunicodechar{ᵖ}{\ifmmode^{p}\else\textsuperscript{\ensuremath{p}}\fi}
\newunicodechar{ᴺ}{\ifmmode^{N}\else\textsuperscript{\ensuremath{N}}\fi}
\newunicodechar{ᶜ}{\ifmmode^{c}\else\textsuperscript{\ensuremath{c}}\fi}
\newunicodechar{ʰ}{\ifmmode^{h}\else\textsuperscript{\ensuremath{h}}\fi}
\newunicodechar{ᵠ}{\ifmmode^{q}\else\textsuperscript{\ensuremath{q}}\fi}
\newunicodechar{ₙ}{\ifmmode_{n}\else\textsubscript{\ensuremath{n}}\fi}
\newunicodechar{ₐ}{\ifmmode_{a}\else\textsubscript{\ensuremath{a}}\fi}
\newunicodechar{ᵢ}{\ifmmode_{i}\else\textsubscript{\ensuremath{i}}\fi}
\newunicodechar{ᵣ}{\ifmmode_{r}\else\textsubscript{\ensuremath{r}}\fi}
\newunicodechar{ₖ}{\ifmmode_{k}\else\textsubscript{\ensuremath{k}}\fi}
\newunicodechar{ᵦ}{\ifmmode_{\beta}\else\textsubscript{\ensuremath{\beta}}\fi}
\newunicodechar{ₓ}{\ifmmode_{x}\else\textsubscript{\ensuremath{x}}\fi}
\newfontfamily\popdisplay{ArchivoBlack-Regular.ttf}[Path=../../../fonts/]
\newfontfamily\poplogo{SpaceGrotesk-Bold.ttf}[Path=../../../fonts/]
\newfontfamily\popsemi{PlusJakartaSans-SemiBold.ttf}[Path=../../../fonts/]
\newfontfamily\popxbold{PlusJakartaSans-ExtraBold.ttf}[Path=../../../fonts/]
\newfontfamily\popxboldls{PlusJakartaSans-ExtraBold.ttf}[Path=../../../fonts/,LetterSpace=3.0]

\setlist[enumerate]{leftmargin=1.7em,itemsep=5pt,topsep=4pt}
\setlist[itemize]{leftmargin=1.7em,itemsep=4pt,topsep=4pt,label={\color{poporange}\textbullet}}
\setlength{\parindent}{0pt}
\setlength{\parskip}{5pt}
\renewcommand{\arraystretch}{1.25}

% pgfplots : style Pop par défaut
\pgfplotsset{
  every axis/.append style={axis line style={popink,line width=1pt},tick style={popink},
    grid style={popline,line width=0.5pt},label style={font=\small},tick label style={font=\footnotesize}},
  popcurve/.style={poporange,line width=1.8pt,no markers,smooth},
}
% Même style disponible dans un \draw TikZ simple (hors pgfplots)
\tikzset{popcurve/.style={poporange,line width=1.8pt}}
\tikzset{popgrid/.style={popline,very thin}}
\colorlet{popcurve}{poporange} % le nom du style est parfois utilisé comme couleur

% ============================================================
% Pill / squiggle
% ============================================================
\newcommand{\poppill}[3]{%
  \tikz[baseline=-0.55ex]\node[draw=popink,line width=1.2pt,rounded corners=2.6pt,%
    fill=#2,inner xsep=6.5pt,inner ysep=3pt]{%
    \popxboldls\fontsize{7}{7}\selectfont\color{#3}\MakeUppercase{#1}};%
}
\newcommand{\popsquiggle}[1]{%
  \tikz[baseline=-0.4ex]{
    \draw[#1,line width=2.6pt,line cap=round]
      plot [smooth] coordinates {(0,0.09) (0.34,0.01) (0.68,0.21) (1.02,0.01) (1.36,0.10)};
  }%
}
% Mot-clé surligné (marqueur orange sous le texte)
\newcommand{\cle}[1]{{\popxbold\color{popdark}#1}}

% ============================================================
% En-tête / pied
% ============================================================
\pagestyle{fancy}
\fancyhf{}
\renewcommand{\headrulewidth}{0pt}
\renewcommand{\footrulewidth}{0pt}
\lhead{\raisebox{-0.15ex}{\poplogo\fontsize{13}{13}\selectfont\color{popink}OnBuch\color{poporange}+}}
\rhead{\poppill{\DOCMATIERE\ \textbullet\ \DOCNIVEAU\ \textbullet\ Chapter \DOCLECON}{white}{popink}}
\lfoot{\popxbold\fontsize{7.6}{7.6}\selectfont\color{popmuted}\MakeUppercase{Signed course OnBuch+}\ \textbullet\ {\popsemi\DOCTITRE}}
\rfoot{\tikz[baseline=-0.6ex]\node[rounded corners=8.5pt,fill=popink,minimum height=17pt,minimum width=17pt,inner xsep=5pt,inner sep=0]{\popxbold\fontsize{8}{8}\selectfont\color{popcream}\thepage\,/\,\pageref*{LastPage}};}

% ============================================================
% Titres
% ============================================================
\newcommand{\chaptitre}[2]{%
  \begin{tcolorbox}[enhanced,colback=poporange,colframe=popink,boxrule=2.6pt,arc=11pt,
      drop shadow=popink,left=15pt,right=15pt,top=12pt,bottom=11pt,before skip=3pt,after skip=18pt]
    \poppill{#2}{popink}{popcream}\par\vspace{9pt}
    {\raggedright\hyphenpenalty=10000\exhyphenpenalty=10000\popdisplay\fontsize{22}{24}\selectfont\color{popcream}\MakeUppercase{#1}\par}
  \end{tcolorbox}
}
% Section numérotée : pastille orange avec le numéro + titre Archivo
\newcounter{popsec}
\newcommand{\coursec}[1]{%
  \stepcounter{popsec}\setcounter{popsub}{0}%
  \par\vspace{16pt}\noindent
  \begin{minipage}{\linewidth}
  \tikz[baseline=-0.7ex]\node[draw=popink,line width=1.6pt,fill=poporange,rounded corners=4pt,minimum size=22pt,inner sep=0]{\popdisplay\fontsize{12}{12}\selectfont\color{popcream}\thepopsec};%
  \hspace{8pt}{\popdisplay\fontsize{15}{17}\selectfont\color{popink}\MakeUppercase{#1}}\par
  \vspace{4pt}\noindent{\color{poporange}\rule{36pt}{3.6pt}}
  \end{minipage}\par\nopagebreak\vspace{8pt}
}
\newcounter{popsub}
\newcommand{\courssub}[1]{%
  \stepcounter{popsub}%
  \par\vspace{9pt}\noindent{\popxbold\fontsize{11.5}{13}\selectfont\color{popink}{\color{poporange}\thepopsec.\thepopsub}\hspace{6pt}#1}\par\nopagebreak\vspace{4pt}
}

% ============================================================
% Boîtes pédagogiques — même recette : fond blanc, bordure épaisse colorée,
% ombre dure encre, badge plein. Titre optionnel [#1].
% ============================================================
\tcbset{popbase/.style={breakable,enhanced,colback=white,boxrule=2.3pt,arc=9pt,
  shadow={3pt}{-3pt}{0pt}{popink},left=14pt,right=14pt,top=11pt,bottom=11pt,before skip=11pt,after skip=15pt,
  fontupper=\popsemi\fontsize{10.6}{14.5}\selectfont\color{popink},
  fontlower=\popsemi\fontsize{10.6}{14.5}\selectfont\color{popink}}}
% Coche dessinée (le glyphe ✓ n'existe pas dans les polices du document)
\newcommand{\popcheck}[1]{\tikz[baseline=-0.5ex]\draw[#1,line width=1.6pt,line cap=round,line join=round](-0.13,0.0)--(-0.04,-0.09)--(0.13,0.1);}
\providecommand{\coche}{\popcheck{popgreen}}
\providecommand{\resultat}[1]{\par\smallskip\noindent\fcolorbox{popgreen}{popgreenL}{\parbox{\dimexpr\linewidth-2\fboxsep-2\fboxrule\relax}{\textbf{Result.} #1}}\par\smallskip}
\newcommand{\popboxhead}[3]{\poppill{#1}{#2}{white}\ifx\relax#3\relax\else\hspace{6pt}{\popxbold\fontsize{10.6}{12}\selectfont\color{popink}#3}\fi\par\vspace{7pt}}

\newtcolorbox{aretenir}[1][]{popbase,colframe=popgreen,before upper={\popboxhead{Key points}{popgreen}{#1}}}
\newtcolorbox{exemplebox}[1][]{popbase,colframe=poporange,before upper={\popboxhead{Exemple}{poporange}{#1}}}
\newtcolorbox{definition}[1][]{popbase,colframe=popblue,before upper={\popboxhead{Definition}{popblue}{#1}}}
\newtcolorbox{propriete}[1][]{popbase,colframe=poppurple,before upper={\popboxhead{Property}{poppurple}{#1}}}
\newtcolorbox{methode}[1][]{popbase,colframe=popgold,before upper={\popboxhead{Method}{popgold}{#1}}}
\newtcolorbox{attention}[1][]{popbase,colframe=poppink,before upper={\popboxhead{Warning}{poppink}{#1}}}
\newtcolorbox{experience}[1][]{popbase,colframe=popblue,colback=popblueL,before upper={\popboxhead{Experiment}{popblue}{#1}}}
% « Le savais-tu ? » : carte violette pleine (contexte, culture, Cameroun)
\newtcolorbox{savaistu}[1][]{popbase,colback=poppurple,colframe=popink,
  fontupper=\popsemi\fontsize{10.4}{14.2}\selectfont\color{white},
  before upper={\poppill{Did you know?}{white}{poppurple}\ifx\relax#1\relax\else\hspace{6pt}{\popxbold\color{white}#1}\fi\par\vspace{7pt}}}
% Exercice résolu : énoncé en haut, \tcblower puis la solution sur fond crème
\newcounter{popexo}\newcounter{popexr}
\newtcolorbox{exoresolu}[1][]{popbase,colframe=poporange,colbacklower=popcream,
  segmentation style={popink,line width=1.2pt,dashed},
  before upper={\stepcounter{popexr}\popboxhead{Worked example \thepopexr}{poporange}{#1}},
  before lower={\poppill{Solution}{popink}{popcream}\par\vspace{6pt}}}
% Exercice (non résolu en place) — la correction est dans la partie Corrigés
\newtcolorbox{exercice}[1][]{popbase,colframe=popink,
  before upper={\stepcounter{popexo}\popboxhead{Exercise \thepopexo}{popink}{#1}}}
% Corrigé
\newtcolorbox{corrige}[1][]{popbase,colframe=popgreen,colback=popgreenL,
  before upper={\popboxhead{Solution}{popgreen}{#1}}}
% Objectifs / prérequis / bilan
\newtcolorbox{objectifs}{popbase,colframe=popink,colback=poporangeL,
  before upper={\popboxhead{Chapter objectives}{popink}{}}}
\newtcolorbox{prerequis}{popbase,colframe=popink,colback=popblueL,
  before upper={\popboxhead{Prerequisites}{popblue}{}}}
\newtcolorbox{bilan}{popbase,colframe=popink,colback=white,boxrule=2.8pt,
  before upper={\popboxhead{Fiche bilan}{poporange}{L'essentiel à connaître pour l'examen}}}
% Figure encadrée avec légende
\newtcolorbox{popfigure}[1][]{enhanced,breakable=false,colback=white,colframe=popline,boxrule=1.4pt,arc=8pt,
  left=10pt,right=10pt,top=10pt,bottom=8pt,before skip=10pt,after skip=12pt,halign=center,
  lower separated=false,halign lower=center,
  fontlower=\popsemi\fontsize{9}{11.5}\selectfont\color{popmuted},
  #1}
\newcommand{\legende}[1]{\par\vspace{6pt}{\popsemi\fontsize{9}{11.5}\selectfont\color{popmuted}#1}}

% ============================================================
% Signature OnBuch+
% ============================================================
\newcommand{\poplogoinline}[1]{{\poplogo\fontsize{#1}{#1}\selectfont\color{popink}OnBuch\color{poporange}+}}
\newcommand{\signatureonbuch}{%
  \begin{tcolorbox}[enhanced,colback=popink,colframe=popink,boxrule=2.2pt,arc=10pt,
      drop shadow=poporange,left=14pt,right=14pt,top=10pt,bottom=10pt,before skip=0pt,after skip=0pt]
    \tikz[baseline=-0.7ex]\node[circle,fill=popgreen,draw=popcream,line width=1.3pt,minimum size=22pt,inner sep=0]{\popcheck{white}};%
    \hspace{9pt}\begin{minipage}[c]{0.62\linewidth}
      {\popxbold\fontsize{12}{13}\selectfont\color{popcream}Signed\ \textbullet\ {\poplogo OnBuch\color{poporange}+}}\par\vspace{2pt}
      {\raggedright\popsemi\fontsize{8}{10.5}\selectfont\color{popline}\MakeUppercase{OnBuch+ teaching team}\\\MakeUppercase{Follows the official MINESEC syllabus}\par}
    \end{minipage}\hfill
    {\poplogo\fontsize{22}{22}\selectfont\color{popcream}OnBuch\color{poporange}+}
  \end{tcolorbox}%
}

% ============================================================
% Page de garde
% #1 titre · #2 matière · #3 niveau · #4 module · #5 n° leçon · #6 durée ·
% #7 description / objectifs (texte ou itemize)
% ============================================================
\newcommand{\pagegarde}[7]{%
  \thispagestyle{empty}
  \newgeometry{a4paper,margin=1.7cm,top=1.6cm,bottom=1.4cm}%
  \begin{tikzpicture}[remember picture,overlay]
    \fill[poporange!18] ([xshift=-3.2cm,yshift=-3.6cm]current page.north east) circle (4.6cm);
    \fill[poppurple!14] ([xshift=2.4cm,yshift=7.6cm]current page.south west) circle (3.8cm);
  \end{tikzpicture}%
  {\poplogo\fontsize{30}{30}\selectfont\color{popink}OnBuch\color{poporange}+}\hfill
  \raisebox{0.6ex}{\poppill{Signed course \textbullet\ Premium}{popink}{popcream}}\par
  \vspace{1pt}\popsquiggle{poppurple}\par
  \vspace{8pt}
  \begin{tcolorbox}[enhanced,colback=poporange,colframe=popink,boxrule=2.8pt,arc=13pt,
      drop shadow=popink,left=18pt,right=18pt,top=12pt,bottom=13pt,after skip=10pt]
    \poppill{#2 \textbullet\ #3}{popink}{popcream}\hspace{5pt}\poppill{#4}{white}{popink}\par\vspace{8pt}
    {\popdisplay\fontsize{14}{14}\selectfont\color{popink}CHAPTER #5}\par\vspace{4pt}
    {\raggedright\hyphenpenalty=10000\exhyphenpenalty=10000\popdisplay\fontsize{25}{27}\selectfont\color{popcream}\MakeUppercase{#1}\par}
    \vspace{8pt}
    {\popxbold\fontsize{9.5}{9.5}\selectfont\color{popink}\MakeUppercase{Suggested time: #6}}
  \end{tcolorbox}
  \begin{tcolorbox}[enhanced,colback=white,colframe=popink,boxrule=2.2pt,arc=10pt,
      drop shadow=popink,left=16pt,right=16pt,top=10pt,bottom=10pt,after skip=8pt]
    \poppill{In this chapter}{poppurple}{white}\par\vspace{5pt}
    \popsemi\fontsize{9.3}{12.6}\selectfont\color{popink}#7
  \end{tcolorbox}
  \noindent\begin{minipage}[t]{0.487\linewidth}
    \begin{tcolorbox}[enhanced,colback=popgreen,colframe=popink,boxrule=2.2pt,arc=10pt,
        drop shadow=popink,left=11pt,right=11pt,top=10pt,bottom=10pt,height=3.1cm,valign=center]
      \tcbox[colback=white,colframe=popink,boxrule=1.3pt,arc=5pt,boxsep=0pt,nobeforeafter,
        left=3pt,right=3pt,top=3pt,bottom=3pt,valign=center]{\qrcode[height=1.9cm]{https://play.google.com/store/apps/details?id=com.luvvix.onbuch&hl=en}}%
      \hspace{8pt}\begin{minipage}[c]{0.5\linewidth}
        {\popdisplay\fontsize{11}{12.5}\selectfont\color{white}\MakeUppercase{Download OnBuch}}\par\vspace{4pt}
        {\raggedright\popsemi\fontsize{8.4}{11}\selectfont\color{white}Free \textbullet\ Google Play\\AI tutor, past papers, results\par}
      \end{minipage}
    \end{tcolorbox}
  \end{minipage}\hfill
  \begin{minipage}[t]{0.487\linewidth}
    \begin{tcolorbox}[enhanced,colback=poppurple,colframe=popink,boxrule=2.2pt,arc=10pt,
        drop shadow=popink,left=11pt,right=11pt,top=10pt,bottom=10pt,height=3.1cm,valign=center]
      \tikz[baseline=-0.7ex]\node[circle,fill=white,draw=popink,line width=1.3pt,minimum size=1.6cm,inner sep=0]{\popdisplay\fontsize{24}{24}\selectfont\color{poppurple}+};%
      \hspace{8pt}\begin{minipage}[c]{0.6\linewidth}
        {\popdisplay\fontsize{11}{12.5}\selectfont\color{white}\MakeUppercase{Go OnBuch+}}\par\vspace{4pt}
        {\raggedright\popsemi\fontsize{8.4}{11}\selectfont\color{white}All signed courses, unlimited AI tutor, PDF sheets — OnBuch+ tab in the app\par}
      \end{minipage}
    \end{tcolorbox}
  \end{minipage}
  \vspace{8pt}
  \popcontenu
  \vfill
  \signatureonbuch
  \restoregeometry
  \newpage
}

% Rangée « Ce cours contient » (page de garde)
\newcommand{\poptile}[3]{% #1 couleur #2 gros texte #3 légende
  \begin{tcolorbox}[enhanced,colback=#1,colframe=popink,boxrule=2pt,arc=9pt,shadow={2.5pt}{-2.5pt}{0pt}{popink},
      left=6pt,right=6pt,top=9pt,bottom=9pt,halign=center,valign=center,height=1.75cm,width=\linewidth,nobeforeafter]
    {\popdisplay\fontsize{12.5}{13}\selectfont\color{white}\MakeUppercase{#2}}\par\vspace{3pt}
    {\centering\popsemi\fontsize{7.8}{9.5}\selectfont\color{white}#3\par}
  \end{tcolorbox}}
\newcommand{\popcontenu}{%
  \par\noindent{\popxboldls\fontsize{7.5}{7.5}\selectfont\color{popmuted}THIS SIGNED COURSE CONTAINS}\par\vspace{7pt}
  \noindent\begin{minipage}[t]{0.235\linewidth}\poptile{popblue}{Course}{complete, illustrated, step by step}\end{minipage}\hfill
  \begin{minipage}[t]{0.235\linewidth}\poptile{poporange}{Methods}{and worked examples}\end{minipage}\hfill
  \begin{minipage}[t]{0.235\linewidth}\poptile{poppink}{Exercises}{exam style + solutions}\end{minipage}\hfill
  \begin{minipage}[t]{0.235\linewidth}\poptile{popgreen}{Summary sheet}{the essentials to remember}\end{minipage}\par}

% Sommaire
\newtcolorbox{sommaire}{enhanced,colback=white,colframe=popink,boxrule=2.2pt,arc=10pt,
  drop shadow=popink,left=18pt,right=18pt,top=13pt,bottom=13pt,before skip=6pt,after skip=20pt,
  before upper={\poppill{Contents}{poppurple}{white}\par\vspace{9pt}\popsemi\fontsize{10.6}{14.5}\selectfont\color{popink}}}

% Signature de fin de cours — poussée tout en bas de la dernière page
\newcommand{\findecours}{%
  \par\vfill\nopagebreak
  \begin{center}\popsquiggle{poporange}\end{center}\vspace{4pt}
  \signatureonbuch
}
\AtBeginDocument{\def\llbracket{\mathopen{[\mkern-3.5mu[}}\def\rrbracket{\mathclose{]\mkern-3.5mu]}}} % intervalles d'entiers (glyphe ⟦ absent de la police)
% Glyphes absents de Fira Math : redéfinis à partir de glyphes disponibles
\AtBeginDocument{%
  \def\setminus{\mathbin{\backslash}}\let\smallsetminus\setminus
  \def\vdots{\mathord{\vbox{\baselineskip3.2pt\lineskiplimit0pt\kern4pt\hbox{.}\hbox{.}\hbox{.}}}}%
  \def\ddots{\mathinner{\mkern1mu\raise6.5pt\hbox{.}\mkern2mu\raise3.6pt\hbox{.}\mkern2mu\raise0.7pt\hbox{.}\mkern1mu}}%
  \def\triangle{\mathord{\scalebox{0.9}{$\symup{\Delta}$}}}%
  \def\nabla{\mathord{\raisebox{\depth}{\scalebox{1}[-1]{$\symup{\Delta}$}}}}%
  \def\checkmark{\popcheck{popdark}}%
  \def\star{\mathbin{\ast}}\def\bigstar{\mathord{\ast}}%
  \def\diamond{\mathbin{\scalebox{0.6}{$\Diamond$}}}%
  \def\bigcup{\mathop{\mathchoice{\raisebox{-0.3em}{\scalebox{1.7}{$\cup$}}}{\scalebox{1.25}{$\cup$}}{\cup}{\cup}}\displaylimits}%
  \def\bigcap{\mathop{\mathchoice{\raisebox{-0.3em}{\scalebox{1.7}{$\cap$}}}{\scalebox{1.25}{$\cap$}}{\cap}{\cap}}\displaylimits}%
  \def\lesseqqgtr{\mathrel{\lessgtr}}\def\bigoplus{\mathop{\oplus}\displaylimits}%
}
\providecommand{\ssi}{if and only if} % abréviation fréquente
\AtBeginDocument{\providecommand{\wideparen}{\overparen}} % arc de cercle

% Fonctions usuelles que les modèles utilisent sans que amsmath les fournisse
\providecommand{\arsinh}{\operatorname{arsinh}}\providecommand{\arcosh}{\operatorname{arcosh}}
\providecommand{\artanh}{\operatorname{artanh}}\providecommand{\arsech}{\operatorname{arsech}}
\providecommand{\arcsch}{\operatorname{arcsch}}\providecommand{\arcoth}{\operatorname{arcoth}}
\providecommand{\arcsinh}{\operatorname{arsinh}}\providecommand{\arccosh}{\operatorname{arcosh}}
\providecommand{\arctanh}{\operatorname{artanh}}\providecommand{\sech}{\operatorname{sech}}
\providecommand{\csch}{\operatorname{csch}}\providecommand{\arcsec}{\operatorname{arcsec}}
\providecommand{\arccsc}{\operatorname{arccsc}}\providecommand{\arccot}{\operatorname{arccot}}
\providecommand{\sgn}{\operatorname{sgn}}\providecommand{\lcm}{\operatorname{lcm}}
\providecommand{\Var}{\operatorname{Var}}\providecommand{\Cov}{\operatorname{Cov}}

%% FIGFIX-BEGIN v3 — correctifs globaux des figures (docs/onbuch-generation/tools/figfix)
\usepackage{adjustbox}\usepackage{etoolbox}
% toute figure TikZ/pgfplots plus large que la ligne est réduite à la largeur disponible
\BeforeBeginEnvironment{tikzpicture}{\begin{adjustbox}{max width=\linewidth}}
\AfterEndEnvironment{tikzpicture}{\end{adjustbox}}
% légendes pgfplots lisibles par-dessus les courbes
\pgfplotsset{every axis/.append style={legend style={fill=white,fill opacity=0.92,text opacity=1,draw=black!15,inner sep=3pt}}}
% étiquettes d'arêtes (syntaxe edge["..."]) avec fond blanc
\tikzset{every edge quotes/.append style={fill=white,inner sep=1.2pt}}
% bulles de cartes mentales : texte à la ligne dans la bulle, bulle un peu plus grande
\tikzset{every concept/.append style={text width=2.0cm,align=center,minimum size=2.3cm,inner sep=2pt}}
%% FIGFIX-END
\begin{document}

\pagegarde{Simple and Damped Harmonic motion}{Further Mathematics}{Upper Sixth}{Upper Sixth — Further mathematics}{9}{3 periods}{This chapter studies simple harmonic motion (SHM) and damped harmonic motion through their defining differential equations, their solutions, and their energy properties. It closes with modelling real oscillating systems — springs, pendulums, vehicle suspensions — in the spirit of GCE Advanced Level problems.
\vspace{4pt}
\begin{multicols}{2}\small
\begin{itemize}
\item By the end of this chapter I can recognise a physical situation that leads to simple harmonic motion and derive the equation ẍ = -ω²x from Newton's second law.
\item By the end of this chapter I can solve the SHM equation and write the solution in the forms x = a sin(ωt + ε) and x = a cos(ωt + φ), choosing the form adapted to the initial conditions.
\item By the end of this chapter I can determine the amplitude, period, frequency and phase of an SHM and compute the velocity and acceleration at any point of the motion.
\item By the end of this chapter I can use the relation v² = ω²(a² - x²) to link speed and position, and analyse the energy of an oscillator.
\item By the end of this chapter I can set up the damped harmonic motion equation mẍ = -kx - λẋ and identify the damping coefficient.
\item By the end of this chapter I can solve the damped equation using the auxiliary (characteristic) equation and distinguish light, critical and heavy damping.
\end{itemize}\end{multicols}}

\begin{sommaire}
\begin{multicols}{2}
\begin{enumerate}
\item Simple Harmonic Motion: Definition and Equation
\item Further Results on SHM: Energy, Pendulums and Composite Problems
\item Damped Harmonic Motion
\item Applications and GCE Problem-Solving
\item Integration activity
\item Methods and worked examples
\item Exercises
\item Solutions to the exercises
\item Summary sheet
\end{enumerate}
\end{multicols}
\end{sommaire}

\begin{objectifs}
\begin{itemize}
\item By the end of this chapter I can recognise a physical situation that leads to simple harmonic motion and derive the equation ẍ = -ω²x from Newton's second law.
\item By the end of this chapter I can solve the SHM equation and write the solution in the forms x = a sin(ωt + ε) and x = a cos(ωt + φ), choosing the form adapted to the initial conditions.
\item By the end of this chapter I can determine the amplitude, period, frequency and phase of an SHM and compute the velocity and acceleration at any point of the motion.
\item By the end of this chapter I can use the relation v² = ω²(a² - x²) to link speed and position, and analyse the energy of an oscillator.
\item By the end of this chapter I can set up the damped harmonic motion equation mẍ = -kx - λẋ and identify the damping coefficient.
\item By the end of this chapter I can solve the damped equation using the auxiliary (characteristic) equation and distinguish light, critical and heavy damping.
\item By the end of this chapter I can sketch and interpret the displacement-time graphs of under-damped, critically damped and over-damped motions.
\item By the end of this chapter I can model concrete systems (mass-spring, simple pendulum for small angles, vehicle suspension) and compute their periods or damping behaviour.
\item By the end of this chapter I can solve full GCE-style problems combining derivation of the equation of motion, solution, and interpretation of results.
\end{itemize}
\end{objectifs}

\begin{prerequis}
\begin{itemize}
\item Second-order linear differential equations with constant coefficients and the auxiliary equation (Lower Sixth / Upper Sixth pure mathematics)
\item Newton's second law of motion and resolution of forces (Lower Sixth mechanics)
\item Hooke's law for springs and elastic strings: T = λx/l (Lower Sixth mechanics)
\item Trigonometric functions, their derivatives, and the R sin(ωt + ε) compound-angle form
\item Kinematics: velocity and acceleration as first and second derivatives of displacement
\end{itemize}
\end{prerequis}

\newpage

\coursec{Simple Harmonic Motion: Definition and Equation}

Every evening in Bamenda, a trader hangs her spring scale and watches the pointer bounce up and down before settling on the weight of a bag of beans. On the road to Douala, passengers feel their bus bounce after each pothole, yet the oscillations die out quickly thanks to the shock absorbers. Why does the pointer oscillate, why does the motion gradually fade, and how can we predict the period of these vibrations? This chapter develops the mathematics — differential equations of the form $\ddot{x} = -\omega^2 x$ and $\ddot{x} + 2k\dot{x} + \omega^2 x = 0$ — that answer these questions and appear every year in the GCE Further Mathematics paper.

\courssub{What is Simple Harmonic Motion?}

\begin{definition}[Simple Harmonic Motion (SHM)]
A particle $P$ describes \textbf{simple harmonic motion} about a fixed centre $O$ if its acceleration is always directed towards $O$ and is proportional to its displacement from $O$. If $x$ is the displacement from $O$ at time $t$, then
\[
\ddot{x} = -\omega^2 x
\]
where $\omega > 0$ is a constant called the \cle{angular frequency}.
\end{definition}

The negative sign is crucial: it tells us that the acceleration is \emph{restoring} — it always points back towards the centre $O$. The constant $\omega^2$ is the constant of proportionality between the magnitude of the acceleration and the displacement.

\courssub{Derivation from a Horizontal Mass–Spring System}

Consider a particle of mass $m$ attached to one end of a light elastic spring of natural length $l$ and modulus of elasticity $\lambda$. The other end of the spring is fixed at a point $A$ on a smooth horizontal table. Let $O$ be the point where the spring is unstretched (natural length). At time $t$, let the displacement of $P$ from $O$ be $x$ (positive in the direction $A$ to $P$). The extension of the spring is then $x$.

\begin{popfigure}
\begin{tikzpicture}[scale=1.2]
% Wall
\draw[fill=popdark!20] (-0.5,0.5) rectangle (0,-0.5);
\draw[thick] (-0.5,0.5) -- (0,0.5) -- (0,-0.5) -- (-0.5,-0.5);
% Spring (coil)
\draw[popblue, thick] (0,0) -- (0.3,0);
\foreach \i in {1,...,8} {
    \draw[popblue, thick] (0.3+0.3*\i,0) -- (0.3+0.3*\i+0.15,0.2) -- (0.3+0.3*\i+0.3,0) -- (0.3+0.3*\i+0.45,-0.2) -- (0.3+0.3*\i+0.6,0);
}
\draw[popblue, thick] (0.3+0.3*8+0.6,0) -- (3.5,0);
% Mass
\draw[fill=poporange!30, thick] (3.5,-0.4) rectangle (4.5,0.4);
\node at (4,0) {$m$};
% Natural length marker
\draw[<->, popgreen] (0, -0.7) -- (3.5, -0.7);
\node[popgreen] at (1.75, -1.0) {Natural length $l$};
% Displacement x
\draw[<->, poppurple] (3.5, 0.7) -- (4.5, 0.7);
\node[poppurple] at (4, 1.0) {$x$};
% Tension force
\draw[->, thick, popred] (4.5,0) -- (5.2,0);
\node[popred] at (4.85, 0.3) {$T = \frac{\lambda x}{l}$};
% Acceleration vector
\draw[->, thick, poppink] (4.5,-0.4) -- (5.2,-0.4);
\node[poppink] at (4.85, -0.7) {$m\ddot{x}$};
% Axes
\draw[->] (-0.5,0) -- (5.5,0) node[right] {$x$};
\end{tikzpicture}
\legende{Mass–spring system on a smooth horizontal table. The tension $T$ acts as the restoring force.}
\end{popfigure}

By Hooke's law, the tension in the spring is $T = \dfrac{\lambda x}{l}$. This tension acts towards $A$ (i.e. in the negative $x$ direction). Applying Newton's second law in the positive $x$ direction:
\[
-m\ddot{x} = T = \frac{\lambda x}{l} \quad \Longrightarrow \quad \ddot{x} = -\frac{\lambda}{ml} x.
\]
This is exactly the SHM equation $\ddot{x} = -\omega^2 x$ with
\[
\omega^2 = \frac{\lambda}{ml}.
\]

\begin{savaistu}[Modulus of Elasticity]
The modulus $\lambda$ (measured in newtons) is the force required to double the length of the spring. For a spring of natural length $l$, the tension for an extension $x$ is $T = \frac{\lambda x}{l}$. This is Hooke's law in the form used in the GCE syllabus. Do not confuse $\lambda$ with the spring constant $k = \lambda/l$ (stiffness, in N/m).
\end{savaistu}

\courssub{Equilibrium Position and Change of Variable}

In many problems the spring is not at its natural length when the system is at rest. For example, a vertical spring stretches under the weight of the mass. The \cle{equilibrium position} $E$ is where the net force is zero. If we measure displacement $X$ from $E$ (i.e. $X = x - e$ where $e$ is the equilibrium extension), then the equation of motion becomes $\ddot{X} = -\omega^2 X$, proving SHM about $E$.

\begin{methode}[Proving a Motion is SHM]
\begin{enumerate}
    \item Identify the forces acting on the particle.
    \item Find the equilibrium position $E$ (where resultant force $= 0$).
    \item Let $X$ be the displacement from $E$ (choose a positive direction).
    \item Apply $F = ma$ in the positive direction, expressing all forces in terms of $X$.
    \item Show that the equation reduces to $\ddot{X} = -\omega^2 X$ and read off $\omega$.
\end{enumerate}
\end{methode}

\begin{attention}[Amplitude vs. Initial Displacement]
The \cle{amplitude} $a$ is the maximum displacement from the \textbf{centre of oscillation} (equilibrium position), not from the spring's natural length. Always locate the equilibrium position first! A common mistake is to take the initial extension from the natural length as the amplitude.
\end{attention}

\courssub{General Solution and Initial Conditions}

The differential equation $\ddot{x} = -\omega^2 x$ has general solution
\[
x = a \sin(\omega t + \varepsilon)
\]
where $a \ge 0$ is the \cle{amplitude} and $\varepsilon$ is the \cle{phase constant} (or initial phase). Equivalently,
\[
x = A \cos \omega t + B \sin \omega t \quad \text{with} \quad a = \sqrt{A^2+B^2},\; \tan\varepsilon = \frac{A}{B}.
\]

The constants are determined by the initial conditions $x(0)$ and $\dot{x}(0)$:
\[
x(0) = a\sin\varepsilon,\qquad \dot{x}(0) = a\omega\cos\varepsilon.
\]

\begin{aretenir}[GCE Shortcut]
\begin{itemize}
    \item If the particle is \textbf{released from rest} at displacement $x = a$ (maximum displacement), then $\varepsilon = \frac{\pi}{2}$ and $x = a \cos \omega t$.
    \item If the particle is \textbf{projected from the centre} $x=0$ with speed $v_0$, then $\varepsilon = 0$ and $x = \frac{v_0}{\omega}\sin\omega t$.
\end{itemize}
This avoids solving for $\varepsilon$ explicitly.
\end{aretenir}

\courssub{Key Vocabulary and Quantities}

\begin{propriete}[Parameters of SHM]
For a particle executing SHM $x = a\sin(\omega t + \varepsilon)$:
\begin{itemize}
    \item \textbf{Amplitude} $a$: maximum displacement from the centre.
    \item \textbf{Angular frequency} $\omega$: $\omega = \sqrt{\frac{\lambda}{ml}}$ (rad/s).
    \item \textbf{Period} $T = \dfrac{2\pi}{\omega}$: time for one complete oscillation.
    \item \textbf{Frequency} $f = \dfrac{1}{T} = \dfrac{\omega}{2\pi}$: number of oscillations per second (Hz).
    \item \textbf{Phase} at time $t$: $\phi = \omega t + \varepsilon$.
\end{itemize}
\end{propriete}

\begin{popfigure}
\begin{tikzpicture}
\begin{axis}[
    width=0.9\linewidth,
    height=6cm,
    axis lines=middle,
    xlabel=$t$, ylabel=$x$,
    xmin=0, xmax=6.5,
    ymin=-1.5, ymax=1.5,
    xtick={0,1.57,3.14,4.71,6.28},
    xticklabels={$0$, $\frac{\pi}{\omega}$, $\frac{2\pi}{\omega}$, $\frac{3\pi}{\omega}$, $\frac{4\pi}{\omega}$},
    ytick={-1,0,1},
    yticklabels={$-a$, $0$, $a$},
    domain=0:6.28,
    samples=200,
    tick label style={font=\small},
    label style={font=\small},
    clip=false
]
\addplot[popblue, thick] {sin(deg(x))};
% Amplitude annotation
\draw[poporange, <->] (axis cs:0,0) -- (axis cs:0,1) node[midway, right] {$a$};
% Period annotation
\draw[popgreen, <->] (axis cs:0,-1.3) -- (axis cs:6.28,-1.3) node[midway, below] {$T = \frac{2\pi}{\omega}$};
% Phase shift example (dashed)
\addplot[poppurple, dashed, thick] {sin(deg(x+1))};
\node[poppurple, anchor=west] at (axis cs:6.5,0.8) {$x = a\sin(\omega t+\varepsilon)$};
\node[popblue, anchor=west] at (axis cs:6.5,0.5) {$x = a\sin\omega t$};
\end{axis}
\end{tikzpicture}
\legende{Graph of $x = a\sin(\omega t + \varepsilon)$ showing amplitude $a$ and period $T$. The dashed curve shows a phase shift $\varepsilon$.}
\end{popfigure}

\courssub{Velocity and Acceleration in SHM}

Differentiating $x = a\sin(\omega t + \varepsilon)$:
\[
\dot{x} = a\omega \cos(\omega t + \varepsilon), \qquad \ddot{x} = -a\omega^2 \sin(\omega t + \varepsilon) = -\omega^2 x.
\]

\begin{propriete}[Maximum Values]
\begin{itemize}
    \item Maximum speed: $v_{\max} = a\omega$ (occurs at the centre $x=0$).
    \item Maximum acceleration: $|\ddot{x}|_{\max} = a\omega^2$ (occurs at the extremes $x = \pm a$).
\end{itemize}
\end{propriete}

\courssub{The Speed–Displacement Relation}

Eliminating $t$ between $x = a\sin(\omega t + \varepsilon)$ and $\dot{x} = a\omega\cos(\omega t + \varepsilon)$ using $\sin^2\theta + \cos^2\theta = 1$ gives the fundamental relation:
\[
v^2 = \omega^2(a^2 - x^2) \quad \text{where } v = \dot{x}.
\]

This tells us that the speed depends only on the displacement, not on the direction of motion.

\begin{popfigure}
\begin{tikzpicture}
\begin{axis}[
    width=0.9\linewidth,
    height=6cm,
    axis lines=middle,
    xlabel=$x$, ylabel=$v$,
    xmin=-1.2, xmax=1.2,
    ymin=-1.2, ymax=1.2,
    xtick={-1,0,1},
    xticklabels={$-a$, $0$, $a$},
    ytick={-1,0,1},
    yticklabels={$-a\omega$, $0$, $a\omega$},
    domain=-1:1,
    samples=200,
    tick label style={font=\small},
    label style={font=\small},
    axis equal image,
    clip=false
]
% Upper half
\addplot[popblue, thick] {sqrt(max(0,1-x^2))};
% Lower half
\addplot[popblue, thick] {-sqrt(max(0,1-x^2))};
% Key points
\addplot[only marks, mark=*, mark size=3, popred] coordinates {(0,1) (0,-1) (1,0) (-1,0)};
\node[popred, anchor=south west] at (axis cs:0,1) {$(0, a\omega)$};
\node[popred, anchor=north west] at (axis cs:0,-1) {$(0, -a\omega)$};
\node[popred, anchor=north] at (axis cs:1,0) {$(a, 0)$};
\node[popred, anchor=north] at (axis cs:-1,0) {$(-a, 0)$};
% Direction arrows
\draw[->, thick, poporange] (axis cs:-0.5,0.866) -- (axis cs:-0.4,0.916);
\draw[->, thick, poporange] (axis cs:0.5,-0.866) -- (axis cs:0.4,-0.916);
\node[poporange] at (axis cs:-0.8,0.5) {Clockwise};
\end{axis}
\end{tikzpicture}
\legende{The ellipse $v^2 = \omega^2(a^2 - x^2)$ in the phase plane. The motion proceeds clockwise.}
\end{popfigure}

\courssub{Vertical Oscillations of a Spring}

A light spring of natural length $l$ and modulus $\lambda$ hangs vertically from a fixed point. A mass $m$ is attached to the lower end. At equilibrium, the extension is $e$ where $mg = \dfrac{\lambda e}{l}$. Let $x$ be the displacement \textbf{downwards} from the equilibrium position. The total extension is $e+x$.

\begin{popfigure}
\begin{tikzpicture}[scale=1.2]
% Ceiling
\draw[fill=popdark!20] (-0.5,2.5) rectangle (2,2);
\draw[thick] (-0.5,2.5) -- (2,2.5) -- (2,2) -- (-0.5,2);
% Spring at equilibrium
\draw[popblue, thick] (1,2) -- (1,1.2);
\foreach \i in {1,...,6} {
    \draw[popblue, thick] (1,2-0.15*\i) -- (1.15,2-0.15*\i-0.075) -- (1,2-0.15*\i-0.15) -- (0.85,2-0.15*\i-0.225) -- (1,2-0.15*\i-0.3);
}
\draw[popblue, thick] (1,2-0.15*6) -- (1,0.5);
% Mass at equilibrium
\draw[fill=poporange!30, thick] (0.5,0.5) rectangle (1.5,0);
\node at (1,0.25) {$m$};
% Equilibrium label
\draw[<->, popgreen] (1.7,2) -- (1.7,0.5);
\node[popgreen, right] at (1.7,1.25) {$e$};
% Displaced position (dashed)
\draw[dashed, thick] (1,0.5) -- (1,-0.3);
\foreach \i in {1,...,3} {
    \draw[dashed, thick] (1,0.5-0.15*\i) -- (1.15,0.5-0.15*\i-0.075) -- (1,0.5-0.15*\i-0.15) -- (0.85,0.5-0.15*\i-0.225) -- (1,0.5-0.15*\i-0.3);
}
\draw[dashed, thick] (1,0.5-0.15*3) -- (1,-0.3);
\draw[dashed, fill=poporange!30, thick] (0.5,-0.3) rectangle (1.5,-0.8);
\node at (1,-0.55) {$m$};
% x arrow
\draw[<->, poppurple] (1.7,0.5) -- (1.7,-0.3);
\node[poppurple, right] at (1.7,0.1) {$x$};
% Forces at displaced position
\draw[->, thick, popred] (1,-0.8) -- (1,-1.2) node[below] {$mg$};
\draw[->, thick, popblue] (1,-0.3) -- (1,0.1) node[above] {$T = \frac{\lambda(e+x)}{l}$};
\end{tikzpicture}
\legende{Vertical mass–spring system. The equilibrium extension is $e$; $x$ is measured downwards from equilibrium.}
\end{popfigure}

Applying Newton's second law downwards:
\[
m\ddot{x} = mg - T = mg - \frac{\lambda(e+x)}{l}.
\]
But $mg = \dfrac{\lambda e}{l}$, so
\[
m\ddot{x} = -\frac{\lambda x}{l} \quad \Longrightarrow \quad \ddot{x} = -\frac{\lambda}{ml} x.
\]
This is SHM about the equilibrium position with the \textbf{same} angular frequency $\omega^2 = \dfrac{\lambda}{ml}$ as the horizontal case. The period is independent of $g$!

\begin{aretenir}[Key Results for Vertical Spring]
\begin{itemize}
    \item Equilibrium extension: $e = \dfrac{mgl}{\lambda}$.
    \item Equation of motion about equilibrium: $\ddot{x} = -\dfrac{\lambda}{ml}x$.
    \item Angular frequency: $\omega = \sqrt{\dfrac{\lambda}{ml}}$.
    \item Period: $T = 2\pi\sqrt{\dfrac{ml}{\lambda}}$.
\end{itemize}
\end{aretenir}

\courssub{Worked Examples}

\begin{exoresolu}[Horizontal Spring on a Smooth Table]
A particle $P$ of mass $0.5\,\mathrm{kg}$ is attached to one end of a light elastic spring of natural length $1.2\,\mathrm{m}$ and modulus $20\,\mathrm{N}$. The other end is fixed at $A$ on a smooth horizontal table. $P$ is held at rest with $AP = 1.5\,\mathrm{m}$ and then released.
\begin{enumerate}
    \item Show that $P$ moves with SHM and find the period.
    \item Find the maximum speed of $P$.
    \item Find the speed of $P$ when $AP = 1.3\,\mathrm{m}$.
\end{enumerate}

\tcblower
\textbf{Solution.}
\begin{enumerate}
    \item The equilibrium position is at natural length $l = 1.2\,\mathrm{m}$ (no other horizontal forces). Let $x$ be the displacement from equilibrium (positive towards $A$). Extension $= x$. Tension $T = \dfrac{\lambda x}{l} = \dfrac{20x}{1.2} = \dfrac{50}{3}x\,\mathrm{N}$.
    Equation of motion: $-T = m\ddot{x} \Rightarrow -\frac{50}{3}x = 0.5\ddot{x} \Rightarrow \ddot{x} = -\frac{100}{3}x$.
    This is SHM with $\omega^2 = \frac{100}{3}$, so $\omega = \frac{10}{\sqrt{3}}\,\mathrm{rad/s}$.
    Period $T = \frac{2\pi}{\omega} = 2\pi\frac{\sqrt{3}}{10} = \frac{\pi\sqrt{3}}{5} \approx 1.088\,\mathrm{s}$.
    \item Amplitude $a$ = initial displacement from equilibrium $= 1.5 - 1.2 = 0.3\,\mathrm{m}$.
    Maximum speed $v_{\max} = a\omega = 0.3 \times \frac{10}{\sqrt{3}} = \frac{3}{\sqrt{3}} = \sqrt{3} \approx 1.732\,\mathrm{m/s}$.
    \item When $AP = 1.3\,\mathrm{m}$, displacement from equilibrium $x = 1.3 - 1.2 = 0.1\,\mathrm{m}$.
    Using $v^2 = \omega^2(a^2 - x^2) = \frac{100}{3}(0.3^2 - 0.1^2) = \frac{100}{3}(0.09 - 0.01) = \frac{100}{3} \times 0.08 = \frac{8}{3}$.
    Speed $v = \sqrt{\frac{8}{3}} = \frac{2\sqrt{2}}{\sqrt{3}} \approx 1.633\,\mathrm{m/s}$.
\end{enumerate}
\end{exoresolu}

\begin{exoresolu}[Vertical Spring — Finding Amplitude and Period]
A light elastic string of natural length $0.8\,\mathrm{m}$ and modulus $12\,\mathrm{N}$ has one end fixed at $A$. A particle of mass $0.4\,\mathrm{kg}$ is attached to the other end and hangs in equilibrium. The particle is pulled down a further $0.2\,\mathrm{m}$ and released from rest.
\begin{enumerate}
    \item Find the equilibrium extension.
    \item Show that the subsequent motion is SHM and find its period.
    \item Find the maximum speed of the particle.
\end{enumerate}

\tcblower
\textbf{Solution.}
\begin{enumerate}
    \item At equilibrium, $mg = \frac{\lambda e}{l} \Rightarrow 0.4 \times 9.8 = \frac{12e}{0.8} \Rightarrow 3.92 = 15e \Rightarrow e = \frac{3.92}{15} = 0.2613\,\mathrm{m}$ (approx).
    \item Let $x$ be displacement downwards from equilibrium. Total extension $= e+x$. Tension $T = \frac{\lambda(e+x)}{l}$.
    Equation downwards: $m\ddot{x} = mg - T = mg - \frac{\lambda(e+x)}{l} = mg - \frac{\lambda e}{l} - \frac{\lambda x}{l} = -\frac{\lambda x}{l}$.
    So $\ddot{x} = -\frac{\lambda}{ml}x = -\frac{12}{0.4 \times 0.8}x = -37.5x$.
    SHM with $\omega^2 = 37.5$, $\omega = \sqrt{37.5} = \frac{5\sqrt{6}}{2} \approx 6.124\,\mathrm{rad/s}$.
    Period $T = \frac{2\pi}{\omega} = \frac{4\pi}{5\sqrt{6}} \approx 1.026\,\mathrm{s}$.
    \item Amplitude $a = 0.2\,\mathrm{m}$ (the additional displacement from equilibrium).
    Maximum speed $v_{\max} = a\omega = 0.2 \times \sqrt{37.5} = \frac{\sqrt{37.5}}{5} \approx 1.225\,\mathrm{m/s}$.
\end{enumerate}
\end{exoresolu}

\courssub{Exercises}

\begin{exercice}[Exercise 1]
A particle of mass $2\,\mathrm{kg}$ is attached to one end of a light elastic spring of natural length $0.5\,\mathrm{m}$ and modulus $50\,\mathrm{N}$. The other end is fixed on a smooth horizontal table. The particle is pulled until the spring length is $0.8\,\mathrm{m}$ and released from rest.
\begin{enumerate}
    \item Prove that the motion is SHM and find $\omega$.
    \item Find the period and amplitude.
    \item Calculate the maximum acceleration.
    \item Find the speed when the spring length is $0.6\,\mathrm{m}$.
\end{enumerate}
\end{exercice}

\begin{exercice}[Exercise 2]
A light elastic string of natural length $1\,\mathrm{m}$ and modulus $20\,\mathrm{N}$ hangs vertically from a fixed point. A mass of $0.5\,\mathrm{kg}$ is attached to the lower end. The mass is pulled down $0.15\,\mathrm{m}$ below the equilibrium position and released from rest.
\begin{enumerate}
    \item Find the equilibrium extension.
    \item Show that the motion is SHM and find its period.
    \item Find the greatest speed of the mass.
    \item Find the time taken for the mass to first reach a point $0.05\,\mathrm{m}$ above the equilibrium position.
\end{enumerate}
\end{exercice}

\begin{exercice}[Exercise 3]
A particle $P$ moves with SHM about centre $O$ with period $\frac{2\pi}{3}\,\mathrm{s}$. When $P$ is at a distance $0.4\,\mathrm{m}$ from $O$, its speed is $1.2\,\mathrm{m/s}$. Find the amplitude of the motion and the maximum speed.
\end{exercice}

\begin{corrige}[Exercise 1]
\begin{enumerate}
    \item Equilibrium at natural length $l=0.5$. Let $x$ be extension. $T = \frac{50x}{0.5}=100x$. $-T = 2\ddot{x} \Rightarrow \ddot{x} = -50x$. $\omega^2=50$, $\omega=5\sqrt{2}$.
    \item $T = \frac{2\pi}{5\sqrt{2}} = \frac{\pi\sqrt{2}}{5}\,\mathrm{s}$. Amplitude $a = 0.8-0.5 = 0.3\,\mathrm{m}$.
    \item $|\ddot{x}|_{\max} = a\omega^2 = 0.3 \times 50 = 15\,\mathrm{m/s^2}$.
    \item Spring length $0.6 \Rightarrow x=0.1$. $v^2 = 50(0.3^2-0.1^2) = 50 \times 0.08 = 4 \Rightarrow v = 2\,\mathrm{m/s}$.
\end{enumerate}
\end{corrige}

\begin{corrige}[Exercise 2]
\begin{enumerate}
    \item $mg = \frac{\lambda e}{l} \Rightarrow 0.5 \times 9.8 = 20e \Rightarrow e = 0.245\,\mathrm{m}$.
    \item $\ddot{x} = -\frac{\lambda}{ml}x = -\frac{20}{0.5 \times 1}x = -40x$. $\omega = \sqrt{40}=2\sqrt{10}$. $T = \frac{2\pi}{2\sqrt{10}} = \frac{\pi}{\sqrt{10}}\,\mathrm{s}$.
    \item $a = 0.15$. $v_{\max} = a\omega = 0.15 \times 2\sqrt{10} = 0.3\sqrt{10} \approx 0.949\,\mathrm{m/s}$.
    \item $x = a\cos\omega t$ (released from rest at $x=a$). $0.05 = 0.15\cos(2\sqrt{10}t) \Rightarrow \cos(2\sqrt{10}t) = 1/3 \Rightarrow 2\sqrt{10}t = \arccos(1/3) \approx 1.231 \Rightarrow t \approx 0.194\,\mathrm{s}$.
\end{enumerate}
\end{corrige}

\begin{corrige}[Exercise 3]
$T = \frac{2\pi}{3} \Rightarrow \omega = 3$. Given $x=0.4$, $v=1.2$. $v^2 = \omega^2(a^2-x^2) \Rightarrow 1.44 = 9(a^2-0.16) \Rightarrow a^2-0.16 = 0.16 \Rightarrow a^2 = 0.32 \Rightarrow a = 0.4\sqrt{2} \approx 0.566\,\mathrm{m}$. $v_{\max} = a\omega = 1.2\sqrt{2} \approx 1.697\,\mathrm{m/s}$.
\end{corrige}

\coursec{Further Results on SHM: Energy, Pendulums and Composite Problems}

In the previous section we established the definition and the differential equation of simple harmonic motion (SHM): $\ddot{x} = -\omega^2 x$. We also learned how to write the displacement as $x = a \cos(\omega t + \varepsilon)$ or $x = a \sin(\omega t + \varepsilon)$, and how to obtain the velocity $\dot{x} = \pm\omega\sqrt{a^2-x^2}$ and the acceleration. We now deepen our understanding by studying energy in SHM, the classic example of the simple pendulum, systems with two springs or strings, motions in which the SHM is interrupted (for example when a string goes slack), and the important technique of calculating the time taken to move between two points. These are precisely the tools that turn the abstract equation into a powerful model for GCE examination problems.

\courssub{Energy in Simple Harmonic Motion}

Consider a particle of mass $m$ attached to a light elastic string or spring of natural length $l$ and modulus of elasticity $\lambda$, performing SHM on a smooth horizontal table. The only force doing work is the elastic tension (or thrust), which is conservative. Therefore the total mechanical energy is conserved.

The kinetic energy is $KE = \frac{1}{2} m \dot{x}^2$. The elastic potential energy stored in a string or spring stretched by $x$ is $PE = \dfrac{\lambda x^2}{2l}$. Writing Hooke's law as $T = \dfrac{\lambda x}{l} = kx$ with stiffness $k = \lambda/l$, this becomes $PE = \frac{1}{2} k x^2$.

From the SHM equation $\ddot{x} = -\dfrac{\lambda}{ml}x$ we read off $\omega^2 = \dfrac{\lambda}{ml}$ (equivalently $\omega^2 = \dfrac{k}{m}$). Hence $\lambda = m\omega^2 l$ and the potential energy measured from the centre of oscillation is $PE = \frac{1}{2} m \omega^2 x^2$.

\begin{propriete}[Conservation of Energy in SHM]
In SHM with amplitude $a$ and angular frequency $\omega$, the total mechanical energy $E$ is constant and is given by
\[
E = \frac{1}{2} m \omega^2 a^2.
\]
At displacement $x$ from the centre,
\[
KE = \frac{1}{2} m \omega^2 (a^2 - x^2), \qquad PE = \frac{1}{2} m \omega^2 x^2.
\]
Kinetic energy is maximum at the centre ($x=0$) and zero at the extremes ($x=\pm a$); potential energy is zero at the centre and maximum at the extremes.
\end{propriete}

\textbf{Proof (examinable as a ``show that'').} Take $x = a \cos(\omega t + \varepsilon)$, so $\dot{x} = -a\omega \sin(\omega t + \varepsilon)$. Using $\dot{x}^2 = \omega^2(a^2 - x^2)$,
\[
KE = \frac{1}{2} m \omega^2 (a^2 - x^2), \qquad PE = \frac{1}{2} m \omega^2 x^2.
\]
Adding,
\[
E = KE + PE = \frac{1}{2} m \omega^2 (a^2 - x^2) + \frac{1}{2} m \omega^2 x^2 = \frac{1}{2} m \omega^2 a^2,
\]
which is independent of $x$ and of $t$: the total energy is constant. Equivalently, substituting the trigonometric forms gives $E = \frac{1}{2} m \omega^2 a^2(\sin^2(\omega t+\varepsilon) + \cos^2(\omega t+\varepsilon)) = \frac{1}{2} m \omega^2 a^2$. \hfill$\blacksquare$

\begin{aretenir}[Energy Summary]
\begin{itemize}
\item Total energy $E = \frac{1}{2}m\omega^2 a^2$ (constant).
\item $KE = \frac{1}{2}m\omega^2(a^2-x^2)$; $PE = \frac{1}{2}m\omega^2 x^2$.
\item $KE$ and $PE$ each vary sinusoidally with time at \emph{twice} the frequency of the motion (period $T/2$), because they involve $\sin^2$ and $\cos^2$.
\item The average kinetic energy over one period equals the average potential energy, each being half the total: $\langle KE \rangle = \langle PE \rangle = \frac{1}{4} m \omega^2 a^2$.
\end{itemize}
\end{aretenir}

\begin{popfigure}
\begin{tikzpicture}
\begin{axis}[
    width=\linewidth, height=7cm,
    xlabel={$t$}, ylabel={Energy},
    xmin=0, xmax=6.28, ymin=0, ymax=1.15,
    xtick={0,1.57,3.14,4.71,6.28},
    xticklabels={$0$, $\frac{T}{4}$, $\frac{T}{2}$, $\frac{3T}{4}$, $T$},
    ytick={0,0.5,1},
    yticklabels={$0$, $\frac{E}{2}$, $E$},
    legend pos=north east,
    grid=major
]
\addplot[popblue, thick, domain=0:6.28, samples=200] {0.5*(1-cos(deg(2*x)))};
\addlegendentry{Kinetic energy $KE$}
\addplot[poporange, thick, domain=0:6.28, samples=200] {0.5*(1+cos(deg(2*x)))};
\addlegendentry{Potential energy $PE$}
\addplot[popgreen, thick, dashed, domain=0:6.28] {1};
\addlegendentry{Total energy $E$}
\end{axis}
\end{tikzpicture}
\legende{Variation of kinetic, potential and total energy with time over one period $T$ (particle released from $x=+a$ at $t=0$). Note that $KE$ and $PE$ oscillate at twice the frequency of the motion, while their sum stays constant.}
\end{popfigure}

\begin{exoresolu}[Energy in a vertical string]
A particle of mass $0.5\,\mathrm{kg}$ hangs in equilibrium from a light elastic string of natural length $1\,\mathrm{m}$ and modulus $20\,\mathrm{N}$. The particle is pulled down a further $0.2\,\mathrm{m}$ and released from rest. Find (a) the angular frequency $\omega$, (b) the total energy of the motion, (c) the maximum speed. (Take $g = 9.8\,\mathrm{m\,s^{-2}}$.)
\tcblower
\textbf{(a)} At equilibrium the tension balances the weight: $mg = \dfrac{\lambda e}{l}$, giving $e = \dfrac{mgl}{\lambda} = \dfrac{0.5 \times 9.8 \times 1}{20} = 0.245\,\mathrm{m}$. For vertical oscillations about the equilibrium position the equation of motion reduces to $\ddot{x} = -\dfrac{\lambda}{ml}x$ (the weight cancels the equilibrium tension), so
\[
\omega^2 = \frac{\lambda}{ml} = \frac{20}{0.5 \times 1} = 40 \quad\Rightarrow\quad \omega = 2\sqrt{10}\,\mathrm{rad\,s^{-1}}.
\]
\textbf{(b)} The amplitude is $a = 0.2\,\mathrm{m}$ (released from rest $0.2\,\mathrm{m}$ below equilibrium). Hence
\[
E = \frac{1}{2} m \omega^2 a^2 = \frac{1}{2} \times 0.5 \times 40 \times (0.2)^2 = 0.4\,\mathrm{J}.
\]
\textbf{(c)} The maximum speed occurs at the centre: $v_{\max} = a\omega = 0.2 \times 2\sqrt{10} = 0.4\sqrt{10} \approx 1.26\,\mathrm{m\,s^{-1}}$.
\end{exoresolu}

\courssub{The Simple Pendulum}

The simple pendulum is the historical icon of SHM. It consists of a particle (the bob) of mass $m$ attached to a light inextensible string of length $l$, swinging in a vertical plane about a fixed point.

\begin{popfigure}
\begin{tikzpicture}[scale=1.3]
% Pivot and vertical reference
\draw[fill=popdark] (0,0) circle (0.05);
\draw[dashed, popmuted] (0,0.2) -- (0,-2.6);
% Displaced pendulum
\draw[popblue, thick] (0,0) -- (1.2,-1.6) node[midway, above right] {$l$};
\draw[popblue, fill=popblue!30] (1.2,-1.6) circle (0.15);
% Angle
\draw[poppurple, ->] (0,-0.7) arc (-90:-53:0.7) node[midway, below right] {$\theta$};
% Forces on bob
\begin{scope}[shift={(1.2,-1.6)}]
\draw[->, poporange, thick] (0,0) -- (0,-1.1) node[right] {$mg$};
\draw[->, popblue, thick] (0,0) -- (-0.75,1.0) node[above left] {$T$};
\draw[dashed, popmuted] (0,0) -- (0.75,-1.0);
\draw[->, poppink, thick] (0,0) -- (0.75,-1.0) node[below right] {$mg\sin\theta$};
\end{scope}
% Arc of motion
\draw[popmuted, dashed] (0,-2.45) arc (-90:-53:2.45);
\end{tikzpicture}
\legende{Simple pendulum displaced through an angle $\theta$. The tension $T$ has no tangential component; the restoring force along the arc is $-mg\sin\theta$.}
\end{popfigure}

Let $\theta$ be the angular displacement from the vertical, measured in \textbf{radians}. The bob moves on a circular arc of radius $l$; its tangential acceleration is $l\ddot{\theta}$. The tension is perpendicular to the motion, so the only tangential force is the component of the weight, $-mg\sin\theta$ (negative because it acts towards $\theta = 0$). Newton's second law along the tangent gives
\[
m l \ddot{\theta} = -mg \sin\theta \quad \Rightarrow \quad \ddot{\theta} = -\frac{g}{l} \sin\theta.
\]
This is \emph{not} exactly SHM, because $\sin\theta \neq \theta$ in general. However, for \emph{small oscillations} (in practice $\theta \lesssim 10^\circ \approx 0.17\,\mathrm{rad}$) we may use $\sin\theta \approx \theta$, giving
\[
\ddot{\theta} \approx -\frac{g}{l}\,\theta,
\]
which is the SHM equation with $\omega^2 = \dfrac{g}{l}$.

\begin{propriete}[Period of a Simple Pendulum]
For small oscillations, a simple pendulum of length $l$ performs approximate SHM with period
\[
T = 2\pi \sqrt{\frac{l}{g}}.
\]
The period is independent of the mass $m$ and of the (small) amplitude: the oscillations are \emph{isochronous}.
\end{propriete}

\begin{attention}[Radians Only!]
The approximation $\sin\theta \approx \theta$ is \textbf{only valid when $\theta$ is measured in radians}. In degrees the derivative of $\sin\theta$ is not $\cos\theta$, and the whole derivation collapses. Always convert angles to radians before applying any small-angle approximation.
\end{attention}

\begin{savaistu}[The Seconds Pendulum]
A \emph{seconds pendulum} is one whose half-period (the time of one swing) is exactly one second, i.e. $T = 2\,\mathrm{s}$. Its length is $l = g/\pi^2$, about $0.99\,\mathrm{m}$ in Cameroon. Pendulum clocks based on this principle were the most accurate timekeepers in the world for nearly three centuries, and pendulum measurements were historically used to map the variation of $g$ over the Earth's surface.
\end{savaistu}

\begin{exoresolu}[Determining $g$ with a pendulum]
In an experiment to determine $g$, a student measures the time for 20 complete oscillations of a simple pendulum of length $0.800\,\mathrm{m}$ as $36.2\,\mathrm{s}$. Calculate the value of $g$ obtained, and suggest two reasons why it might differ from the accepted value.
\tcblower
The period is $T = \dfrac{36.2}{20} = 1.81\,\mathrm{s}$. From $T = 2\pi\sqrt{l/g}$,
\[
g = \frac{4\pi^2 l}{T^2} = \frac{4\pi^2 \times 0.800}{(1.81)^2} \approx 9.64\,\mathrm{m\,s^{-2}}.
\]
\textbf{Possible reasons for the discrepancy:}
\begin{enumerate}
\item The amplitude may have been too large; the true period increases slightly with amplitude, so the measured $T$ is too large and $g$ comes out too small.
\item Timing errors: reaction time in starting and stopping the stopwatch (reduced by timing many oscillations, but not eliminated).
\item The string has mass and the bob is not a point mass, so the system behaves as a compound pendulum with a slightly different period.
\end{enumerate}
\end{exoresolu}

\courssub{Oscillations with Two Springs or Elastic Strings}

Many GCE problems involve a particle attached to two springs or elastic strings on a smooth horizontal surface. The key step is always to find the \textbf{net restoring force} as a function of the displacement $x$ from the \emph{equilibrium} position.

\begin{popfigure}
\begin{tikzpicture}[scale=1.2]
\draw[popline, thick] (-3.4,-0.35) -- (3.4,-0.35);
% Left string (zigzag)
\draw[popblue, thick] (-3,0) -- (-2.6,0.12) -- (-2.2,-0.12) -- (-1.8,0.12) -- (-1.4,-0.12) -- (-1.0,0.12) -- (-0.6,0);
% Right string
\draw[popblue, thick] (0.6,0) -- (1.0,0.12) -- (1.4,-0.12) -- (1.8,0.12) -- (2.2,-0.12) -- (2.6,0.12) -- (3,0);
% Particle
\draw[fill=poporange, draw=popdark] (0,0) circle (0.22) node[below=0.35cm] {$m$};
% Equilibrium mark
\draw[dashed, popmuted] (0,-0.5) -- (0,0.6);
% Displacement
\draw[<->, poppurple] (0,0.45) -- (0.9,0.45) node[midway, above] {$x$};
% Labels
\node[above] at (-1.8,0.35) {$l_1,\ \lambda_1$};
\node[above] at (1.8,0.35) {$l_2,\ \lambda_2$};
% Fixed walls
\draw[fill=popdark] (-3.3,-0.35) rectangle (-3,0.5);
\draw[fill=popdark] (3,-0.35) rectangle (3.3,0.5);
\node[left] at (-3.3,0.1) {$A$};
\node[right] at (3.3,0.1) {$B$};
\end{tikzpicture}
\legende{Particle of mass $m$ between two light elastic strings fixed at $A$ and $B$ on a smooth table, displaced a distance $x$ to the right of equilibrium.}
\end{popfigure}

\textbf{Particle between two fixed points (the standard GCE configuration).}
Let the left string have natural length $l_1$ and modulus $\lambda_1$, the right string $l_2$ and $\lambda_2$, and let $AB = L$. At equilibrium the extensions $e_1, e_2$ satisfy
\[
(l_1 + e_1) + (l_2 + e_2) = L, \qquad \frac{\lambda_1 e_1}{l_1} = \frac{\lambda_2 e_2}{l_2} \quad \text{(tensions balance).}
\]
Displace the particle a distance $x$ to the right. The new extensions are $e_1 + x$ (left) and $e_2 - x$ (right). The net force, taking right as positive, is
\[
F = \frac{\lambda_2 (e_2 - x)}{l_2} - \frac{\lambda_1 (e_1 + x)}{l_1}
= \underbrace{\left(\frac{\lambda_2 e_2}{l_2} - \frac{\lambda_1 e_1}{l_1}\right)}_{=\,0\ \text{(equilibrium)}} - \left( \frac{\lambda_1}{l_1} + \frac{\lambda_2}{l_2} \right) x.
\]
The constant terms cancel — this always happens when displacement is measured from equilibrium. The motion is therefore SHM with
\[
\omega^2 = \frac{\lambda_1}{m l_1} + \frac{\lambda_2}{m l_2}.
\]
For identical strings ($\lambda_1 = \lambda_2 = \lambda$, $l_1 = l_2 = l$), $\omega^2 = \dfrac{2\lambda}{ml}$.

\begin{attention}[Strings can go slack]
An elastic \emph{string} (unlike a spring) cannot push: it exerts a force only when taut. The formula above is valid only while \emph{both} strings remain taut, i.e. for $|x| \le \min(e_1, e_2)$. If the amplitude exceeds the smaller equilibrium extension, one string goes slack during the motion and the problem becomes an ``interrupted motion'' problem (see below). Always check this condition.
\end{attention}

\begin{methode}[Analysing Two-String Systems]
\begin{enumerate}
\item Draw a clear diagram showing the equilibrium position and a displaced position $x$.
\item Find the equilibrium extensions $e_1, e_2$ using the geometry and the balance of tensions.
\item Write the tensions in the displaced position in terms of $x$.
\item Write the net force $F(x)$; the constant (equilibrium) terms must cancel.
\item Read off $\omega^2$ as the coefficient of $-x/m$, then $T = 2\pi/\omega$.
\item Check that both strings remain taut throughout the motion.
\end{enumerate}
\end{methode}

\begin{exoresolu}[Two strings on a smooth table]
A particle of mass $0.2\,\mathrm{kg}$ is attached to two light elastic strings of natural lengths $1\,\mathrm{m}$ and $1.5\,\mathrm{m}$ and moduli $10\,\mathrm{N}$ and $20\,\mathrm{N}$ respectively. The other ends are fixed to points $A$ and $B$ on a smooth horizontal table, where $AB = 4\,\mathrm{m}$. The particle is displaced $0.3\,\mathrm{m}$ from its equilibrium position and released from rest. Show that the motion is SHM and find its period.
\tcblower
Let $e_1, e_2$ be the equilibrium extensions. Geometry: $(1+e_1) + (1.5+e_2) = 4 \Rightarrow e_1 + e_2 = 1.5$.
Tension balance: $\dfrac{10 e_1}{1} = \dfrac{20 e_2}{1.5} \Rightarrow 10 e_1 = \dfrac{40}{3} e_2 \Rightarrow e_1 = \dfrac{4}{3} e_2$.
Substituting: $\dfrac{4}{3}e_2 + e_2 = 1.5 \Rightarrow \dfrac{7}{3}e_2 = 1.5 \Rightarrow e_2 = \dfrac{9}{14}\,\mathrm{m}$, $e_1 = \dfrac{6}{7}\,\mathrm{m}$.
Both strings remain taut since the amplitude $0.3 < \min(e_1, e_2) = 9/14 \approx 0.64\,\mathrm{m}$.
Displace $x$ to the right: extensions are $e_1 + x$ and $e_2 - x$, so
\[
T_1 = 10(e_1 + x), \qquad T_2 = \frac{20}{1.5}(e_2 - x) = \frac{40}{3}(e_2 - x).
\]
Net force (taking right as positive):
\[
F = T_2 - T_1 = \frac{40}{3}e_2 - 10 e_1 - \left(10 + \frac{40}{3}\right)x = -\frac{70}{3}x,
\]
since $\frac{40}{3}e_2 = 10e_1$ at equilibrium. Hence
\[
0.2\,\ddot{x} = -\frac{70}{3}x \quad\Rightarrow\quad \ddot{x} = -\frac{350}{3}x,
\]
which is SHM with $\omega^2 = \dfrac{350}{3}$, so $\omega = \sqrt{350/3} \approx 10.8\,\mathrm{rad\,s^{-1}}$ and
\[
T = \frac{2\pi}{\omega} = 2\pi\sqrt{\frac{3}{350}} \approx 0.58\,\mathrm{s}.
\]
\end{exoresolu}

\courssub{Interrupted Motion: Strings Going Slack}

A favourite GCE theme: the SHM is interrupted because a string goes slack, or the particle leaves a platform, or a collision occurs. The governing principle is:

\begin{aretenir}[Interrupted Motion Principle]
While the string is taut (or the spring is attached), the motion is SHM about the equilibrium position. When a string goes slack, its tension vanishes, the restoring force disappears, and the particle moves with \textbf{constant velocity} (horizontal motion) or as a \textbf{projectile / under gravity alone} (vertical motion) until the string becomes taut again. The two phases must be analysed separately.
\end{aretenir}

\begin{attention}[Slack string means no SHM]
The instant a string goes slack, $\ddot{x} = -\omega^2 x$ ceases to apply. Do not continue to use SHM formulas (such as $v = \omega\sqrt{a^2-x^2}$) beyond the slack point. Use the velocity at the slack point as the initial velocity for the free phase.
\end{attention}

\begin{exoresolu}[Horizontal string going slack]
A particle $P$ of mass $0.5\,\mathrm{kg}$ is attached to one end of a light elastic string of natural length $1\,\mathrm{m}$ and modulus $20\,\mathrm{N}$. The other end is fixed to a point $O$ of a smooth horizontal table. The particle is held with the string stretched to a length $1.5\,\mathrm{m}$ and released from rest. Find (a) the time until the string first becomes slack, (b) the speed of $P$ at that instant, (c) the distance travelled by $P$ in the $0.2\,\mathrm{s}$ after the string goes slack.
\tcblower
Measure $x$ as the extension of the string. While taut, $\ddot{x} = -\dfrac{\lambda}{ml}x$ with
\[
\omega^2 = \frac{\lambda}{ml} = \frac{20}{0.5 \times 1} = 40 \quad\Rightarrow\quad \omega = 2\sqrt{10}\,\mathrm{rad\,s^{-1}}.
\]
The amplitude is $a = 0.5\,\mathrm{m}$ and, since $P$ is released from rest at maximum extension, $x = 0.5\cos(\omega t)$.
\textbf{(a)} The string goes slack when the extension first becomes zero: $\cos(\omega t) = 0 \Rightarrow \omega t = \dfrac{\pi}{2}$, so
\[
t = \frac{\pi}{2\omega} = \frac{\pi}{4\sqrt{10}} \approx 0.248\,\mathrm{s}.
\]
\textbf{(b)} $\dot{x} = -0.5\,\omega \sin(\omega t)$; at $\omega t = \pi/2$, $|\dot{x}| = 0.5\omega = \sqrt{10} \approx 3.16\,\mathrm{m\,s^{-1}}$.
\textbf{(c)} Once slack, there is no horizontal force on $P$ (smooth table), so it moves with constant speed $\sqrt{10}\,\mathrm{m\,s^{-1}}$. In $0.2\,\mathrm{s}$ it travels
\[
s = \sqrt{10} \times 0.2 \approx 0.63\,\mathrm{m}.
\]
\end{exoresolu}

\begin{exoresolu}[Vertical string going slack]
A particle of mass $0.2\,\mathrm{kg}$ hangs in equilibrium from a light elastic string of natural length $0.5\,\mathrm{m}$ and modulus $10\,\mathrm{N}$. It is pulled down a further $0.3\,\mathrm{m}$ and released from rest. Find the time from release until the string first becomes slack. (Take $g = 9.8\,\mathrm{m\,s^{-2}}$.)
\tcblower
The equilibrium extension is $e = \dfrac{mgl}{\lambda} = \dfrac{0.2 \times 9.8 \times 0.5}{10} = 0.098\,\mathrm{m}$.
About equilibrium, $\omega^2 = \dfrac{\lambda}{ml} = \dfrac{10}{0.2 \times 0.5} = 100$, so $\omega = 10\,\mathrm{rad\,s^{-1}}$.
Take $x$ downwards from equilibrium; then $x = 0.3\cos(10t)$ (released from rest at $x = +0.3$).
The total extension of the string is $e + x$, so the tension is
\[
T = \frac{\lambda}{l}(e + x) = 20(0.098 + x).
\]
The string goes slack when $T = 0$, i.e. $x = -0.098\,\mathrm{m}$ ($0.098\,\mathrm{m}$ above equilibrium). Note this point is reached because the amplitude $0.3 > e = 0.098$.
Solve $0.3\cos(10t) = -0.098$: $\cos(10t) = -0.3267$, so the first positive solution is
\[
10t = \arccos(-0.3267) \approx 1.904\,\mathrm{rad} \quad\Rightarrow\quad t \approx 0.190\,\mathrm{s}.
\]
\end{exoresolu}

\begin{exoresolu}[One of two strings going slack]
A particle of mass $0.1\,\mathrm{kg}$ is attached to two identical light elastic strings, each of natural length $0.5\,\mathrm{m}$ and modulus $5\,\mathrm{N}$. The other ends are fixed to points $A$ and $B$ on a smooth horizontal table with $AB = 2\,\mathrm{m}$. The particle is displaced $0.4\,\mathrm{m}$ from the centre towards $B$ and released from rest. Show that one string goes slack during the motion, and find the speed of the particle at the instant of slackening.
\tcblower
By symmetry the equilibrium position is the midpoint, where each string has length $1\,\mathrm{m}$ and extension $e = 0.5\,\mathrm{m}$.
While both strings are taut, displacement $x$ from the centre gives (identical strings)
\[
\omega^2 = \frac{2\lambda}{ml} = \frac{2 \times 5}{0.1 \times 0.5} = 200 \quad\Rightarrow\quad \omega = 10\sqrt{2}\,\mathrm{rad\,s^{-1}}.
\]
The amplitude is $a = 0.4\,\mathrm{m}$. A string goes slack when its extension vanishes, i.e. when $|x| = e = 0.5\,\mathrm{m}$. Since $a = 0.4 < 0.5$, \emph{neither} string goes slack: the motion is pure SHM throughout, and the premise fails.
\medskip

\textbf{Corrected data:} suppose instead the particle is displaced $0.6\,\mathrm{m}$ towards $B$. Then $a = 0.6 > e = 0.5$, so the string attached to $A$ goes slack when $x = +0.5\,\mathrm{m}$ (moving towards $B$). Using $v^2 = \omega^2(a^2 - x^2)$,
\[
v^2 = 200\,(0.6^2 - 0.5^2) = 200 \times 0.11 = 22 \quad\Rightarrow\quad v = \sqrt{22} \approx 4.69\,\mathrm{m\,s^{-1}}.
\]
Thereafter the particle continues under the tension of the $B$-string alone, which is a new SHM about a new centre — a typical extension for the final part of such a question.
\end{exoresolu}

\courssub{Time Taken to Travel Between Two Points}

A direct and very common examination task: given two displacements $x_1$ and $x_2$ measured from the centre of oscillation, find the time taken to travel from one to the other.

\begin{methode}[Time from $x_1$ to $x_2$ in SHM]
\begin{enumerate}
\item Choose the form of the solution to match the initial condition: $x = a\cos(\omega t)$ if the particle starts at an extreme ($x = \pm a$ at $t=0$); $x = a\sin(\omega t)$ if it starts at the centre moving positively; otherwise $x = a\cos(\omega t + \varepsilon)$.
\item Find the phase angles $\theta_1, \theta_2$ corresponding to $x_1$ and $x_2$, choosing the quadrants consistent with the \emph{direction of motion} (the phase increases steadily with time).
\item The time taken is $\Delta t = \dfrac{|\theta_2 - \theta_1|}{\omega}$.
\end{enumerate}
\end{methode}

\begin{exoresolu}[Time between two points]
A particle performs SHM with amplitude $0.4\,\mathrm{m}$ and period $2\,\mathrm{s}$. It is released from rest at the positive extreme. Find the time taken to travel (a) from $x = 0.4$ to $x = 0.2$, (b) from $x = 0.2$ to $x = -0.2$, (c) from $x = 0.2$ to the centre.
\tcblower
$\omega = \dfrac{2\pi}{T} = \pi\,\mathrm{rad\,s^{-1}}$ and $x = 0.4\cos(\pi t)$, with the particle moving towards the centre for $0 < t < 1$.
\textbf{(a)} $x_1 = 0.4 \Rightarrow \pi t_1 = 0$; $x_2 = 0.2 \Rightarrow \cos(\pi t_2) = 0.5 \Rightarrow \pi t_2 = \dfrac{\pi}{3}$.
Time $= \dfrac{\pi/3 - 0}{\pi} = \dfrac{1}{3}\,\mathrm{s}$.
\textbf{(b)} $\pi t_1 = \dfrac{\pi}{3}$; $x_2 = -0.2 \Rightarrow \cos(\pi t_2) = -0.5 \Rightarrow \pi t_2 = \dfrac{2\pi}{3}$ (still in the first half-oscillation, moving negatively).
Time $= \dfrac{2\pi/3 - \pi/3}{\pi} = \dfrac{1}{3}\,\mathrm{s}$.
\textbf{(c)} $\pi t_1 = \dfrac{\pi}{3}$; $x_2 = 0 \Rightarrow \pi t_2 = \dfrac{\pi}{2}$.
Time $= \dfrac{\pi/2 - \pi/3}{\pi} = \dfrac{1}{6}\,\mathrm{s}$.
\medskip

Note the check: (b) $=$ (a) $+$ 2$\times$(c)? Indeed $\tfrac13 = \tfrac13$? No — the correct decomposition is (b) $=$ (c) $+$ time from centre to $-0.2$, and by symmetry of the cosine the time from $0.2$ to $0$ equals the time from $0$ to $-0.2$, so (b) $= 2\times$(c) $= \tfrac13\,\mathrm{s}$. Consistent.
\end{exoresolu}

\courssub{Approximate Simple Harmonic Motion}

Many systems are not \emph{exactly} simple harmonic but become so for small displacements. The recipe is always the same:

\begin{methode}[Proving Approximate SHM]
\begin{enumerate}
\item Identify the equilibrium position.
\item Displace the system by a small amount $x$ (or angle $\theta$).
\item Write the exact equation of motion (Newton's second law, or torque $= I\ddot{\theta}$).
\item Linearise using the small-quantity approximations: $\sin\theta \approx \theta$, $\tan\theta \approx \theta$, $\cos\theta \approx 1$ (to first order), $(1+x)^n \approx 1 + nx$.
\item Show the equation reduces to $\ddot{x} = -\omega^2 x$ (or $\ddot{\theta} = -\omega^2\theta$) and identify $\omega^2$.
\end{enumerate}
\end{methode}

\begin{exoresolu}[Bead on a smooth circular wire]
A bead of mass $m$ is threaded on a smooth circular wire of radius $R$ fixed in a vertical plane. The bead is slightly displaced from the lowest point. Show that it performs approximate SHM and find the period.
\tcblower
Let $\theta$ be the angular displacement from the downward vertical. The tangential component of the weight is $-mg\sin\theta$ and the tangential acceleration is $R\ddot{\theta}$, so
\[
mR\ddot{\theta} = -mg\sin\theta \quad\Rightarrow\quad \ddot{\theta} = -\frac{g}{R}\sin\theta.
\]
For small $\theta$ (radians), $\sin\theta \approx \theta$: $\ddot{\theta} \approx -\dfrac{g}{R}\theta$ — SHM with $\omega^2 = \dfrac{g}{R}$, hence
\[
T = 2\pi\sqrt{\frac{R}{g}}.
\]
This is identical to a simple pendulum of length $R$: the wire supplies the constraint that the string supplied before.
\end{exoresolu}

\begin{exoresolu}[Floating cylinder]
A uniform solid cylinder of cross-sectional area $A$, height $h$ and density $\rho$ floats vertically in a liquid of density $\sigma$ ($\rho < \sigma$). It is pushed down a small distance $x$ and released. Show that the motion is SHM and find the period.
\tcblower
At equilibrium, weight $=$ upthrust (Archimedes): $\rho A h g = \sigma A h_0 g$, where $h_0$ is the submerged depth, so $h_0 = \dfrac{\rho}{\sigma}h$.
Push the cylinder down by $x$ (small enough that it remains partially immersed and not fully submerged). The submerged depth becomes $h_0 + x$, so the upthrust is $\sigma A (h_0 + x) g$. Taking downwards as positive, the net force is
\[
F = \underbrace{\rho A h g - \sigma A h_0 g}_{=\,0\ \text{(equilibrium)}} - \sigma A g\, x = -\sigma A g\, x.
\]
With mass $m = \rho A h$,
\[
\rho A h\, \ddot{x} = -\sigma A g\, x \quad\Rightarrow\quad \ddot{x} = -\frac{\sigma g}{\rho h} x,
\]
SHM with $\omega^2 = \dfrac{\sigma g}{\rho h}$, so $T = 2\pi\sqrt{\dfrac{\rho h}{\sigma g}}$.
\end{exoresolu}

\begin{attention}[Approximate SHM Checklist]
When asked to ``show that the motion is approximately simple harmonic'':
\begin{itemize}
\item State explicitly the approximation used (e.g. $\sin\theta \approx \theta$, $\theta$ in radians).
\item Show that the restoring force (or torque) is proportional to the displacement and directed towards equilibrium.
\item Identify $\omega^2$ clearly and state the period if required.
\item Mention that the result requires the amplitude to be small.
\end{itemize}
\end{attention}

\begin{aretenir}[Key Formulas for this Section]
\begin{itemize}
\item Energy: $E = \frac{1}{2}m\omega^2 a^2$, $KE = \frac{1}{2}m\omega^2(a^2-x^2)$, $PE = \frac{1}{2}m\omega^2 x^2$.
\item Simple pendulum (small oscillations): $T = 2\pi\sqrt{\dfrac{l}{g}}$.
\item Particle between two strings/springs: $\omega^2 = \dfrac{\lambda_1}{ml_1} + \dfrac{\lambda_2}{ml_2}$ (both taut).
\item Time from $x_1$ to $x_2$: match phase angles in $x = a\cos(\omega t)$ or $a\sin(\omega t)$; $\Delta t = \Delta\theta/\omega$.
\item Slack string $\Rightarrow$ constant velocity (horizontal) or motion under gravity alone (vertical).
\item Approximate SHM: linearise with $\sin\theta \approx \theta$, $\tan\theta \approx \theta$, $(1+x)^n \approx 1 + nx$.
\end{itemize}
\end{aretenir}

\begin{exercice}[GCE-Style Practice]
\begin{enumerate}
\item A particle of mass $0.3\,\mathrm{kg}$ is attached to a light elastic string of natural length $0.8\,\mathrm{m}$ and modulus $12\,\mathrm{N}$ on a smooth horizontal table. The particle is pulled until the extension is $0.4\,\mathrm{m}$ and released from rest. (a) Show that the motion is SHM and find $\omega$. (b) Find the total energy. (c) Find the speed when the extension is $0.1\,\mathrm{m}$. (d) Find the time taken to go from the point of release to the point where the extension is $0.1\,\mathrm{m}$.
\item A simple pendulum of length $1.2\,\mathrm{m}$ makes small oscillations where $g = 9.8\,\mathrm{m\,s^{-2}}$. (a) Find its period. (b) The pendulum is taken to a planet where $g = 4.0\,\mathrm{m\,s^{-2}}$. By what factor does the period change? (c) If the amplitude is $5^\circ$, is the small-angle approximation justified? (Use $\sin 5^\circ \approx 0.0872$ and $5^\circ \approx 0.0873\,\mathrm{rad}$.)
\item A particle of mass $m$ is attached to two identical light elastic strings, each of natural length $l$ and modulus $\lambda$. The other ends are fixed to points $A$ and $B$ on a smooth horizontal table with $AB = 4l$. The particle is held at the midpoint of $AB$ and released from rest. Show that the motion is SHM and find the period in terms of $m, \lambda, l$.
\item A uniform rod of length $2L$ and mass $M$ is pivoted smoothly at one end and hangs vertically. It is displaced through a small angle and released. By considering the torque about the pivot, show that the rod performs approximate SHM and find its period. (Moment of inertia of a uniform rod of length $2L$ about one end: $I = \frac{4}{3}ML^2$; the centre of mass is at distance $L$ from the pivot.)
\item A cylindrical buoy of mass $m$ and cross-sectional area $A$ floats vertically in water of density $\rho$. It is pushed down a small distance and released. Show that it performs SHM and find the period. If a bird of mass $m_b$ lands gently on the buoy, how does the period change?
\end{enumerate}
\end{exercice}

\begin{corrige}[Exercise 1]
(a) With extension $x$, the only horizontal force is the tension $T = \dfrac{\lambda x}{l}$ directed towards the natural-length position, so $m\ddot{x} = -\dfrac{\lambda}{l}x$, i.e. $\ddot{x} = -\dfrac{\lambda}{ml}x$: SHM. $\omega^2 = \dfrac{12}{0.3 \times 0.8} = 50$, $\omega = 5\sqrt{2} \approx 7.07\,\mathrm{rad\,s^{-1}}$.
(b) $E = \frac{1}{2}m\omega^2 a^2 = \frac{1}{2} \times 0.3 \times 50 \times (0.4)^2 = 1.2\,\mathrm{J}$.
(c) $v^2 = \omega^2(a^2 - x^2) = 50(0.16 - 0.01) = 7.5 \Rightarrow v = \sqrt{7.5} \approx 2.74\,\mathrm{m\,s^{-1}}$.
(d) $x = 0.4\cos(\omega t)$. At $x = 0.1$: $\cos(\omega t) = 0.25 \Rightarrow \omega t = \arccos(0.25) \approx 1.318\,\mathrm{rad}$, so $t = \dfrac{1.318}{5\sqrt{2}} \approx 0.186\,\mathrm{s}$.
\end{corrige}

\begin{corrige}[Exercise 2]
(a) $T = 2\pi\sqrt{\dfrac{1.2}{9.8}} \approx 2.20\,\mathrm{s}$.
(b) $T \propto 1/\sqrt{g}$, so the period is multiplied by $\sqrt{9.8/4.0} \approx 1.57$ (it increases).
(c) $\sin 5^\circ \approx 0.0872$ while $\theta = 0.0873\,\mathrm{rad}$: a relative error of about $0.1\%$. The approximation is excellent.
\end{corrige}

\begin{corrige}[Exercise 3]
At the midpoint each string has length $2l$ and extension $l$; the tensions are equal, so the midpoint is the equilibrium position. Displace $x$ towards $B$: extensions are $l + x$ and $l - x$ (both positive for $|x| < l$, so both strings stay taut).
\[
T_L = \frac{\lambda}{l}(l + x), \qquad T_R = \frac{\lambda}{l}(l - x).
\]
Net force towards the centre: $F = T_R - T_L = -\dfrac{2\lambda}{l}x$. Hence $\ddot{x} = -\dfrac{2\lambda}{ml}x$: SHM with $\omega^2 = \dfrac{2\lambda}{ml}$ and
\[
T = 2\pi\sqrt{\frac{ml}{2\lambda}}.
\]
\end{corrige}

\begin{corrige}[Exercise 4]
At angular displacement $\theta$, the weight $Mg$ acts at the centre of mass, distance $L$ from the pivot, giving a restoring torque $\tau = -MgL\sin\theta \approx -MgL\theta$ for small $\theta$. With $I = \frac{4}{3}ML^2$,
\[
\frac{4}{3}ML^2\,\ddot{\theta} = -MgL\,\theta \quad\Rightarrow\quad \ddot{\theta} = -\frac{3g}{4L}\theta.
\]
SHM with $\omega^2 = \dfrac{3g}{4L}$, so $T = 2\pi\sqrt{\dfrac{4L}{3g}}$.
\end{corrige}

\begin{corrige}[Exercise 5]
At equilibrium $mg = \rho A h_0 g$, giving submerged depth $h_0 = \dfrac{m}{\rho A}$. Pushed down by $x$, the extra upthrust is $\rho A g x$ upwards, so $m\ddot{x} = -\rho A g\, x$: SHM with $\omega^2 = \dfrac{\rho A g}{m}$ and
\[
T = 2\pi\sqrt{\frac{m}{\rho A g}}.
\]
With the bird, the mass becomes $m + m_b$ while $\rho, A, g$ are unchanged, so
\[
T' = 2\pi\sqrt{\frac{m + m_b}{\rho A g}} = T\sqrt{\frac{m + m_b}{m}} > T:
\]
the period increases.
\end{corrige}

\coursec{Damped Harmonic Motion}

In the previous section we studied simple harmonic motion in its ideal form: a mass on a spring, a pendulum swinging in a vacuum, oscillations that continue forever with constant amplitude because no energy is ever lost. But walk through any mechanic's workshop in Douala or Bamenda and push a car down on its suspension: the car bounces a few times and then settles. A pendulum clock eventually stops unless it is wound. A plucked guitar string fades into silence. Real oscillators \emph{lose energy}, and their motion dies away. This section builds the mathematics of that reality: \cle{damped harmonic motion}. It is one of the most rewarding parts of the syllabus, because the same differential equation describes car suspensions, earthquake-resistant buildings and the needle of an analogue voltmeter.

\courssub{The Physical Origin of Damping}

When a body oscillates in a real medium, it experiences a \cle{resistive force} opposing its motion. The sources are many:

\begin{itemize}
\item \textbf{Air resistance:} a swinging pendulum pushes air aside and is slowed by it.
\item \textbf{Fluid friction:} a \cle{dashpot} (a piston moving in a cylinder of oil, exactly the device inside a car's \cle{shock absorber}) converts mechanical energy into heat.
\item \textbf{Internal friction:} the repeated flexing of a spring warms it slightly.
\end{itemize}

For the moderate speeds met in oscillating systems, experiments show that the resistance is well modelled as \emph{proportional to the velocity and directed opposite to it}:
\[
R = -\lambda \dot{x}, \qquad \lambda > 0,
\]
where $\lambda$ is the \cle{damping coefficient} (in $\mathrm{kg\,s^{-1}}$). The minus sign is essential: the resistance always acts against the direction of motion, whatever that direction is.

\begin{popfigure}
\begin{center}
\begin{tikzpicture}[scale=1.0, every node/.style={font=\small}]
% wall
\fill[poppurple!30] (-0.3,-1.2) rectangle (0,1.6);
\draw[thick] (0,-1.2) -- (0,1.6);
\foreach \y in {-1.1,-0.7,...,1.5}{\draw (0,\y) -- (-0.28,\y+0.18);}
% spring
\draw[thick, popblue, decorate, decoration={coil, segment length=5pt, amplitude=4pt}] (0,0.6) -- (3.2,0.6);
\node[popblue, above] at (1.6,0.85) {spring, stiffness $k$};
% dashpot
\draw[thick, popdark] (0,-0.6) -- (2.0,-0.6);
\draw[thick, popdark] (2.0,-0.95) rectangle (3.2,-0.25);
\draw[thick, popdark] (3.2,-0.6) -- (4.2,-0.6);
\fill[popblueL] (2.05,-0.9) rectangle (3.15,-0.3);
\node[popdark, below] at (2.6,-1.0) {dashpot: resistance $\lambda\dot{x}$};
% mass
\fill[poporange!80] (3.2,0.0) rectangle (4.2,1.2);
\node[white] at (3.7,0.6) {$m$};
% displacement arrow
\draw[->, thick, popgreen] (4.4,0.6) -- (5.4,0.6) node[right] {$x$};
% ground
\draw[thick] (-0.3,-1.2) -- (5.6,-1.2);
\foreach \x in {-0.2,0.3,...,5.5}{\draw (\x,-1.2) -- (\x+0.2,-1.45);}
\end{tikzpicture}
\end{center}
\legende{A mass $m$ attached to a spring of stiffness $k$ and a dashpot providing resistance $\lambda\dot{x}$.}
\end{popfigure}

\courssub{Setting Up the Equation of Motion}

Consider a particle of mass $m$ attached to a spring of natural length $l$ and modulus of elasticity $\lambda_e$ (stiffness $k = \lambda_e/l$), moving horizontally on a smooth table, subject to a resistance $\lambda\dot{x}$. Measuring $x$ from the \emph{equilibrium} position (so that the weight and any static extension cancel, exactly as in our SHM work), Newton's second law gives
\[
m\ddot{x} = -kx - \lambda\dot{x}.
\]
Dividing by $m$ and renaming the constants, we obtain the standard form.

\begin{definition}[Damped harmonic motion]
A particle performs \cle{damped harmonic motion} when its displacement $x$ from equilibrium satisfies
\[
\ddot{x} + 2k\dot{x} + \omega^{2}x = 0,
\]
where $\omega^{2} = \dfrac{k_{\text{spring}}}{m}$ is the square of the \cle{natural angular frequency} (the frequency the system \emph{would} have without damping) and $k = \dfrac{\lambda}{2m} > 0$ measures the \cle{strength of the damping}. When $k = 0$ the equation reduces to simple harmonic motion, $\ddot{x} + \omega^{2}x = 0$.
\end{definition}

\begin{attention}[Notation alert]
In the standard form $\ddot{x} + 2k\dot{x} + \omega^{2}x = 0$, the letter $k$ denotes the \emph{damping parameter} $\lambda/(2m)$, \textbf{not} the spring stiffness. Some textbooks write $2n$ or $2\gamma$ instead of $2k$. Always check the notation of the question before quoting formulae, and state your own notation clearly at the start of your solution.
\end{attention}

\courssub{Solving the Equation: the Auxiliary Equation}

The equation is linear with constant coefficients, so we try $x = e^{rt}$. Substituting and dividing by $e^{rt} \neq 0$ gives the \cle{auxiliary equation}
\[
r^{2} + 2kr + \omega^{2} = 0,
\qquad\text{with roots}\qquad
r = -k \pm \sqrt{k^{2} - \omega^{2}}.
\]
Everything now hinges on the sign of the discriminant $k^{2} - \omega^{2}$, i.e.\ on how the damping $k$ compares with the natural frequency $\omega$.

\begin{propriete}[The three regimes of damping — proof examinable]
The nature of the motion governed by $\ddot{x} + 2k\dot{x} + \omega^{2}x = 0$ depends on the sign of $k^{2} - \omega^{2}$:
\begin{itemize}
\item $k < \omega$ (\cle{light} or \emph{under} damping): complex roots $-k \pm i\Omega$; the motion \emph{oscillates} with decaying amplitude.
\item $k = \omega$ (\cle{critical} damping): repeated root $-k$; the motion returns to equilibrium \emph{as fast as possible without oscillating}.
\item $k > \omega$ (\cle{heavy} or \emph{over} damping): two real negative roots; the motion returns to equilibrium \emph{slowly, without oscillating}.
\end{itemize}
\end{propriete}

\courssub{Light Damping ($k < \omega$): Oscillations that Fade Away}

When $k < \omega$, write $\Omega = \sqrt{\omega^{2} - k^{2}} > 0$. The roots are $r = -k \pm i\Omega$, and the general solution is
\[
\boxed{\;x = e^{-kt}\bigl(A\cos\Omega t + B\sin\Omega t\bigr), \qquad \Omega = \sqrt{\omega^{2} - k^{2}}.\;}
\]
This is a beautiful formula: it is \emph{SHM of angular frequency $\Omega$ multiplied by a dying exponential}. The motion still crosses equilibrium regularly, but each swing is smaller than the last. The curve is trapped between the two \cle{exponential envelopes} $x = \pm a\,e^{-kt}$.

\begin{popfigure}
\begin{center}
\begin{tikzpicture}
\begin{axis}[
  width=0.85\linewidth, height=6.2cm,
  axis lines=middle, xlabel={$t$}, ylabel={$x$},
  xmin=0, xmax=13, ymin=-1.3, ymax=1.3,
  xtick=\empty, ytick=\empty,
  domain=0:12.5, samples=300,
  legend style={at={(0.98,0.95)}, anchor=north east, font=\small}
]
\addplot[popblue, thick] {exp(-0.25*x)*cos(deg(1.5*x))};
\addplot[poporange, dashed, thick] {exp(-0.25*x)};
\addplot[poporange, dashed, thick] {-exp(-0.25*x)};
\legend{$x=e^{-kt}\cos\Omega t$, $+e^{-kt}$, $-e^{-kt}$}
\end{axis}
\end{tikzpicture}
\end{center}
\legende{Lightly damped oscillation: the curve oscillates inside the dashed exponential envelopes $\pm a e^{-kt}$.}
\end{popfigure}

The damped motion is not strictly periodic (the amplitude changes), but successive passages through equilibrium in the same direction are separated by a constant time, the \cle{damped period}:
\[
T_{d} = \frac{2\pi}{\Omega} = \frac{2\pi}{\sqrt{\omega^{2} - k^{2}}}.
\]

\begin{attention}[The damped period is longer]
The period of a lightly damped oscillator is $\dfrac{2\pi}{\sqrt{\omega^{2}-k^{2}}}$, \textbf{not} $\dfrac{2\pi}{\omega}$. Since $\sqrt{\omega^{2}-k^{2}} < \omega$, damping always \emph{lengthens} the period. Quoting $2\pi/\omega$ for a damped system is a classic GCE error.
\end{attention}

\begin{exemplebox}[Finding the constants for light damping]
A particle moves with $\ddot{x} + 2\dot{x} + 10x = 0$, with $x = 2$ and $\dot{x} = 0$ at $t = 0$. Here $k = 1$, $\omega^{2} = 10$, so $k^{2} - \omega^{2} = -9 < 0$: light damping with $\Omega = \sqrt{10-1} = 3$. Hence
\[
x = e^{-t}\bigl(A\cos 3t + B\sin 3t\bigr).
\]
From $x(0) = 2$: $A = 2$. Differentiating,
\[
\dot{x} = e^{-t}\bigl[(-3A\sin 3t + 3B\cos 3t) - (A\cos 3t + B\sin 3t)\bigr],
\]
so $\dot{x}(0) = 3B - A = 0$, giving $B = \tfrac{2}{3}$. Therefore
\[
x = e^{-t}\left(2\cos 3t + \tfrac{2}{3}\sin 3t\right), \qquad T_{d} = \frac{2\pi}{3}\ \mathrm{s}.
\]
\end{exemplebox}

\courssub{Critical Damping ($k = \omega$): the Fastest Non-Oscillatory Return}

When $k = \omega$ the auxiliary equation has the \emph{repeated} root $r = -k$. A repeated root of a linear equation always contributes solutions $e^{-kt}$ \emph{and} $t\,e^{-kt}$, so
\[
\boxed{\;x = (A + Bt)e^{-kt}.\;}
\]
The factor $e^{-kt}$ eventually dominates the linear factor $A + Bt$, so $x \to 0$. The particle returns to equilibrium without oscillating (it may cross equilibrium \emph{at most once}, when $A$ and $B$ have opposite signs), and it does so \emph{faster} than any over-damped system. This is exactly how a good car shock absorber is tuned: after hitting a pothole on the Yaound\'e--Bafoussam road, the wheel should settle immediately, not bounce.

\begin{attention}[GCE exam tip — the repeated root]
For critical damping the solution is $(A + Bt)e^{-kt}$. The extra factor $t$ comes from the \emph{repeated} root of the auxiliary equation. Writing only $Ae^{-kt}$ (a single exponential) is the most common error in this topic and costs nearly all the marks for the solution, because the general solution then has only one arbitrary constant and cannot satisfy two initial conditions.
\end{attention}

\courssub{Heavy Damping ($k > \omega$): a Slow Creep Back}

When $k > \omega$, the roots are real and distinct: $r = -k \pm \gamma$ with $\gamma = \sqrt{k^{2} - \omega^{2}} < k$, so \emph{both roots are negative}. The general solution is
\[
\boxed{\;x = e^{-kt}\bigl(Ae^{\gamma t} + Be^{-\gamma t}\bigr) = Ae^{-(k-\gamma)t} + Be^{-(k+\gamma)t}, \qquad \gamma = \sqrt{k^{2}-\omega^{2}}.\;}
\]
Both exponentials decay; there is \emph{no oscillation at all}. The slowest-decaying term, $e^{-(k-\gamma)t}$, governs the long-term behaviour, and since $k - \gamma$ is small when damping is heavy, the return to equilibrium is \emph{slower} than in the critical case. A door-closer that is too stiff is over-damped: the door creeps shut frustratingly slowly.

\begin{popfigure}
\begin{center}
\begin{tikzpicture}
\begin{axis}[
  width=0.85\linewidth, height=6.5cm,
  axis lines=middle, xlabel={$t$}, ylabel={$x$},
  xmin=0, xmax=10, ymin=-0.2, ymax=1.25,
  xtick=\empty, ytick=\empty,
  domain=0:10, samples=250,
  legend style={at={(0.98,0.95)}, anchor=north east, font=\small}
]
% under-damped: k=0.4, omega=1 => Omega=sqrt(1-0.16)=0.9165, x0=1,v0=0 => A=1,B=k/Omega=0.4364
\addplot[popblue, thick] {exp(-0.4*x)*(cos(deg(0.9165*x)) + 0.4364*sin(deg(0.9165*x)))};
% critical: k=1, x=(1+t)e^{-t}
\addplot[poporange, thick] {(1+x)*exp(-x)};
% over-damped: k=2, omega=1 => gamma=sqrt(3)=1.732; A=(v0+(k+gamma)x0)/(2 gamma)= (2+1.732)/3.464=1.0774, B=1-A=-0.0774
\addplot[popgreen, thick] {1.0774*exp(-0.2679*x) - 0.0774*exp(-3.732*x)};
\legend{light ($k<\omega$), critical ($k=\omega$), heavy ($k>\omega$)}
\end{axis}
\end{tikzpicture}
\end{center}
\legende{The three regimes for the same initial conditions $x(0)=1$, $\dot{x}(0)=0$ (with $\omega = 1$). Critical damping returns to equilibrium fastest without oscillating.}
\end{popfigure}

\courssub{The Logarithmic Decrement: Measuring Damping Experimentally}

For a lightly damped oscillator, consider two \emph{successive} maximum displacements on the same side, occurring one damped period $T_{d}$ apart. Since the cosine factor repeats while the exponential decays,
\[
\frac{x_{n}}{x_{n+1}} = \frac{e^{-kt}}{e^{-k(t+T_{d})}} = e^{kT_{d}},
\qquad\text{so}\qquad
\ln\!\left(\frac{x_{n}}{x_{n+1}}\right) = kT_{d}.
\]

\begin{definition}[Logarithmic decrement]
The \cle{logarithmic decrement} of a lightly damped oscillation is
\[
\Lambda = \ln\!\left(\frac{x_{n}}{x_{n+1}}\right) = kT_{d} = \frac{2\pi k}{\sqrt{\omega^{2}-k^{2}}},
\]
the natural logarithm of the ratio of two successive amplitudes on the same side. It is constant, so measuring any two successive amplitudes experimentally yields the damping constant $k$.
\end{definition}

This is a GCE-useful extension: given a table of experimental amplitudes, compute $\ln(x_n/x_{n+1})$ for each pair, average, and deduce $k$. Note also the practical trick $\ln(x_0/x_n) = n\Lambda$: using amplitudes many periods apart reduces measurement error.

\courssub{Method and Worked Examination Problem}

\begin{methode}[Solving a damped harmonic motion problem]
\begin{enumerate}
\item \textbf{Set up the equation} from Newton's second law (resistance $-\lambda\dot{x}$, restoring force $-k_{\text{spring}}x$) and reduce it to the standard form $\ddot{x} + 2k\dot{x} + \omega^{2}x = 0$, identifying $k$ and $\omega^{2}$.
\item \textbf{Write the auxiliary equation} $r^{2} + 2kr + \omega^{2} = 0$ and \textbf{compute $k^{2} - \omega^{2}$} to identify the regime.
\item \textbf{Write the standard solution} for that regime:
\[
\begin{cases}
k<\omega: & x = e^{-kt}(A\cos\Omega t + B\sin\Omega t),\ \Omega=\sqrt{\omega^{2}-k^{2}},\\[2pt]
k=\omega: & x = (A+Bt)e^{-kt},\\[2pt]
k>\omega: & x = Ae^{-(k-\gamma)t} + Be^{-(k+\gamma)t},\ \gamma=\sqrt{k^{2}-\omega^{2}}.
\end{cases}
\]
\item \textbf{Apply the initial conditions} $x(0)$ and $\dot{x}(0)$ to find $A$ and $B$ (differentiate carefully using the product rule).
\item \textbf{Answer the question asked}: damped period, logarithmic decrement, number of crossings, time to halve the amplitude, etc.
\end{enumerate}
\end{methode}

\begin{exoresolu}[A complete GCE-style problem]
A particle $P$ of mass $0.5\ \mathrm{kg}$ is attached to one end of a light elastic spring of stiffness $8\ \mathrm{N\,m^{-1}}$, the other end being fixed. $P$ moves in a straight line on a smooth horizontal table and experiences a resistance of magnitude $4v$ newtons, where $v$ is its speed. The particle is pulled $0.3\ \mathrm{m}$ from equilibrium and released from rest.
\smallskip

(i) Show that the displacement $x$ satisfies $\ddot{x} + 8\dot{x} + 16x = 0$.

(ii) Solve the equation and describe the motion.

(iii) Find the time at which the speed is greatest, and state whether $P$ crosses the equilibrium position.
\tcblower
\textbf{(i)} Newton's second law with restoring force $-8x$ and resistance $-4\dot{x}$:
\[
0.5\,\ddot{x} = -8x - 4\dot{x}
\;\Longrightarrow\;
\ddot{x} + 8\dot{x} + 16x = 0. \qquad\blacksquare
\]

\textbf{(ii)} Standard form: $2k = 8 \Rightarrow k = 4$; $\omega^{2} = 16 \Rightarrow \omega = 4$. Since $k = \omega$, the damping is \textbf{critical}. Auxiliary equation: $r^{2} + 8r + 16 = (r+4)^{2} = 0$, repeated root $r = -4$. General solution:
\[
x = (A + Bt)e^{-4t}.
\]
Initial conditions: $x(0) = 0.3 \Rightarrow A = 0.3$. Differentiate:
\[
\dot{x} = Be^{-4t} - 4(A+Bt)e^{-4t}, \qquad \dot{x}(0) = B - 4A = 0 \Rightarrow B = 1.2.
\]
\[
\boxed{\;x = (0.3 + 1.2t)e^{-4t}\;}
\]
The motion is \emph{non-oscillatory}: $P$ moves directly towards equilibrium, $x \to 0$ as $t \to \infty$.

\textbf{(iii)} Since $0.3 + 1.2t > 0$ for all $t \ge 0$, $x$ never changes sign: $P$ \textbf{never crosses} the equilibrium position. The speed is greatest when $\ddot{x} = 0$. From the equation, $\ddot{x} = -8\dot{x} - 16x$; setting $\ddot{x} = 0$:
\[
\dot{x} = (1.2 - 4(0.3+1.2t))e^{-4t} = -4.8t\,e^{-4t},
\]
so $\ddot{x} = -4.8e^{-4t} + 19.2t\,e^{-4t} = 0 \Rightarrow t = \dfrac{4.8}{19.2} = 0.25\ \mathrm{s}$.
The greatest speed is $|\dot{x}(0.25)| = 4.8(0.25)e^{-1} = 1.2e^{-1} \approx 0.44\ \mathrm{m\,s^{-1}}$.
\end{exoresolu}

\begin{exercice}[Light damping and the logarithmic decrement]
A particle performs damped harmonic motion with $\ddot{x} + 2\dot{x} + 26x = 0$. At $t = 0$, $x = 0.1\ \mathrm{m}$ and $\dot{x} = 0$.
\begin{enumerate}
\item Identify the regime and find $x$ in terms of $t$.
\item Find the damped period $T_{d}$.
\item Show that the ratio of successive maximum displacements is $e^{\pi/5}$, and find after how many complete oscillations the amplitude first falls below $0.01\ \mathrm{m}$.
\end{enumerate}
\end{exercice}

\begin{corrige}[Exercise 1]
\textbf{1.} $k = 1$, $\omega^{2} = 26$; $k^{2} - \omega^{2} = -25 < 0$: light damping, $\Omega = \sqrt{26-1} = 5$. So $x = e^{-t}(A\cos 5t + B\sin 5t)$. $x(0) = 0.1 \Rightarrow A = 0.1$; $\dot{x}(0) = 5B - A = 0 \Rightarrow B = 0.02$. Hence $x = e^{-t}(0.1\cos 5t + 0.02\sin 5t)$.

\textbf{2.} $T_{d} = \dfrac{2\pi}{5}\ \mathrm{s} \approx 1.26\ \mathrm{s}$.

\textbf{3.} Ratio $= e^{kT_{d}} = e^{2\pi/5}$... note $kT_d = 1 \times \tfrac{2\pi}{5} = \tfrac{2\pi}{5}$, so the ratio is $e^{2\pi/5}$. Amplitude after $n$ periods: $a_n = 0.1\,e^{-nT_d} = 0.1e^{-2\pi n/5}$. We need $e^{-2\pi n/5} < 0.1$, i.e.\ $\tfrac{2\pi n}{5} > \ln 10 \approx 2.303$, giving $n > \dfrac{5 \times 2.303}{2\pi} \approx 1.83$. So the amplitude first falls below $0.01\ \mathrm{m}$ after $\mathbf{2}$ complete oscillations.
\end{corrige}

\begin{aretenir}[Key points of the section]
\begin{itemize}
\item Damping adds a resistance $-\lambda\dot{x}$, giving $\ddot{x} + 2k\dot{x} + \omega^{2}x = 0$ with $k = \lambda/(2m)$.
\item The auxiliary equation $r^{2} + 2kr + \omega^{2} = 0$ decides everything via the sign of $k^{2} - \omega^{2}$.
\item Light damping ($k<\omega$): $x = e^{-kt}(A\cos\Omega t + B\sin\Omega t)$, $\Omega = \sqrt{\omega^{2}-k^{2}}$, damped period $T_d = 2\pi/\Omega > 2\pi/\omega$.
\item Critical damping ($k=\omega$): $x = (A+Bt)e^{-kt}$ — fastest return without oscillation; do not forget the $t$.
\item Heavy damping ($k>\omega$): $x = Ae^{-(k-\gamma)t} + Be^{-(k+\gamma)t}$, $\gamma = \sqrt{k^{2}-\omega^{2}}$ — slow, non-oscillatory.
\item Logarithmic decrement: $\Lambda = \ln(x_n/x_{n+1}) = kT_d$; the practical tool for measuring damping.
\end{itemize}
\end{aretenir}

\begin{savaistu}[Damping all around us]
Engineers choose the damping regime deliberately. Car shock absorbers are designed to be very nearly \emph{critically} damped so the wheels recover instantly from bumps. The needles of analogue ammeters and the moving parts of traditional balances are critically damped so the reading settles without swinging. Buildings in earthquake zones are fitted with giant tuned mass dampers — huge pendulums with dashpots — that absorb the oscillation energy of the structure. Conversely, a \emph{lightly} damped system is wanted when oscillations must persist, as in the quartz crystal of a watch. The single equation of this section governs them all.
\end{savaistu}

We now have the complete picture: undamped SHM, the energy and pendulum results, and the three damping regimes. In the next and final section of this chapter we put all of this to work on applications and full GCE examination problems, including composite systems and questions that mix SHM with damping.

\coursec{Applications and GCE Problem-Solving}

In the previous section we learnt that when a resistive force proportional to velocity acts on an oscillator, the equation of motion becomes $\ddot{x} + 2k\dot{x} + \omega^{2}x = 0$, and that the nature of the motion — underdamped, critically damped or overdamped — depends on the sign of $k^{2} - \omega^{2}$. We now put the whole machinery of this chapter to work. Real systems — a car suspension, a market spring balance, a log bobbing on a river, a pendulum in the school laboratory — are all modelled by the equations we have derived. The GCE examiner's favourite skill is precisely this: \emph{translate a physical situation into a differential equation, solve it, and interpret the answer}. This section trains exactly that skill.

\courssub{The Modelling Checklist}

Almost every structured question in this chapter follows the same logical path. Learn it as a routine.

\begin{methode}[Modelling checklist for oscillation problems]
\begin{enumerate}
\item \textbf{Find the equilibrium position first.} Resolve forces with the system at rest. This gives a relation (e.g.\ $mg = \lambda e/l$ for a spring) that will simplify the equation later.
\item \textbf{Define the displacement variable $x$ from equilibrium}, with a clearly stated positive direction.
\item \textbf{List all forces} in the displaced position: weight, tension or spring force (Hooke's law $T = \lambda(\text{extension})/l$), upthrust, resistance $-2mk\dot{x}$ if present.
\item \textbf{Apply $F = ma$} along the positive direction: $m\ddot{x} = \sum F$.
\item \textbf{Use the equilibrium relation} to cancel the constant terms. What remains should be $\ddot{x} = -\omega^{2}x$ (SHM) or $\ddot{x} + 2k\dot{x} + \omega^{2}x = 0$ (damped).
\item \textbf{Identify $\omega^{2}$ (and $k$)} by comparing coefficients.
\item \textbf{Solve} using the standard forms, then \textbf{apply the initial conditions}.
\item \textbf{Interpret}: period, amplitude, times, speeds — and state the limits of the model.
\end{enumerate}
\end{methode}

\begin{attention}[The sign of the damping term]
A resistance proportional to velocity \emph{opposes} the motion, so it must enter the equation with a \cle{negative sign}: $-2mk\dot{x}$. If you obtain $\ddot{x} - 2k\dot{x} + \omega^{2}x = 0$ (positive damping), the auxiliary equation gives roots with positive real part and the amplitude \emph{grows} exponentially — physically absurd for a passive system. A positive sign is a reliable signal that you have made a sign error in your force diagram.
\end{attention}

\courssub{Vehicle Suspension: Why Near-Critical Damping?}

A car wheel, together with the share of the chassis it carries, can be modelled as a mass $m$ supported by a spring of stiffness $s$ (force $= s \times$ compression) and a \emph{shock absorber} (damper) exerting a resistance $r\dot{x}$.

\begin{popfigure}
\begin{center}
\begin{tikzpicture}[scale=0.9, every node/.style={font=\small}]
% road
\draw[thick, popdark] (-2.2,0) -- (2.2,0);
\foreach \x in {-2.0,-1.6,...,2.0} \draw[popmuted] (\x,0) -- (\x-0.25,-0.25);
\node[popmuted] at (2.9,-0.15) {road};
% wheel
\draw[fill=popmuted!40] (0,0.45) circle (0.45);
\draw[fill=white] (0,0.45) circle (0.18);
% spring
\draw[thick, popblue] (0.7,0.9) -- (0.7,1.15);
\draw[thick, popblue] (0.7,1.15) -- (0.35,1.35) -- (1.05,1.6) -- (0.35,1.85) -- (1.05,2.1) -- (0.35,2.35) -- (1.05,2.6) -- (0.7,2.8);
\draw[thick, popblue] (0.7,2.8) -- (0.7,3.05);
\node[popblue, right] at (1.15,2.0) {spring $s$};
% damper
\draw[thick, poporange] (-0.7,0.9) -- (-0.7,1.7);
\draw[thick, poporange] (-1.0,1.7) rectangle (-0.4,2.5);
\draw[thick, poporange] (-0.7,2.5) -- (-0.7,2.35);
\draw[thick, poporange] (-0.95,2.35) -- (-0.45,2.35);
\draw[thick, poporange] (-0.7,2.35) -- (-0.7,3.05);
\node[poporange, left] at (-1.1,2.1) {damper $r\dot{x}$};
% chassis mass
\draw[fill=popblueL, thick, popdark] (-1.6,3.05) rectangle (1.6,3.75);
\node at (0,3.4) {chassis, mass $m$};
% displacement arrow
\draw[<->, poppink, thick] (2.1,3.05) -- (2.1,3.75) node[midway, right] {$x$};
\end{tikzpicture}
\end{center}
\legende{Quarter-car model: the chassis mass $m$ is connected to the wheel (and road) by a spring of stiffness $s$ and a damper resisting with force $r\dot{x}$.}
\end{popfigure}

Measuring $x$ from the equilibrium compression of the spring, the weight cancels the static spring force exactly as in the vertical mass--spring of Section 1, and the equation of motion is
\[
m\ddot{x} = -sx - r\dot{x}
\quad\Longrightarrow\quad
\ddot{x} + 2k\dot{x} + \omega^{2}x = 0,
\qquad
\omega^{2} = \frac{s}{m},\quad 2k = \frac{r}{m}.
\]
When the wheel hits a pothole, the chassis is displaced and the suspension must bring it back to rest. The three regimes feel very different to the passenger:

\begin{itemize}
\item \textbf{Underdamped} ($k < \omega$): the car bounces several times after the bump — uncomfortable and dangerous, since the tyres lose grip during the oscillation.
\item \textbf{Overdamped} ($k > \omega$): the return is slow and the suspension feels ``dead''; the car cannot respond to the next bump in time.
\item \textbf{Critically damped} ($k = \omega$): the chassis returns to equilibrium in the \emph{shortest possible time without oscillating}.
\end{itemize}

This is why engineers tune shock absorbers to be \cle{near-critical}: $k \approx \omega$, i.e.\ $r \approx 2\sqrt{ms}$. Critical damping is not a mathematical curiosity — it is a design target.

\begin{exemplebox}[Tuning a shock absorber]
A car's suspension carries a mass $m = 300\ \mathrm{kg}$ per wheel on a spring of stiffness $s = 27\,000\ \mathrm{N\,m^{-1}}$. Then
\[
\omega = \sqrt{\frac{s}{m}} = \sqrt{90} \approx 9.49\ \mathrm{rad\,s^{-1}}.
\]
Critical damping requires $k = \omega$, i.e.\ $r = 2m\omega \approx 2(300)(9.49) \approx 5\,690\ \mathrm{N\,s\,m^{-1}}$. A damper much weaker than this lets the car bounce; much stronger makes the ride harsh.
\end{exemplebox}

\courssub{The Spring Scale and the Market Balance}

A spring balance hanging in a Bamenda market stall is a vertical mass--spring system. When a bag of beans of unknown mass $m$ is hung on the hook and set bobbing gently, the period of oscillation is
\[
T = 2\pi\sqrt{\frac{m}{s}},
\]
where $s$ is the stiffness of the spring. Measuring $T$ (time ten oscillations with a phone stopwatch and divide by ten) therefore \emph{measures the mass}: $m = sT^{2}/(4\pi^{2})$. This is the principle of the \emph{inertial balance}, used where gravity is unreliable — for instance to ``weigh'' astronauts in orbit, where ordinary scales read zero but inertia is unchanged.

\begin{exemplebox}[Mass from the period]
A spring scale oscillates with period $T = 0.80\ \mathrm{s}$ when loaded with an unknown mass; its stiffness is $s = 50\ \mathrm{N\,m^{-1}}$. Then
\[
m = \frac{sT^{2}}{4\pi^{2}} = \frac{50 \times 0.64}{4\pi^{2}} \approx 0.81\ \mathrm{kg}.
\]
\end{exemplebox}

\courssub{Buoyancy Oscillations: a Floating Cylinder}

A cylindrical log of mass $m$ and uniform cross-sectional area $A$ floats vertically in water of density $\rho$. Push it down a little and release: it bobs up and down. Is this SHM? Let us prove it.

At equilibrium, the weight is balanced by the upthrust (Archimedes' principle): if $d$ is the submerged depth at equilibrium, $mg = \rho g A d$. Now displace the cylinder a further distance $x$ \emph{downwards} from equilibrium. The submerged depth becomes $d + x$, so the upthrust is $\rho g A (d+x)$, and the resultant force (taking downwards as positive) is
\[
m\ddot{x} = mg - \rho g A (d + x) = \underbrace{mg - \rho g A d}_{=\,0\ \text{(equilibrium)}} - \rho g A x = -\rho g A x.
\]

\begin{propriete}[Vertical oscillations of a floating body]
A body of mass $m$ and uniform cross-sectional area $A$, floating vertically in a liquid of density $\rho$ and displaced by $x$ from its floating equilibrium, satisfies
\[
\ddot{x} = -\frac{\rho g A}{m}\,x,
\]
and hence performs \cle{simple harmonic motion} with
\[
\omega^{2} = \frac{\rho g A}{m}, \qquad T = 2\pi\sqrt{\frac{m}{\rho g A}}.
\]
(The proof above — equilibrium condition plus Archimedes' principle — is examinable.)
\end{propriete}

\begin{popfigure}
\begin{center}
\begin{tikzpicture}[scale=1.0, every node/.style={font=\small}]
% water
\fill[popblueL] (-3,0) rectangle (3,-2.4);
\draw[thick, popblue] (-3,0) -- (3,0);
\node[popblue] at (2.5,0.25) {water};
% cylinder
\draw[fill=popgold!40, thick, popdark] (-0.6,-1.4) rectangle (0.6,0.9);
\node at (0,0.35) {mass $m$};
\node[popdark] at (0,-1.75) {area $A$};
% displacement
\draw[<->, poppink, thick] (1.0,0) -- (1.0,-0.6) node[midway, right] {$x$};
\draw[dashed, popmuted] (-0.6,-0.6) -- (1.0,-0.6);
\node[popmuted, left] at (-0.7,-0.6) {equil.\ level};
% upthrust (upwards)
\draw[->, very thick, popgreen] (0,-1.4) -- (0,-0.4);
\node[popgreen, right] at (0.05,-1.1) {upthrust $\rho g A(d+x)$};
% weight (downwards)
\draw[->, very thick, poporange] (0,0.9) -- (0,-0.1);
\node[poporange, left] at (-0.05,0.55) {weight $mg$};
\end{tikzpicture}
\end{center}
\legende{A floating cylinder pushed down by $x$: the upthrust exceeds the equilibrium value by $\rho g A x$, giving a restoring force proportional to $x$.}
\end{popfigure}

\begin{attention}[Units, units, units]
Always check units before substituting: $\omega$ in $\mathrm{rad\,s^{-1}}$, modulus of elasticity $\lambda$ in $\mathrm{N}$, natural length and extension in $\mathrm{m}$, density in $\mathrm{kg\,m^{-3}}$, area in $\mathrm{m^{2}}$. Mixing centimetres and metres is the single most frequent source of wrong periods at the GCE. Convert everything to SI units \emph{before} writing the equation of motion.
\end{attention}

\courssub{The Piston and Crankshaft (Extension)}

A car engine's piston is driven by a crank of radius $r$ rotating at angular speed $\omega$, connected to the piston by a rod of length $L$. Taking the crank angle as $\theta = \omega t$, geometry gives the piston displacement from the top of its stroke as
\[
x = r(1-\cos\omega t) + L - \sqrt{L^{2} - r^{2}\sin^{2}\omega t}.
\]
This is \emph{not} exactly SHM. But in a real engine $r \ll L$, so $\sqrt{L^{2}-r^{2}\sin^{2}\omega t} \approx L - \frac{r^{2}}{2L}\sin^{2}\omega t \approx L$, and
\[
x \approx r(1 - \cos\omega t),
\]
which is SHM of amplitude $r$ about the mid-stroke position. Engineers use this approximation to estimate piston speeds and the forces on the connecting rod; the small correction term is responsible for the high-frequency ``buzz'' of a fast engine. \emph{This result is an extension: you will not be asked to derive it, but it shows SHM appearing as an approximation inside real machines.}

\courssub{Measuring $g$ with a Simple Pendulum}

For small oscillations, a simple pendulum of length $l$ has period
\[
T = 2\pi\sqrt{\frac{l}{g}}
\quad\Longrightarrow\quad
T^{2} = \frac{4\pi^{2}}{g}\,l.
\]
This is a \emph{linear} relation between $T^{2}$ and $l$, with gradient $4\pi^{2}/g$. This \cle{linearisation} is the standard experimental strategy: instead of computing $g$ from one measurement (one bad reading ruins everything), plot $T^{2}$ against $l$ for several lengths and take the gradient of the best-fit line.

\begin{methode}[Experimental determination of $g$]
\begin{enumerate}
\item Suspend a small dense bob on a light inextensible string; measure $l$ from the support to the \emph{centre} of the bob.
\item Displace the bob through a small angle (less than about $10^{\circ}$) and release.
\item Time $20$ complete oscillations, repeat, and average; divide by $20$ to get $T$.
\item Repeat for at least five different lengths.
\item Plot $T^{2}$ (vertical) against $l$ (horizontal); draw the best-fit straight line through the origin.
\item Measure the gradient $G$; then $g = \dfrac{4\pi^{2}}{G}$.
\end{enumerate}
\end{methode}

\begin{popfigure}
\begin{center}
\begin{tikzpicture}
\begin{axis}[
    width=11cm, height=7.5cm,
    xlabel={$l$ (m)}, ylabel={$T^{2}$ (s$^{2}$)},
    xmin=0, xmax=1.1, ymin=0, ymax=4.6,
    xtick={0,0.2,0.4,0.6,0.8,1.0},
    ytick={0,1,2,3,4},
    grid=both, grid style={popline!40},
    axis line style={popdark},
    label style={popdark},
]
\addplot[only marks, mark=*, mark size=2pt, poporange] coordinates {
    (0.20,0.82) (0.40,1.63) (0.60,2.41) (0.80,3.28) (1.00,4.00)
};
\addplot[domain=0:1.1, thick, popblue] {4.03*x};
\addplot[dashed, popmuted] coordinates {(0.6,2.418) (1.0,2.418)};
\addplot[dashed, popmuted] coordinates {(1.0,2.418) (1.0,4.03)};
\node[popblue, anchor=south west, font=\small] at (axis cs:0.12,3.6) {gradient $G = \dfrac{4\pi^{2}}{g} \approx 4.03\ \mathrm{s^{2}\,m^{-1}}$};
\node[popdark, font=\small] at (axis cs:0.55,1.1) {$g = \dfrac{4\pi^{2}}{G} \approx 9.79\ \mathrm{m\,s^{-2}}$};
\end{axis}
\end{tikzpicture}
\end{center}
\legende{Plot of $T^{2}$ against $l$ for pendulum data. The gradient of the best-fit line is $4\pi^{2}/g$, so $g = 4\pi^{2}/G$.}
\end{popfigure}

\begin{exoresolu}[Computing $g$ from the graph]
A student measures the gradient of her $T^{2}$--$l$ graph as $G = 4.06\ \mathrm{s^{2}\,m^{-1}}$. Find the value of $g$ she should quote.
\tcblower
Since $T^{2} = \dfrac{4\pi^{2}}{g}l$, the gradient is $G = \dfrac{4\pi^{2}}{g}$, hence
\[
g = \frac{4\pi^{2}}{G} = \frac{4\pi^{2}}{4.06} = \frac{39.478}{4.06} \approx 9.72\ \mathrm{m\,s^{-2}}.
\]
She should quote $g \approx 9.7\ \mathrm{m\,s^{-2}}$, and could comment that the small discrepancy from the accepted value may come from reaction time in the timing or from measuring $l$ to the top rather than the centre of the bob.
\end{exoresolu}

\courssub{SHM or Damped? Reading the Problem Statement}

Before solving, decide \emph{which model} the question intends. The wording tells you.

\begin{center}
\begin{tabularx}{\linewidth}{|l|X|X|}
\hline
\rowcolor{popblueL}
\textbf{Clue in the statement} & \textbf{Model} & \textbf{Equation} \\
\hline
``light resistance neglected'', ``smooth'', no mention of friction & SHM & $\ddot{x} = -\omega^{2}x$ \\
\hline
``resistance proportional to speed'', ``dashpot'', ``shock absorber'', ``damping force $2mkv$'' & Damped & $\ddot{x} + 2k\dot{x} + \omega^{2}x = 0$ \\
\hline
``comes to rest without oscillating as quickly as possible'' & Critical damping & $k = \omega$ \\
\hline
``oscillations die away'', ``amplitude decreases'' & Underdamped & $k < \omega$ \\
\hline
\end{tabularx}
\end{center}

If a resistance is mentioned, it \emph{must} appear in your force equation; if none is mentioned, do not invent one. Conversely, if a question later says ``in practice the amplitude decreases, criticise the model'', the expected answer is that real resistance (air drag, internal friction) was neglected.

\courssub{A Full GCE-Type Structured Problem}

\begin{exoresolu}[Floating buoy (structured question)]
A cylindrical buoy of mass $m = 12\ \mathrm{kg}$ and cross-sectional area $A = 0.030\ \mathrm{m^{2}}$ floats vertically in sea water of density $\rho = 1025\ \mathrm{kg\,m^{-3}}$. Take $g = 9.81\ \mathrm{m\,s^{-2}}$.

\textbf{(a)} Show that, when the buoy is displaced vertically by $x$ metres from equilibrium and resistance is neglected, $\ddot{x} = -25.1\,x$ approximately, and state the value of $\omega$.

\textbf{(b)} The buoy is pushed down $0.20\ \mathrm{m}$ and released from rest. Find the period and the maximum speed.

\textbf{(c)} Find the time at which the buoy first passes through its equilibrium position.

\textbf{(d)} In practice the oscillations die away. Criticise the model.
\tcblower
\textbf{(a)} At equilibrium, $mg = \rho g A d$ where $d$ is the submerged depth. Displaced a further $x$ downwards, the resultant downward force is
\[
m\ddot{x} = mg - \rho g A(d+x) = -\rho g A x.
\]
Hence
\[
\ddot{x} = -\frac{\rho g A}{m}\,x = -\frac{1025 \times 9.81 \times 0.030}{12}\,x = -25.14\,x \approx -25.1\,x.
\]
So $\omega^{2} = 25.14$, giving $\omega \approx 5.01\ \mathrm{rad\,s^{-1}}$.

\textbf{(b)} Period: $T = \dfrac{2\pi}{\omega} = \dfrac{2\pi}{5.01} \approx 1.25\ \mathrm{s}$. The motion starts from rest at maximum displacement, so $x = 0.20\cos(5.01\,t)$ and the amplitude is $a = 0.20\ \mathrm{m}$. Maximum speed:
\[
v_{\max} = a\omega = 0.20 \times 5.01 \approx 1.00\ \mathrm{m\,s^{-1}}.
\]

\textbf{(c)} The buoy first passes through $x = 0$ when $\cos(5.01\,t) = 0$, i.e.\ $5.01\,t = \pi/2$, giving
\[
t = \frac{\pi}{2\omega} = \frac{\pi}{10.03} \approx 0.31\ \mathrm{s} \quad \Bigl(= \tfrac{T}{4}\Bigr).
\]

\textbf{(d)} The model neglects the resistance of the water (viscous drag, roughly proportional to velocity). Including a term $-2mk\dot{x}$ gives damped harmonic motion, whose amplitude decays like $e^{-kt}$ — matching the observed dying away. The SHM model is therefore only valid for the first few oscillations.
\end{exoresolu}

\begin{attention}[GCE exam tip: use the stated result even if you cannot prove it]
Structured questions very often carry the instruction ``\emph{show that} $\ddot{x} = -\omega^{2}x$'' with a given numerical $\omega$. If you cannot complete the proof, \textbf{do not abandon the question}: the later parts (period, amplitude, times, speeds) only need the \emph{stated} value of $\omega$. Write ``from part (a), $\omega = \ldots$'' and continue. You keep nearly all the marks.
\end{attention}

\courssub{Checking and Criticising a Model}

A complete GCE answer often ends with a discussion mark. Here is what examiners expect.

\textbf{Effect of neglected damping.} If the statement says resistance is negligible but the real system loses energy, the SHM model overestimates the amplitude and the speeds after the first few cycles; the true period of a damped oscillator, $2\pi/\sqrt{\omega^{2}-k^{2}}$, is slightly \emph{longer} than $2\pi/\omega$.

\textbf{Validity of the small-angle approximation.} The pendulum equation is really $\ddot{\theta} = -\frac{g}{l}\sin\theta$; replacing $\sin\theta$ by $\theta$ is only accurate for small amplitudes (a few degrees). For swings of $30^{\circ}$ or more, the true period is measurably longer than $2\pi\sqrt{l/g}$, and a $g$ measurement made with large swings will be slightly too small.

\textbf{Units and orders of magnitude.} Before quoting a final answer, sanity-check it: a pendulum of length $1\ \mathrm{m}$ must have a period near $2\ \mathrm{s}$; a car suspension should oscillate at roughly one cycle per second; a floating log bobs with a period of order one second. An answer of $T = 0.003\ \mathrm{s}$ for a pendulum signals a units slip — almost always centimetres substituted as metres.

\begin{exercice}[GCE practice]
A mass of $0.5\ \mathrm{kg}$ hangs from a light spring of natural length $0.6\ \mathrm{m}$ and modulus $30\ \mathrm{N}$, extending it to equilibrium. The mass is pulled down a further $0.05\ \mathrm{m}$ and released. A resistance of magnitude $2v$ newtons acts, where $v$ is the speed.
\begin{enumerate}
\item Show that $\ddot{x} + 4\dot{x} + 100x = 0$.
\item State the type of damping and write down the general solution.
\item Find the period of the damped oscillations and the time for the amplitude to fall to half its initial value.
\end{enumerate}
\end{exercice}

\begin{corrige}[Exercise 1]
\textbf{1.} Equilibrium: $mg = \lambda e/l \Rightarrow 0.5g = 30e/0.6$, so $e = 0.0981\ \mathrm{m}$. With a further displacement $x$ downwards, tension $= \lambda(e+x)/l = 50(e+x)$, and
\[
0.5\ddot{x} = 0.5g - 50(e+x) - 2\dot{x} = -50x - 2\dot{x},
\]
using $0.5g = 50e$. Dividing by $0.5$: $\ddot{x} + 4\dot{x} + 100x = 0$.

\textbf{2.} Here $k = 2$, $\omega = 10$, and $k^{2} - \omega^{2} = 4 - 100 < 0$: \textbf{underdamped}. General solution:
\[
x = e^{-2t}\bigl(A\cos(\sqrt{96}\,t) + B\sin(\sqrt{96}\,t)\bigr), \qquad \sqrt{96} = 4\sqrt{6} \approx 9.80.
\]

\textbf{3.} Damped period: $T = \dfrac{2\pi}{\sqrt{96}} \approx 0.641\ \mathrm{s}$. Amplitude decays as $e^{-2t}$; halving requires $e^{-2t} = \tfrac12$, so $t = \dfrac{\ln 2}{2} \approx 0.347\ \mathrm{s}$.
\end{corrige}

\begin{aretenir}[Key points of Section 4]
\begin{itemize}
\item Follow the modelling checklist: equilibrium first, then $F = ma$ from equilibrium, then identify $\omega^{2}$ (and $k$).
\item A floating body of cross-section $A$ in a liquid of density $\rho$ performs SHM with $\omega^{2} = \rho g A/m$.
\item Engineers design suspensions for \emph{near-critical} damping: fastest return to equilibrium without bouncing.
\item Linearise $T^{2} = 4\pi^{2}l/g$ and take the gradient to measure $g$ experimentally.
\item Read the statement to choose between SHM and damped models; damping always enters with a negative sign.
\item If you cannot prove the displayed equation, still use the stated $\omega$ in later parts.
\item Always check units (SI throughout) and the order of magnitude of your final answer.
\end{itemize}
\end{aretenir}

\begin{savaistu}[Did you know?]
The same mathematics that keeps a car stable on the potholed roads between Bamenda and Douala also protects buildings during earthquakes: engineers fit giant tuned mass--damper systems — effectively huge mass--spring--damper oscillators — inside skyscrapers so that the building's vibrations are absorbed and die out almost critically damped. The equation $\ddot{x} + 2k\dot{x} + \omega^{2}x = 0$ quite literally saves lives.
\end{savaistu}

\newpage

\coursec{Integration activity}

\courssub{Aims of the activity}

This integration activity closes the chapter by asking you to act as a \cle{design engineer}. Working in groups of three or four, you will mobilise \emph{all} the skills developed in this chapter:

\begin{itemize}
\item model a real mechanical system by a differential equation of the form $m\ddot{x}+\lambda\dot{x}+kx=0$;
\item compute the \cle{natural angular frequency} and the undamped period of an oscillator;
\item identify the three damping regimes (underdamped, critically damped, overdamped) from the discriminant of the auxiliary equation;
\item solve a damped equation of motion with given initial conditions and sketch the response;
\item compute the \cle{damped period} and the \cle{logarithmic decrement}, and interpret them physically;
\item justify an engineering recommendation with rigorous mathematics, and defend it orally.
\end{itemize}

Beyond the calculations, the activity trains the GCE examiner's favourite skill: \emph{interpreting a mathematical result in the language of the real world}.

\courssub{Situation}

The road from Bamenda to Bafoussam, winding through the Western highlands, is famous for its beautiful scenery --- and for its potholes. A transport company, \emph{Highland Express}, runs pick-up trucks on this route and has received many complaints: passengers feel shaken long after each pothole, and goods arrive damaged. The company hires your group to advise on the \cle{suspension} of its vehicles.

Each wheel of a pick-up truck carries an effective mass $m = 300\ \mathrm{kg}$ (one quarter of the $1200\ \mathrm{kg}$ vehicle, neglecting the mass of the wheel itself). The suspension of each wheel is modelled as:

\begin{itemize}
\item a \textbf{spring} of stiffness $k = 20\,000\ \mathrm{N\,m^{-1}}$, and
\item a \textbf{damper} (shock absorber) of damping coefficient $\lambda$ (in $\mathrm{N\,s\,m^{-1}}$),
\end{itemize}

mounted \emph{in parallel} between the wheel axle and the chassis.

\begin{popfigure}
\begin{center}
\begin{tikzpicture}[scale=1.0]
% chassis
\fill[popblueL] (-1.6,2.6) rectangle (1.6,3.1);
\node at (0,3.45) {\small chassis (mass $m$)};
% spring
\draw[thick,popblue] (-0.8,2.6) -- (-0.8,2.3) -- (-1.15,2.1) -- (-0.45,1.8) -- (-1.15,1.5) -- (-0.45,1.2) -- (-1.15,0.9) -- (-0.45,0.6) -- (-0.8,0.4) -- (-0.8,0.1);
\node[popblue] at (-2.05,1.35) {\small spring $k$};
% damper
\draw[thick,poporange] (0.8,2.6) -- (0.8,1.9);
\draw[thick,poporange] (0.45,1.9) -- (1.15,1.9);
\draw[thick,poporange] (0.45,1.9) -- (0.45,1.15);
\draw[thick,poporange] (1.15,1.9) -- (1.15,1.15);
\draw[thick,poporange] (0.3,1.15) -- (1.3,1.15);
\draw[thick,poporange] (0.8,1.15) -- (0.8,0.1);
\node[poporange] at (2.15,1.35) {\small damper $\lambda$};
% axle / wheel
\fill[popgreenL] (-1.6,-0.35) rectangle (1.6,0.1);
\node at (0,-0.62) {\small wheel axle};
\draw[thick] (0,-0.35) circle (0.62);
% ground with pothole
\draw[thick] (-3,-1.6) -- (-0.9,-1.6) .. controls (-0.5,-2.15) and (0.5,-2.15) .. (0.9,-1.6) -- (3,-1.6);
\node at (2.1,-2.0) {\small pothole};
% x axis
\draw[->,thick] (2.6,2.9) -- (2.6,1.7) node[midway,right] {\small $x$};
\end{tikzpicture}
\end{center}
\legende{Quarter-car model of the suspension: spring and damper in parallel.}
\end{popfigure}

When the wheel hits a pothole, the chassis is suddenly displaced from its equilibrium position. We take the \cle{equilibrium position} (with the truck at rest, the spring already compressed by the weight) as the origin, and we denote by $x(t)$ the vertical displacement of the chassis from equilibrium, measured \emph{downwards}, in metres, at time $t$ (in seconds) after the pothole impact.

Because the origin is the static equilibrium, the weight $mg$ is exactly balanced by the static compression of the spring, so gravity \emph{does not appear} in the equation of motion: only the \emph{extra} spring force $-kx$ and the damping force $-\lambda\dot{x}$ act.

The company tells you that a typical pothole impact can be modelled as an \textbf{initial displacement} $x(0)=0.05\ \mathrm{m}$ ($5\ \mathrm{cm}$) with \textbf{zero initial velocity}, $\dot{x}(0)=0$.

Your group must recommend a value of $\lambda$ to the company, with full mathematical justification.

\courssub{Tasks}

\textbf{Part A --- Modelling and the undamped oscillator}

\begin{enumerate}
\item Draw a force diagram for the chassis mass at a displacement $x$ below equilibrium, and use Newton's second law to show that
\[ m\ddot{x}+\lambda\dot{x}+kx=0, \qquad\text{i.e.}\qquad 300\,\ddot{x}+\lambda\,\dot{x}+20\,000\,x=0. \]
\item Suppose first that the damper is removed ($\lambda=0$). Show that the motion is simple harmonic, and compute the natural angular frequency $\omega$ and the undamped period $T_0$. Give $T_0$ to 3 decimal places. Comment: is this a comfortable frequency for passengers?
\end{enumerate}

\textbf{Part B --- Critical damping}

\begin{enumerate}
\setcounter{enumi}{2}
\item Write the auxiliary (characteristic) equation of the general equation of motion, and state the condition on $\lambda$ for the damping to be \emph{critical}.
\item Show that the critical damping coefficient is $\lambda_c = 2\sqrt{mk}$, and compute its numerical value.
\item Explain, using the form of the solutions in each regime, why car engineers choose damping \emph{close to critical}: what happens to passengers if $\lambda$ is too small? And if $\lambda$ is too large?
\end{enumerate}

\textbf{Part C --- The actual damper: $\lambda = 3000\ \mathrm{N\,s\,m^{-1}}$}

\begin{enumerate}
\setcounter{enumi}{5}
\item Compute the discriminant of the auxiliary equation and identify the damping regime.
\item Solve the equation of motion with $x(0)=0.05$, $\dot{x}(0)=0$. Write $x(t)$ in the form $x(t)=e^{-\alpha t}\bigl(A\cos\omega_d t+B\sin\omega_d t\bigr)$, giving $\alpha$, $\omega_d$, $A$, $B$ to 3 significant figures.
\item Sketch the graph of $x(t)$ for $0\le t\le 1.5\ \mathrm{s}$, marking the exponential envelope $\pm 0.05\,e^{-\alpha t}$.
\item Compute the damped period $T_d$ and the logarithmic decrement $\delta=\ln\dfrac{x(t)}{x(t+T_d)}=\alpha T_d$.
\item Hence estimate how many \emph{visible} oscillations the passengers feel after one pothole (an oscillation is ``visible'' while the amplitude exceeds about $1\ \mathrm{mm}$). Is this acceptable comfort?
\end{enumerate}

\textbf{Part D --- Recommendation}

\begin{enumerate}
\setcounter{enumi}{10}
\item Write a short report (half a page) to \emph{Highland Express} recommending a value (or range) of $\lambda$, justifying your choice mathematically (regime, settling time, number of oscillations) and in everyday language. Each group presents its recommendation to the class in three minutes.
\end{enumerate}

\courssub{Model answer}

\textbf{Part A.}\quad \textbf{1.} At displacement $x$ below equilibrium, the spring is stretched by $x$ beyond its static compression, so it pulls the chassis \emph{up} with an extra force $kx$; the damper opposes the velocity with force $\lambda\dot{x}$. Newton's second law (positive downwards) gives
\[ m\ddot{x}=-kx-\lambda\dot{x}\quad\Longrightarrow\quad 300\,\ddot{x}+\lambda\,\dot{x}+20\,000\,x=0. \]

\textbf{2.} With $\lambda=0$: $\ddot{x}+\dfrac{20\,000}{300}x=0$, i.e. $\ddot{x}+\omega^2 x=0$ with
\[ \omega=\sqrt{\frac{k}{m}}=\sqrt{\frac{20\,000}{300}}=\sqrt{66.67}\approx 8.16\ \mathrm{rad\,s^{-1}},\qquad
T_0=\frac{2\pi}{\omega}\approx 0.770\ \mathrm{s}. \]
This is SHM: without a damper the truck would bounce \emph{for ever} about once every $0.77\ \mathrm{s}$ --- extremely uncomfortable and dangerous.

\textbf{Part B.}\quad \textbf{3.} The auxiliary equation is $300r^2+\lambda r+20\,000=0$. Critical damping corresponds to a repeated root: $\Delta=\lambda^2-4(300)(20\,000)=0$.

\textbf{4.} $\lambda_c^2=4mk\cdot$... more precisely $\lambda_c=2\sqrt{mk}=2\sqrt{300\times 20\,000}=2\sqrt{6\,000\,000}\approx 4899\ \mathrm{N\,s\,m^{-1}}$, i.e. $\lambda_c\approx 4900\ \mathrm{N\,s\,m^{-1}}$.

\textbf{5.} If $\lambda<\lambda_c$ (underdamped), the solution contains a factor $\cos(\omega_d t-\varphi)$: the chassis \emph{oscillates} through equilibrium several times before settling, so passengers feel repeated bouncing. If $\lambda>\lambda_c$ (overdamped), the return to equilibrium is a sum of two decaying exponentials, the slowest of which decays like $e^{r_2 t}$ with $r_2$ close to $0$: the return is \emph{sluggish} and the suspension feels ``hard'' over successive bumps. Critical damping gives the \emph{fastest} return to equilibrium \emph{without} oscillation, since $x=(A+Bt)e^{-\omega t}$ decays faster than any overdamped slow mode. Hence engineers aim for $\lambda$ slightly below or equal to $\lambda_c$.

\textbf{Part C.}\quad \textbf{6.} With $\lambda=3000$: $\Delta=3000^2-4(300)(20\,000)=9\times10^6-2.4\times10^7=-1.5\times10^7<0$: the motion is \cle{underdamped}. The roots are
\[ r=\frac{-3000\pm\sqrt{-1.5\times10^7}}{600}=-5\pm i\sqrt{\tfrac{1.5\times10^7}{600^2}}=-5\pm i\,6.45. \]
So $\alpha=\dfrac{\lambda}{2m}=5\ \mathrm{s^{-1}}$ and $\omega_d=\sqrt{\omega^2-\alpha^2}=\sqrt{66.67-25}=\sqrt{41.67}\approx 6.45\ \mathrm{rad\,s^{-1}}$.

\textbf{7.} General solution: $x(t)=e^{-5t}\bigl(A\cos 6.455t+B\sin 6.455t\bigr)$.
$x(0)=A=0.05$. Differentiating,
$\dot{x}=e^{-5t}\bigl[(-5A+6.455B)\cos 6.455t+(-5B-6.455A)\sin 6.455t\bigr]$, so
$\dot{x}(0)=-5A+6.455B=0\Rightarrow B=\dfrac{5(0.05)}{6.455}\approx 0.0387$.
\[ \boxed{\;x(t)=e^{-5t}\bigl(0.0500\cos 6.45t+0.0387\sin 6.45t\bigr)\ \text{(m)}\;} \]

\textbf{8.} Sketch:

\begin{popfigure}
\begin{center}
\begin{tikzpicture}
\begin{axis}[width=12cm,height=6.5cm,xlabel={$t$ (s)},ylabel={$x$ (m)},
xmin=0,xmax=1.5,ymin=-0.06,ymax=0.06,grid=both,grid style={popline!40},
samples=200,domain=0:1.5,axis lines=middle]
\addplot[popblue,thick] {exp(-5*x)*(0.05*cos(deg(6.455*x))+0.0387*sin(deg(6.455*x)))};
\addplot[poporange,dashed] {0.05*exp(-5*x)};
\addplot[poporange,dashed] {-0.05*exp(-5*x)};
\end{axis}
\end{tikzpicture}
\end{center}
\legende{Underdamped response $x(t)$ with exponential envelope $\pm 0.05e^{-5t}$.}
\end{popfigure}

\textbf{9.} Damped period: $T_d=\dfrac{2\pi}{\omega_d}=\dfrac{2\pi}{6.455}\approx 0.974\ \mathrm{s}$.
Logarithmic decrement:
\[ \delta=\alpha T_d=5\times 0.974\approx 4.87. \]
The ratio of successive amplitudes on the same side is $e^{-\delta}\approx e^{-4.87}\approx 0.0077$: after \emph{one} full damped period the displacement is less than $1\%$ of its initial value.

\textbf{10.} Amplitude after $n$ periods: $0.05\,e^{-4.87n}$. For $n=1$: $0.05\times0.0077\approx 3.9\times10^{-4}\ \mathrm{m}=0.39\ \mathrm{mm}<1\ \mathrm{mm}$. So the passengers feel essentially \textbf{one single dip-and-recovery}, and no repeated bouncing. The motion settles within about one second. This is comfortable and acceptable.

\textbf{Part D.}\quad \textbf{11.} \emph{Model report.} We recommend a damper coefficient in the range $4000\le\lambda\le 4900\ \mathrm{N\,s\,m^{-1}}$. The present value $\lambda=3000\ \mathrm{N\,s\,m^{-1}}$ is only $61\%$ of the critical value $\lambda_c\approx 4900$: although the oscillation dies out within one period, the first overshoot below equilibrium (about $-0.036\ \mathrm{m}$, read from the model) is large and the ride feels ``soft and wallowing''. Choosing $\lambda$ close to $\lambda_c$ eliminates the overshoot entirely, gives the fastest settling time $x=(A+Bt)e^{-8.16t}$, and protects both passengers and goods. Values above $\lambda_c$ should be avoided: the slow exponential mode would make the suspension react sluggishly to the rapid succession of potholes typical of the Bamenda--Bafoussam road. \emph{Recommendation: fit dampers of $\lambda\approx 4500\ \mathrm{N\,s\,m^{-1}}$ per wheel, about $92\%$ of critical damping.}

\begin{attention}[Marking points for the presentation]
Examiners (and your classmates) will check: correct use of the equilibrium origin (weight cancels the static spring force); correct auxiliary equation and discriminant; units carried throughout ($\mathrm{rad\,s^{-1}}$, $\mathrm{N\,s\,m^{-1}}$, s); and a conclusion that \emph{uses the numbers} rather than vague phrases such as ``the damping is good''.
\end{attention}

\newpage

\coursec{Methods and worked examples}

\courssub{Skill 1: Recognising and Setting Up the Equation of SHM}

\begin{methode}[Proving that a motion is simple harmonic]
To show that a particle moves with simple harmonic motion and to find its period:
\begin{enumerate}
\item Choose a \cle{coordinate} $x$ measured from the \textbf{equilibrium position} of the particle (this is essential: SHM is about oscillation \emph{about equilibrium}).
\item Draw a diagram showing all forces acting at a \textbf{general displaced position} (weight, tension, spring force, etc.).
\item Apply \textbf{Newton's second law} $F = m\ddot{x}$ along the direction of motion, taking care with signs (use Hooke's law $T = \lambda e / l$ for springs or elastic strings).
\item Simplify until the equation takes the standard form
\[ \ddot{x} = -\omega^{2} x . \]
\item Read off $\omega^{2}$, deduce $\omega$, then the \textbf{period} $T = \dfrac{2\pi}{\omega}$ and the \textbf{frequency} $f = \dfrac{1}{T}$.
\end{enumerate}
\textbf{Reflex:} if the equation comes out as $\ddot{x} = -\omega^{2}(x - a)$, the motion is still SHM, but about the shifted centre $x = a$; put $y = x - a$.
\end{methode}

\begin{exoresolu}[Particle on a vertical spring]
A particle of mass $0.5$ kg hangs in equilibrium from a light elastic spring of natural length $0.6$ m and modulus of elasticity $12$ N, the upper end being fixed. The particle is pulled down a further $0.1$ m and released. Taking $g = 10\ \mathrm{m\,s^{-2}}$, show that the particle moves with SHM and find the period of the motion.
\tcblower
\textbf{Step 1: equilibrium.} Let $e_{0}$ be the extension at equilibrium. Then
\[ \frac{\lambda e_{0}}{l} = mg \quad\Longrightarrow\quad \frac{12\,e_{0}}{0.6} = 0.5 \times 10 \quad\Longrightarrow\quad e_{0} = 0.25\ \mathrm{m}. \]
\textbf{Step 2: equation of motion.} Let $x$ be the displacement \emph{below} equilibrium. The total extension is $e_{0} + x$, so the tension is $T = \dfrac{12}{0.6}(0.25 + x) = 5 + 20x$. Taking the downward direction as positive, Newton's second law gives
\[ m\ddot{x} = mg - T = 5 - (5 + 20x) = -20x . \]
Hence
\[ 0.5\,\ddot{x} = -20x \quad\Longrightarrow\quad \ddot{x} = -40x . \]
This is of the form $\ddot{x} = -\omega^{2}x$ with $\omega^{2} = 40$, so the motion is \textbf{simple harmonic} about the equilibrium position.
\textbf{Step 3: period.}
\[ \omega = \sqrt{40} = 2\sqrt{10}, \qquad T = \frac{2\pi}{\omega} = \frac{2\pi}{2\sqrt{10}} = \frac{\pi}{\sqrt{10}} \approx 0.99\ \mathrm{s}. \]
\textbf{Remark:} the amplitude is $0.1$ m; since the spring remains extended throughout ($0.25 - 0.1 > 0$), the spring never goes slack and the SHM is complete.
\end{exoresolu}

\begin{popfigure}
\begin{center}
\begin{tikzpicture}[scale=0.9]
% fixed support
\draw[thick] (-1.2,0) -- (1.2,0);
\foreach \x in {-1.0,-0.6,...,1.0} \draw (\x,0) -- (\x+0.2,0.25);
% spring (zigzag)
\draw[thick,popblue] (0,0) -- (0,-0.3);
\draw[thick,popblue] (0,-0.3) -- (0.18,-0.5) -- (-0.18,-0.7) -- (0.18,-0.9) -- (-0.18,-1.1) -- (0.18,-1.3) -- (-0.18,-1.5) -- (0,-1.7);
\draw[thick,popblue] (0,-1.7) -- (0,-2.0);
% mass
\fill[poporange] (0,-2.35) circle (0.35);
\node at (0,-2.35) {\small $m$};
% equilibrium mark
\draw[dashed,popmuted] (-1.6,-2.35) -- (1.6,-2.35);
\node[left,popmuted] at (-1.6,-2.35) {\small equilibrium};
% displacement arrow
\draw[->,thick,popdark] (1.1,-2.35) -- (1.1,-3.15) node[midway,right] {\small $x$};
% forces
\draw[->,thick,popgreen] (-0.9,-2.35) -- (-0.9,-1.35) node[midway,left] {\small $T$};
\draw[->,thick,popgreen] (-0.55,-2.35) -- (-0.55,-3.35) node[midway,left] {\small $mg$};
\end{tikzpicture}
\end{center}
\legende{Vertical spring: displacement $x$ measured from the equilibrium position, with tension $T$ upwards and weight $mg$ downwards.}
\end{popfigure}

\courssub{Skill 2: Using the Standard Solutions of SHM}

\begin{methode}[Exploiting $x = a\cos(\omega t + \varepsilon)$ and $v^{2} = \omega^{2}(a^{2} - x^{2})$]
Given that a particle performs SHM with $\ddot{x} = -\omega^{2}x$:
\begin{enumerate}
\item Identify the \textbf{amplitude} $a$ (maximum displacement from the centre) and $\omega$.
\item Write $x = a\cos(\omega t + \varepsilon)$, choosing $\varepsilon$ from the \textbf{initial conditions} (if released from rest at $x = a$, simply $x = a\cos\omega t$).
\item Use the velocity formulas:
\[ v = \dot{x} = -\omega a \sin(\omega t + \varepsilon), \qquad v^{2} = \omega^{2}(a^{2} - x^{2}). \]
\item Maximum speed: $v_{\max} = \omega a$ (at the centre); maximum acceleration: $\omega^{2}a$ (at the extremes).
\item For a \textbf{time between two positions}, solve $x = a\cos(\omega t + \varepsilon)$ for $t$ at each position and subtract.
\end{enumerate}
\textbf{Reflex:} the formula $v^{2} = \omega^{2}(a^{2} - x^{2})$ links speed and position \emph{without time} — use it whenever time is not mentioned.
\end{methode}

\begin{exoresolu}[Speeds and times in SHM]
A particle moves with SHM of period $2$ s and amplitude $0.5$ m. Find (a) the maximum speed; (b) the speed when the particle is $0.3$ m from the centre; (c) the time taken to travel directly from $x = 0.3$ m to $x = -0.3$ m.
\tcblower
We have $T = 2$ s, so $\omega = \dfrac{2\pi}{T} = \pi\ \mathrm{rad\,s^{-1}}$, and $a = 0.5$ m.

\textbf{(a) Maximum speed:}
\[ v_{\max} = \omega a = \pi \times 0.5 = \frac{\pi}{2} \approx 1.57\ \mathrm{m\,s^{-1}}. \]

\textbf{(b) Speed at $x = 0.3$ m:} using $v^{2} = \omega^{2}(a^{2} - x^{2})$,
\[ v^{2} = \pi^{2}\left(0.5^{2} - 0.3^{2}\right) = \pi^{2}(0.25 - 0.09) = 0.16\pi^{2}, \]
\[ v = 0.4\pi \approx 1.26\ \mathrm{m\,s^{-1}}. \]

\textbf{(c) Time from $x = 0.3$ to $x = -0.3$:} take $x = 0.5\cos(\pi t)$. At $x = 0.3$:
\[ \cos(\pi t_{1}) = 0.6 \quad\Longrightarrow\quad \pi t_{1} = \arccos(0.6) \approx 0.9273, \quad t_{1} \approx 0.2952\ \mathrm{s}. \]
At $x = -0.3$ (first time after passing the centre):
\[ \cos(\pi t_{2}) = -0.6 \quad\Longrightarrow\quad \pi t_{2} = \pi - 0.9273 = 2.2143, \quad t_{2} \approx 0.7048\ \mathrm{s}. \]
Hence the required time is
\[ t_{2} - t_{1} \approx 0.7048 - 0.2952 = 0.410\ \mathrm{s}. \]
\textbf{Check:} the particle passes the centre ($x = 0$) at $t = 0.5$ s, and indeed $0.2952 < 0.5 < 0.7048$. $\blacksquare$
\end{exoresolu}

\courssub{Skill 3: Determining Amplitude and Phase from Initial Conditions}

\begin{methode}[Finding $a$ and $\varepsilon$ from $x(0)$ and $v(0)$]
\begin{enumerate}
\item Write the general solution $x = a\cos(\omega t + \varepsilon)$ (or equivalently $x = A\cos\omega t + B\sin\omega t$).
\item Substitute $t = 0$: $x(0) = a\cos\varepsilon$ and $v(0) = -\omega a \sin\varepsilon$.
\item Compute the amplitude from the energy-like identity
\[ a^{2} = x(0)^{2} + \frac{v(0)^{2}}{\omega^{2}} . \]
\item Determine $\varepsilon$ from the \textbf{two} equations of step 2 (signs of both sine and cosine fix the quadrant).
\end{enumerate}
\textbf{Reflex:} the form $x = A\cos\omega t + B\sin\omega t$ is often quicker: $A = x(0)$, $B = v(0)/\omega$.
\end{methode}

\begin{exoresolu}[Motion started from an interior point]
A particle performs SHM with $\ddot{x} = -9x$ (units: m, s). At $t = 0$ the particle is at $x = 2$ m moving with speed $6\ \mathrm{m\,s^{-1}}$ in the direction of $x$ increasing. Find the amplitude, the maximum speed, and the first time at which the particle is instantaneously at rest.
\tcblower
Here $\omega^{2} = 9$, so $\omega = 3\ \mathrm{rad\,s^{-1}}$.

\textbf{Amplitude:}
\[ a^{2} = x(0)^{2} + \frac{v(0)^{2}}{\omega^{2}} = 2^{2} + \frac{6^{2}}{3^{2}} = 4 + 4 = 8 \quad\Longrightarrow\quad a = 2\sqrt{2}\ \mathrm{m}. \]

\textbf{Maximum speed:}
\[ v_{\max} = \omega a = 3 \times 2\sqrt{2} = 6\sqrt{2} \approx 8.49\ \mathrm{m\,s^{-1}}. \]

\textbf{Phase and first rest position:} write $x = 2\sqrt{2}\cos(3t + \varepsilon)$. At $t = 0$:
\[ 2\sqrt{2}\cos\varepsilon = 2 \;\Rightarrow\; \cos\varepsilon = \tfrac{1}{\sqrt{2}}, \qquad -3 \cdot 2\sqrt{2}\sin\varepsilon = 6 \;\Rightarrow\; \sin\varepsilon = -\tfrac{1}{\sqrt{2}}. \]
Hence $\varepsilon = -\dfrac{\pi}{4}$, and $x = 2\sqrt{2}\cos\left(3t - \tfrac{\pi}{4}\right)$.

The particle is at rest at the extremes, i.e. when $3t - \tfrac{\pi}{4} = 0$ (first occasion, since it is moving towards $+x$):
\[ t = \frac{\pi}{12} \approx 0.262\ \mathrm{s}. \]
\textbf{Check:} at $t = \pi/12$, $x = 2\sqrt{2}\cos 0 = 2\sqrt{2} = a$. $\blacksquare$
\end{exoresolu}

\courssub{Skill 4: Energy Methods in SHM}

\begin{methode}[Using kinetic and potential energy in SHM]
\begin{enumerate}
\item Recall that in SHM the \textbf{total mechanical energy} is constant:
\[ E = \tfrac{1}{2}m\omega^{2}a^{2} = \tfrac{1}{2}m v_{\max}^{2}. \]
\item At displacement $x$: kinetic energy $E_{k} = \tfrac{1}{2}mv^{2} = \tfrac{1}{2}m\omega^{2}(a^{2} - x^{2})$ and potential energy $E_{p} = \tfrac{1}{2}m\omega^{2}x^{2}$.
\item To find where $E_{k} = E_{p}$ (or any given ratio), equate the expressions and solve for $x$.
\item For a spring, $E_{p} = \dfrac{\lambda x^{2}}{2l}$ measured from the \emph{natural} length; energy conservation often replaces the equation of motion.
\end{enumerate}
\textbf{Reflex:} "find the position where kinetic energy equals potential energy" $\Rightarrow$ $a^{2} - x^{2} = x^{2}$, so $x = \pm \dfrac{a}{\sqrt{2}}$.
\end{methode}

\begin{exoresolu}[Energy distribution in an oscillation]
A particle of mass $0.2$ kg performs SHM with amplitude $0.4$ m and period $\pi$ s. (a) Find the total energy of the motion. (b) Find the kinetic energy when the particle is $0.2$ m from the centre. (c) Find the positions where the kinetic energy is three times the potential energy.
\tcblower
We have $\omega = \dfrac{2\pi}{T} = \dfrac{2\pi}{\pi} = 2\ \mathrm{rad\,s^{-1}}$, $m = 0.2$ kg, $a = 0.4$ m.

\textbf{(a) Total energy:}
\[ E = \tfrac{1}{2}m\omega^{2}a^{2} = \tfrac{1}{2}(0.2)(4)(0.16) = 0.064\ \mathrm{J}. \]

\textbf{(b) Kinetic energy at $x = 0.2$ m:}
\[ E_{k} = \tfrac{1}{2}m\omega^{2}(a^{2} - x^{2}) = \tfrac{1}{2}(0.2)(4)(0.16 - 0.04) = 0.4 \times 0.12 = 0.048\ \mathrm{J}. \]

\textbf{(c) Positions where $E_{k} = 3E_{p}$:}
\[ \tfrac{1}{2}m\omega^{2}(a^{2} - x^{2}) = 3 \times \tfrac{1}{2}m\omega^{2}x^{2} \quad\Longrightarrow\quad a^{2} - x^{2} = 3x^{2} \quad\Longrightarrow\quad x^{2} = \frac{a^{2}}{4}. \]
Hence
\[ x = \pm\frac{a}{2} = \pm 0.2\ \mathrm{m}. \]
\textbf{Interpretation:} at half the amplitude, the energy is shared in the ratio $3:1$ between kinetic and potential energy. $\blacksquare$
\end{exoresolu}

\courssub{Skill 5: The Simple and Compound Pendulum}

\begin{methode}[Analysing pendulum motion]
\begin{enumerate}
\item For a \textbf{simple pendulum} of length $l$, write the tangential equation of motion for small angular displacement $\theta$:
\[ ml\ddot{\theta} = -mg\sin\theta \approx -mg\theta \quad\Longrightarrow\quad \ddot{\theta} = -\frac{g}{l}\theta, \qquad T = 2\pi\sqrt{\frac{l}{g}}. \]
\item For a \textbf{compound pendulum} (rigid body, axis through $O$), take moments about the axis:
\[ I_{O}\ddot{\theta} = -mgh\sin\theta \approx -mgh\,\theta \quad\Longrightarrow\quad T = 2\pi\sqrt{\frac{I_{O}}{mgh}}, \]
where $h$ is the distance from $O$ to the centre of mass and $I_{O}$ the moment of inertia about the axis (use the \textbf{parallel axes theorem} $I_{O} = I_{G} + mh^{2}$).
\item The small-angle approximation $\sin\theta \approx \theta$ (with $\theta$ in \textbf{radians}) is what makes the motion simple harmonic — state it explicitly.
\end{enumerate}
\textbf{Reflex:} in GCE questions, "show that the motion is approximately SHM" always signals the approximation $\sin\theta \approx \theta$.
\end{methode}

\begin{exoresolu}[A rod swinging as a compound pendulum]
A uniform rod $AB$ of length $2a$ and mass $m$ is free to rotate in a vertical plane about a smooth horizontal axis through the end $A$. The rod is displaced through a small angle and released from rest. Given that the moment of inertia of the rod about its centre is $\frac{1}{3}ma^{2}$, show that the rod performs approximate SHM and find the period. Take $g = 9.8\ \mathrm{m\,s^{-2}}$ and $a = 0.5$ m for the numerical value.
\tcblower
\textbf{Step 1: moment of inertia about $A$.} By the parallel axes theorem, with $h = a$ (the centre of mass $G$ is at the midpoint):
\[ I_{A} = I_{G} + mh^{2} = \tfrac{1}{3}ma^{2} + ma^{2} = \tfrac{4}{3}ma^{2}. \]

\textbf{Step 2: equation of motion.} Let $\theta$ be the angle the rod makes with the downward vertical. The weight $mg$ acts at $G$, with perpendicular distance $a\sin\theta$ from the axis, giving a restoring moment $-mga\sin\theta$. Taking moments about $A$:
\[ I_{A}\ddot{\theta} = -mga\sin\theta. \]
For small oscillations, $\sin\theta \approx \theta$, so
\[ \tfrac{4}{3}ma^{2}\,\ddot{\theta} = -mga\,\theta \quad\Longrightarrow\quad \ddot{\theta} = -\frac{3g}{4a}\,\theta. \]
This is SHM with $\omega^{2} = \dfrac{3g}{4a}$.

\textbf{Step 3: period.}
\[ T = 2\pi\sqrt{\frac{4a}{3g}}. \]
With $a = 0.5$ m and $g = 9.8\ \mathrm{m\,s^{-2}}$:
\[ T = 2\pi\sqrt{\frac{4 \times 0.5}{3 \times 9.8}} = 2\pi\sqrt{\frac{2}{29.4}} \approx 2\pi \times 0.2609 \approx 1.64\ \mathrm{s}. \]
\textbf{Remark:} this equals the period of a \emph{simple} pendulum of length $\dfrac{4a}{3}$ (the "equivalent simple pendulum length"). $\blacksquare$
\end{exoresolu}

\courssub{Skill 6: Solving Damped Harmonic Motion}

\begin{methode}[Analysing $\ddot{x} + 2k\dot{x} + \omega^{2}x = 0$]
\begin{enumerate}
\item Write the equation of motion including the resistance $-2mk\dot{x}$ (proportional to speed), and put it in the standard form
\[ \ddot{x} + 2k\dot{x} + \omega^{2}x = 0 . \]
\item Write the \textbf{auxiliary (characteristic) equation} $\lambda^{2} + 2k\lambda + \omega^{2} = 0$, with roots $\lambda = -k \pm \sqrt{k^{2} - \omega^{2}}$.
\item Classify the damping:
\begin{itemize}
\item $k < \omega$ (\textbf{light damping}): $x = e^{-kt}\left(A\cos\mu t + B\sin\mu t\right)$ with $\mu = \sqrt{\omega^{2} - k^{2}}$; oscillations with decaying amplitude, damped period $T_{d} = \dfrac{2\pi}{\mu}$.
\item $k = \omega$ (\textbf{critical damping}): $x = (A + Bt)e^{-kt}$; fastest return to equilibrium without oscillation.
\item $k > \omega$ (\textbf{heavy damping}): $x = Ae^{\lambda_{1}t} + Be^{\lambda_{2}t}$, both roots negative; slow return, no oscillation.
\end{itemize}
\item Use initial conditions to find $A$ and $B$ (remember to \textbf{differentiate} $x(t)$ before substituting $v(0)$).
\item The ratio of successive maximum displacements on the same side is $e^{-kT_{d}}$ (constant); its logarithm $kT_{d}$ is the \textbf{logarithmic decrement}.
\end{enumerate}
\end{methode}

\begin{popfigure}
\begin{center}
\begin{tikzpicture}
\begin{axis}[
  width=0.85\linewidth, height=5.2cm,
  axis lines=middle,
  xlabel={$t$}, ylabel={$x$},
  xmin=0, xmax=6.5, ymin=-2.6, ymax=2.6,
  xtick=\empty, ytick=\empty,
  samples=200, domain=0:6.3,
]
\addplot[popblue, thick] {2*exp(-0.5*x)*cos(deg(2.5*x))};
\addplot[poporange, dashed] {2*exp(-0.5*x)};
\addplot[poporange, dashed] {-2*exp(-0.5*x)};
\end{axis}
\end{tikzpicture}
\end{center}
\legende{Light damping: the displacement oscillates inside the decaying envelope $x = \pm a e^{-kt}$.}
\end{popfigure}

\begin{exoresolu}[Lightly damped oscillations]
A particle of mass $1$ kg moves on a line under a restoring force $25x$ N towards the origin and a resistance of magnitude $4v$ N, where $x$ is its displacement from $O$ and $v$ its speed. Initially $x = 2$ m and the particle is at rest. (a) Show that $\ddot{x} + 4\dot{x} + 25x = 0$. (b) Solve for $x$ in terms of $t$. (c) Find the damped period and the ratio of two successive maximum displacements on the same side of $O$.
\tcblower
\textbf{(a) Equation of motion.} Newton's second law with restoring force $-25x$ and resistance $-4\dot{x}$:
\[ \ddot{x} = -25x - 4\dot{x} \quad\Longrightarrow\quad \ddot{x} + 4\dot{x} + 25x = 0. \]

\textbf{(b) Solution.} Standard form: $2k = 4$, so $k = 2$; $\omega^{2} = 25$, so $\omega = 5$. Since $k = 2 < 5 = \omega$, the damping is \textbf{light}. The auxiliary equation is
\[ \lambda^{2} + 4\lambda + 25 = 0 \quad\Longrightarrow\quad \lambda = \frac{-4 \pm \sqrt{16 - 100}}{2} = -2 \pm \sqrt{21}\,i . \]
So $\mu = \sqrt{21}$ and
\[ x = e^{-2t}\left(A\cos\sqrt{21}\,t + B\sin\sqrt{21}\,t\right). \]
Initial conditions: $x(0) = 2 \Rightarrow A = 2$. Differentiate:
\[ \dot{x} = e^{-2t}\left(-\sqrt{21}A\sin\sqrt{21}t + \sqrt{21}B\cos\sqrt{21}t\right) - 2e^{-2t}\left(A\cos\sqrt{21}t + B\sin\sqrt{21}t\right), \]
\[ \dot{x}(0) = \sqrt{21}\,B - 2A = 0 \quad\Longrightarrow\quad B = \frac{2A}{\sqrt{21}} = \frac{4}{\sqrt{21}}. \]
Hence
\[ x = e^{-2t}\left(2\cos\sqrt{21}\,t + \frac{4}{\sqrt{21}}\sin\sqrt{21}\,t\right). \]

\textbf{(c) Damped period and amplitude ratio.}
\[ T_{d} = \frac{2\pi}{\mu} = \frac{2\pi}{\sqrt{21}} \approx 1.371\ \mathrm{s}. \]
Successive maxima on the same side are separated by one damped period, during which the envelope $e^{-2t}$ is multiplied by
\[ e^{-kT_{d}} = e^{-2 \times 1.371} = e^{-2.742} \approx 0.0645. \]
So each maximum is about $6.4\%$ of the previous one — the oscillations die out very quickly. $\blacksquare$
\end{exoresolu}

\courssub{Skill 7: Critical and Heavy Damping — Applications}

\begin{methode}[Choosing and interpreting the damping regime]
\begin{enumerate}
\item Compute $k$ and $\omega$ from the equation and compare them.
\item \textbf{Critical damping} ($k = \omega$) gives the quickest non-oscillatory return: this is the ideal for car suspensions, door closers and measuring instruments.
\item For $x = (A + Bt)e^{-kt}$, note that $x$ can vanish at most \textbf{once} (when $t = -A/B$): the particle may cross the equilibrium position at most one time.
\item For heavy damping, both exponents are negative: $x \to 0$ without oscillation, dominated at large $t$ by the \emph{slower} exponent (the one closer to $0$).
\end{enumerate}
\textbf{Reflex:} exam questions asking "does the particle pass through $O$?" are answered by solving $x(t) = 0$ and checking whether the root is positive.
\end{methode}

\begin{exoresolu}[Critically damped particle]
A particle of mass $0.5$ kg moves under a restoring force and a resistance such that its equation of motion is $\ddot{x} + 8\dot{x} + 16x = 0$. Initially the particle is at $O$ and is projected with speed $4\ \mathrm{m\,s^{-1}}$ away from $O$. (a) Solve for $x(t)$. (b) Find the maximum displacement from $O$ and the time at which it occurs. (c) State, with justification, whether the particle ever returns through $O$.
\tcblower
\textbf{(a) Solution.} Here $2k = 8$, so $k = 4$, and $\omega^{2} = 16$, so $\omega = 4$. Since $k = \omega$, the damping is \textbf{critical}. The auxiliary equation $\lambda^{2} + 8\lambda + 16 = 0$ has the repeated root $\lambda = -4$, so
\[ x = (A + Bt)e^{-4t}. \]
Initial conditions: $x(0) = 0 \Rightarrow A = 0$. Then
\[ \dot{x} = Be^{-4t} - 4Bt\,e^{-4t}, \qquad \dot{x}(0) = B = 4. \]
Hence
\[ x = 4t\,e^{-4t}. \]

\textbf{(b) Maximum displacement.} Set $\dot{x} = 0$:
\[ \dot{x} = 4e^{-4t}(1 - 4t) = 0 \quad\Longrightarrow\quad t = \frac{1}{4}\ \mathrm{s}. \]
Then
\[ x_{\max} = 4 \times \frac{1}{4} \times e^{-1} = e^{-1} \approx 0.368\ \mathrm{m}. \]

\textbf{(c) Return through $O$?} We need $x(t) = 0$ for $t > 0$, i.e. $4t\,e^{-4t} = 0$. Since $e^{-4t} > 0$ always, this requires $t = 0$ only. Therefore the particle \textbf{never} passes through $O$ again: it slows down, stops momentarily at $t = 0.25$ s, and creeps back towards $O$, reaching it only asymptotically as $t \to \infty$. $\blacksquare$
\end{exoresolu}

\courssub{Skill 8: Composite GCE Problems — Broken and Combined Oscillations}

\begin{methode}[Handling interrupted or two-stage oscillations]
GCE questions often combine several skills: a string going slack, an elastic string becoming taut, a collision, or a mass detaching mid-oscillation.
\begin{enumerate}
\item \textbf{Split the motion into phases} and identify the governing law in each phase (SHM while the string is taut; free fall or uniform motion otherwise).
\item At each \textbf{junction}, compute the position and velocity from the previous phase — they are the initial conditions of the next.
\item Check the \textbf{validity conditions}: an elastic string only pulls (it must stay taut: extension $\geq 0$); a spring can push and pull.
\item The total time of a "broken" oscillation is the sum of the times of the phases, found from $x = a\cos(\omega t + \varepsilon)$ in the SHM phases and from constant-acceleration formulas otherwise.
\end{enumerate}
\textbf{Reflex:} always ask: "is the string still taut / is the contact maintained?" before applying SHM formulas over a whole swing. The string goes slack during the motion exactly when the \textbf{amplitude exceeds the equilibrium extension} $e_{0}$.
\end{methode}

\begin{exoresolu}[Oscillation with a string going slack]
A particle $P$ of mass $0.4$ kg hangs in equilibrium from a light elastic string of natural length $0.8$ m and modulus $8$ N attached to a fixed point $O$. The particle is pulled down until the string has extension $1.0$ m and is then released from rest. Take $g = 10\ \mathrm{m\,s^{-2}}$. (a) Show that while the string is taut, $P$ performs SHM and find its period. (b) Show that the string becomes slack during the motion. (c) Find the time from release until the string first becomes slack, and the speed of $P$ at that instant.
\tcblower
\textbf{(a) SHM while taut.} Equilibrium extension $e_{0}$: $\dfrac{8e_{0}}{0.8} = 0.4 \times 10 \Rightarrow 10e_{0} = 4 \Rightarrow e_{0} = 0.4$ m. Let $x$ be the displacement below equilibrium; the tension is $T = \dfrac{8}{0.8}(0.4 + x) = 4 + 10x$. Newton's second law (downwards positive):
\[ 0.4\ddot{x} = 0.4 \times 10 - (4 + 10x) = -10x \quad\Longrightarrow\quad \ddot{x} = -25x. \]
SHM with $\omega = 5\ \mathrm{rad\,s^{-1}}$, period $T = \dfrac{2\pi}{5} \approx 1.257$ s.

\textbf{(b) Does the string go slack?} The particle is released from rest at extension $1.0$ m, i.e. at $x = 1.0 - 0.4 = 0.6$ m below equilibrium, so the amplitude is $a = 0.6$ m. The SHM would carry the particle up to $x = -0.6$ m, i.e. to extension $0.4 - 0.6 = -0.2$ m $< 0$, which is impossible for a string. Since $a = 0.6 > e_{0} = 0.4$, the string \textbf{does go slack} before the particle reaches the top of the would-be oscillation.

\textbf{(c) Time and speed when slack.} The string goes slack when the extension is $0$, i.e. at $x = -0.4$ m. With release from rest at $x = 0.6$, write $x = 0.6\cos(5t)$:
\[ 0.6\cos(5t) = -0.4 \quad\Longrightarrow\quad \cos(5t) = -\frac{2}{3} \quad\Longrightarrow\quad 5t = \arccos\left(-\tfrac{2}{3}\right) \approx 2.3005. \]
\[ t \approx \frac{2.3005}{5} \approx 0.460\ \mathrm{s}. \]
The speed there follows from $v^{2} = \omega^{2}(a^{2} - x^{2})$:
\[ v = 5\sqrt{0.6^{2} - 0.4^{2}} = 5\sqrt{0.36 - 0.16} = 5\sqrt{0.2} \approx 2.24\ \mathrm{m\,s^{-1}} \quad \text{(upwards).} \]
Thereafter the particle moves as a \textbf{free projectile} under gravity until the string tightens again — the starting point of the next phase. $\blacksquare$
\end{exoresolu}

\begin{attention}[Frequent errors in the GCE]
\begin{itemize}
\item Measuring $x$ from the \emph{top} or from the \emph{natural length} instead of from equilibrium: the equation then contains a constant term and candidates wrongly conclude the motion is not SHM. Shift the origin to equilibrium.
\item Forgetting that $\sin\theta \approx \theta$ requires $\theta$ in \textbf{radians}.
\item Confusing $\omega$ (angular frequency, $\mathrm{rad\,s^{-1}}$) with $f$ (frequency, Hz): $T = \dfrac{2\pi}{\omega} = \dfrac{1}{f}$.
\item In damped motion, substituting $v(0)$ into $x(t)$ instead of into $\dot{x}(t)$.
\item Applying SHM formulas after an elastic string has gone slack.
\end{itemize}
\end{attention}

\newpage

\coursec{Exercises}

\courssub{I check my knowledge}

\begin{exercice}[Multiple choice]
A particle moves with simple harmonic motion according to $x = a\cos(\omega t)$. Its acceleration is
\begin{enumerate}
\item $\ddot{x} = \omega^2 x$, directed away from the centre;
\item $\ddot{x} = -\omega^2 x$, directed towards the centre;
\item $\ddot{x} = -\omega x$, directed towards the centre;
\item $\ddot{x} = \omega^2 a\sin(\omega t)$, directed away from the centre.
\end{enumerate}
Choose the correct answer and justify your choice briefly.
\end{exercice}

\begin{exercice}[True or false]
State whether each of the following statements is true or false, justifying your answer.
\begin{enumerate}
\item In SHM, the speed of the particle is greatest at the extreme positions.
\item The period of a simple pendulum is independent of the mass of the bob.
\item In a critically damped motion, the system returns to equilibrium without oscillating.
\item Light damping reduces the period of an oscillator compared with the undamped period.
\end{enumerate}
\end{exercice}

\begin{exercice}[Definitions]
Define each of the following terms as used in this chapter:
\begin{enumerate}
\item simple harmonic motion;
\item the amplitude of an oscillation;
\item critical damping;
\item the damped (quasi-)period of a lightly damped oscillator.
\end{enumerate}
\end{exercice}

\begin{exercice}[Direct calculations]
A particle performs SHM with $\ddot{x} = -25x$, where $x$ is measured in metres.
\begin{enumerate}
\item Write down the angular frequency $\omega$ and the period $T$ of the motion.
\item Given that the amplitude is $0.4\ \mathrm{m}$, find the maximum speed and the maximum magnitude of the acceleration.
\item Find the speed of the particle when $x = 0.24\ \mathrm{m}$.
\end{enumerate}
\end{exercice}

\courssub{I practise}

\begin{exercice}[Mass on a horizontal spring]
A particle of mass $0.5\ \mathrm{kg}$ moves on a smooth horizontal table. It is attached to one end of a light elastic spring of natural length $0.8\ \mathrm{m}$ and modulus of elasticity $20\ \mathrm{N}$, the other end of the spring being fixed to a point of the table. The particle is pulled $0.3\ \mathrm{m}$ beyond its equilibrium position and released from rest.
\begin{enumerate}
\item Show that the particle performs simple harmonic motion and find the angular frequency.
\item Find the period of the oscillations.
\item Find the maximum speed of the particle.
\item Find the time taken to reach, for the first time, a point at half the amplitude from the centre.
\end{enumerate}
\end{exercice}

\begin{exercice}[Vertical spring]
A mass hangs in equilibrium from a light vertical spring, extending it by $10\ \mathrm{cm}$. The mass is pulled down a further $5\ \mathrm{cm}$ and then projected downward with speed $0.7\ \mathrm{m\,s^{-1}}$. Take $g = 9.8\ \mathrm{m\,s^{-2}}$.
\begin{enumerate}
\item Prove that the subsequent motion is simple harmonic and find its angular frequency.
\item Find the amplitude of the motion.
\item Find the speed of the mass when it is $2\ \mathrm{cm}$ above the equilibrium position.
\end{enumerate}
\end{exercice}

\begin{exercice}[Timing in SHM]
A particle performs simple harmonic motion of period $4\ \mathrm{s}$ and amplitude $0.6\ \mathrm{m}$.
\begin{enumerate}
\item Find the time it takes to travel directly from a point $0.3\ \mathrm{m}$ from the centre on one side to a point $0.3\ \mathrm{m}$ from the centre on the other side.
\item Find the speed of the particle at each of these two points.
\end{enumerate}
\end{exercice}

\begin{exercice}[Classifying damped motion]
The displacement $x$ of a damped oscillator satisfies
\[
\ddot{x} + 2k\dot{x} + 9x = 0,
\]
where $k$ is a positive constant.
\begin{enumerate}
\item Determine the values of $k$ for which the motion is (a) oscillatory (light damping), (b) critically damped, (c) heavily damped.
\item In the critical case, given $x(0) = 0.2$ and $\dot{x}(0) = 0$, find $x(t)$ and show that the particle never crosses the equilibrium position.
\end{enumerate}
\end{exercice}

\courssub{I prepare for the GCE}

\begin{exercice}[The seconds pendulum — GCE style] \hfill (12 marks)

A simple pendulum beats seconds (that is, its half-period is exactly $1\ \mathrm{s}$) in Yaoundé, where $g = 9.78\ \mathrm{m\,s^{-2}}$.
\begin{enumerate}
\item State the equation of motion of a simple pendulum for small oscillations and derive the formula for its period. \hfill (4 marks)
\item Calculate the length of this pendulum, giving your answer to four significant figures. \hfill (3 marks)
\item The pendulum is taken to a mountain station where its period increases by $0.2\%$. Find the new period and the value of $g$ at the mountain station. \hfill (5 marks)
\end{enumerate}
\end{exercice}

\begin{exercice}[A damped oscillator — GCE style] \hfill (14 marks)

A damped oscillator satisfies the equation
\[
\ddot{x} + 4\dot{x} + 13x = 0,
\]
with initial conditions $x(0) = 0.1$ and $\dot{x}(0) = 0$, where $x$ is measured in metres and $t$ in seconds.
\begin{enumerate}
\item Write down the auxiliary equation and solve it. \hfill (3 marks)
\item Hence find the general solution of the differential equation and apply the initial conditions to obtain $x(t)$ explicitly. \hfill (5 marks)
\item State the damped period of the oscillations. \hfill (2 marks)
\item Show that the ratio of two successive maximum displacements on the same side of the equilibrium position is constant, and find its value, giving your answer to three significant figures. \hfill (4 marks)
\end{enumerate}
\end{exercice}

\begin{exercice}[Floating log and car suspension — GCE style] \hfill (16 marks)

\textbf{Part A.} A cylindrical log of mass $40\ \mathrm{kg}$ and uniform cross-sectional area $0.05\ \mathrm{m^2}$ floats vertically in water of density $1000\ \mathrm{kg\,m^{-3}}$. The log is pushed down slightly from its equilibrium position and released. Take $g = 9.8\ \mathrm{m\,s^{-2}}$.
\begin{enumerate}
\item Using Archimedes' principle, find the depth immersed when the log is in equilibrium. \hfill (3 marks)
\item Show that, neglecting resistance, the log performs simple harmonic motion, and find the period of the oscillations. \hfill (5 marks)
\item Discuss qualitatively how the resistance of the water would modify the motion, referring to the terms of a suitable differential equation. \hfill (3 marks)
\end{enumerate}

\textbf{Part B.} A car's suspension is modelled by a mass of $250\ \mathrm{kg}$, a spring of stiffness $18\,000\ \mathrm{N\,m^{-1}}$ and a damper providing a resistance $\lambda\dot{x}$.
\begin{enumerate}
\item[(d)] The suspension must return to equilibrium as fast as possible without oscillating. Write down the equation of motion and find the required value of $\lambda$. \hfill (3 marks)
\item[(e)] If $\lambda$ is only $60\%$ of this value, classify the motion and find the damped period. \hfill (2 marks)
\end{enumerate}
\end{exercice}

\newpage

\coursec{Solutions to the exercises}

\coursec{Worked Exercises: Simple and Damped Harmonic Motion}

\begin{corrige}[Exercise 1 — Multiple Choice]
The correct answer is \textbf{(b)}: $\ddot{x} = -\omega^2 x$, directed towards the centre.

\textbf{Justification.} Differentiating $x = a\cos(\omega t)$ twice:
\[
\dot{x} = -a\omega\sin(\omega t), \qquad \ddot{x} = -a\omega^2\cos(\omega t) = -\omega^2 x.
\]
The acceleration is therefore proportional to the displacement and has the \emph{opposite} sign: when $x > 0$ the acceleration is negative, and conversely. Hence it is always directed towards the centre $O$. This is precisely the defining property of simple harmonic motion.
\end{corrige}

\begin{corrige}[Exercise 2 — True or False]
\begin{enumerate}
\item \textbf{False.} From $v^2 = \omega^2(a^2 - x^2)$, the speed is greatest when $x = 0$, i.e.\ at the \emph{centre}; at the extreme positions ($x = \pm a$) the speed is zero.
\item \textbf{True.} For a simple pendulum, $T = 2\pi\sqrt{\ell/g}$: the mass $m$ cancels when the equation of motion $m\ell\ddot{\theta} = -mg\sin\theta$ is formed, so the period depends only on $\ell$ and $g$.
\item \textbf{True.} Critical damping ($k = \omega$ in $\ddot{x} + 2k\dot{x} + \omega^2 x = 0$) gives $x = (A + Bt)e^{-kt}$, which has at most one zero; generically the system returns to equilibrium without oscillating, in the shortest possible time.
\item \textbf{False.} The damped angular frequency is $\omega_d = \sqrt{\omega^2 - k^2} < \omega$, so the damped period $T_d = 2\pi/\omega_d$ is \emph{longer} than the undamped period $2\pi/\omega$. Light damping \emph{increases} the period (slightly).
\end{enumerate}
\end{corrige}

\begin{corrige}[Exercise 3 — Definitions]
\begin{enumerate}
\item \textbf{Simple harmonic motion} is the motion of a particle along a straight line whose acceleration is always directed towards a fixed point $O$ of the line (the centre) and is proportional to its displacement from $O$: $\ddot{x} = -\omega^2 x$.
\item The \textbf{amplitude} is the maximum distance of the particle from the centre of the motion, i.e.\ half the distance between the two extreme positions.
\item \textbf{Critical damping} is the case of the equation $\ddot{x} + 2k\dot{x} + \omega^2 x = 0$ in which $k = \omega$ (repeated auxiliary root); the system returns to equilibrium in the shortest time without oscillating.
\item The \textbf{damped (quasi-)period} of a lightly damped oscillator ($k < \omega$) is the time between two successive passages through the equilibrium position in the same direction: $T_d = \dfrac{2\pi}{\sqrt{\omega^2 - k^2}}$.
\end{enumerate}
\end{corrige}

\begin{corrige}[Exercise 4 — Direct Calculations]
\begin{enumerate}
\item Comparing $\ddot{x} = -25x$ with $\ddot{x} = -\omega^2 x$: $\omega^2 = 25$, so $\omega = 5\ \mathrm{rad\,s^{-1}}$ and
\[
T = \frac{2\pi}{\omega} = \frac{2\pi}{5} \approx 1.26\ \mathrm{s}.
\]
\item With $a = 0.4\ \mathrm{m}$:
\[
v_{\max} = a\omega = 0.4 \times 5 = 2\ \mathrm{m\,s^{-1}}, \qquad |\ddot{x}|_{\max} = \omega^2 a = 25 \times 0.4 = 10\ \mathrm{m\,s^{-2}}.
\]
\item Using $v^2 = \omega^2(a^2 - x^2)$ with $x = 0.24\ \mathrm{m}$:
\[
v^2 = 25\left(0.4^2 - 0.24^2\right) = 25(0.16 - 0.0576) = 25 \times 0.1024 = 2.56,
\]
so $v = 1.6\ \mathrm{m\,s^{-1}}$.
\end{enumerate}
\end{corrige}

\begin{corrige}[Exercise 5 — Mass on a Horizontal Spring]
\begin{enumerate}
\item Let $x$ be the displacement from the equilibrium position (which, on a smooth horizontal table, is the position where the spring has its natural length). The tension in the spring is $T = \dfrac{\lambda x}{l_0} = \dfrac{20}{0.8}x = 25x$, directed towards equilibrium. Newton's second law gives
\[
0.5\,\ddot{x} = -25x \quad\Longrightarrow\quad \ddot{x} = -50x.
\]
This is of the form $\ddot{x} = -\omega^2 x$ with $\omega^2 = 50$: the motion is SHM with $\omega = \sqrt{50} = 5\sqrt{2} \approx 7.07\ \mathrm{rad\,s^{-1}}$.
\item $T = \dfrac{2\pi}{\omega} = \dfrac{2\pi}{5\sqrt{2}} \approx 0.889\ \mathrm{s}$.
\item The particle is released from rest at $0.3\ \mathrm{m}$ from the centre, so $a = 0.3\ \mathrm{m}$:
\[
v_{\max} = a\omega = 0.3 \times 5\sqrt{2} = 1.5\sqrt{2} \approx 2.12\ \mathrm{m\,s^{-1}}.
\]
\item Measure time from release: $x = a\cos(\omega t) = 0.3\cos(5\sqrt{2}\,t)$. Half the amplitude: $x = 0.15$:
\[
\cos(5\sqrt{2}\,t) = \tfrac{1}{2} \;\Longrightarrow\; 5\sqrt{2}\,t = \frac{\pi}{3} \;\Longrightarrow\; t = \frac{\pi}{15\sqrt{2}} \approx 0.148\ \mathrm{s}.
\]
\end{enumerate}
\end{corrige}

\begin{corrige}[Exercise 6 — Vertical Spring]
\begin{enumerate}
\item At equilibrium, the spring is stretched by $e = 0.1\ \mathrm{m}$, so $mg = \dfrac{\lambda e}{l_0}$. Let $x$ be the displacement measured \emph{downward from equilibrium}. The total extension is $e + x$, so the upward tension is $\dfrac{\lambda(e+x)}{l_0}$. Newton's second law (downward positive):
\[
m\ddot{x} = mg - \frac{\lambda(e+x)}{l_0} = -\frac{\lambda}{l_0}x,
\]
since $mg = \lambda e/l_0$. Hence $\ddot{x} = -\dfrac{g}{e}x = -\dfrac{9.8}{0.1}x = -98x$: SHM with $\omega = \sqrt{98} = 7\sqrt{2} \approx 9.90\ \mathrm{rad\,s^{-1}}$.
\item The mass starts at $x_0 = 0.05\ \mathrm{m}$ with $\dot{x}_0 = 0.7\ \mathrm{m\,s^{-1}}$. Using $v^2 = \omega^2(a^2 - x^2)$:
\[
a^2 = x_0^2 + \frac{v_0^2}{\omega^2} = 0.05^2 + \frac{0.7^2}{98} = 0.0025 + 0.005 = 0.0075,
\]
so $a = \sqrt{0.0075} \approx 0.0866\ \mathrm{m} = 8.66\ \mathrm{cm}$.
\item At $x = -0.02\ \mathrm{m}$ (above equilibrium):
\[
v^2 = 98\left(0.0075 - 0.0004\right) = 98 \times 0.0071 = 0.6958,
\]
so $v \approx 0.834\ \mathrm{m\,s^{-1}}$.
\end{enumerate}
\end{corrige}

\begin{corrige}[Exercise 7 — Timing in SHM]
Here $T = 4\ \mathrm{s}$, so $\omega = \dfrac{2\pi}{4} = \dfrac{\pi}{2}\ \mathrm{rad\,s^{-1}}$, and $a = 0.6\ \mathrm{m}$.
\begin{enumerate}
\item Take $x = 0.6\cos\left(\tfrac{\pi}{2}t\right)$. The particle is at $x = +0.3$ when $\cos\left(\tfrac{\pi}{2}t\right) = \tfrac12$, i.e.\ $\tfrac{\pi}{2}t_1 = \tfrac{\pi}{3}$, $t_1 = \tfrac{2}{3}\ \mathrm{s}$; it reaches $x = -0.3$ when $\tfrac{\pi}{2}t_2 = \tfrac{2\pi}{3}$, $t_2 = \tfrac{4}{3}\ \mathrm{s}$. The direct journey takes
\[
t_2 - t_1 = \frac{2}{3}\ \mathrm{s} \approx 0.667\ \mathrm{s}.
\]
\item At $x = \pm 0.3\ \mathrm{m}$:
\[
v^2 = \omega^2(a^2 - x^2) = \frac{\pi^2}{4}(0.36 - 0.09) = \frac{\pi^2}{4}\times 0.27,
\qquad v = \frac{\pi}{2}\sqrt{0.27} \approx 0.817\ \mathrm{m\,s^{-1}}.
\]
The speed is the same at both points, as expected by symmetry.
\end{enumerate}
\end{corrige}

\begin{corrige}[Exercise 8 — Classifying Damped Motion]
\begin{enumerate}
\item The auxiliary equation is $m^2 + 2km + 9 = 0$, with discriminant $\Delta = 4k^2 - 36 = 4(k^2 - 9)$.
\begin{itemize}
\item (a) Oscillatory (light damping): $\Delta < 0 \iff 0 < k < 3$;
\item (b) Critical damping: $\Delta = 0 \iff k = 3$;
\item (c) Heavy damping: $\Delta > 0 \iff k > 3$.
\end{itemize}
\item For $k = 3$: $\ddot{x} + 6\dot{x} + 9x = 0$, auxiliary equation $(m+3)^2 = 0$, so
\[
x(t) = (A + Bt)e^{-3t}.
\]
$x(0) = A = 0.2$; $\dot{x} = (B - 3A - 3Bt)e^{-3t}$, so $\dot{x}(0) = B - 3A = 0$, giving $B = 0.6$. Hence
\[
x(t) = (0.2 + 0.6t)e^{-3t}.
\]
Since $0.2 + 0.6t > 0$ for all $t \geq 0$ and $e^{-3t} > 0$, we have $x(t) > 0$ for all $t \geq 0$: the particle approaches the equilibrium position asymptotically but never crosses it.
\end{enumerate}
\end{corrige}

\begin{corrige}[Exercise 9 — The Seconds Pendulum (GCE Style)]
\textbf{Indicative mark scheme:} (a) B1 equation of motion, M1 small-angle approximation, A1 SHM form, A1 period. (b) M1 rearranging, M1 substitution, A1 answer to 4 s.f. (c) B1 new period, M1 proportionality $T \propto 1/\sqrt{g}$, M1 substitution, A1 value of $g$, B1 units/comment.

\begin{enumerate}
\item For a simple pendulum of length $\ell$, resolving along the tangent: $m\ell\ddot{\theta} = -mg\sin\theta$. For small oscillations, $\sin\theta \approx \theta$, giving
\[
\ddot{\theta} = -\frac{g}{\ell}\theta,
\]
which is SHM with $\omega^2 = g/\ell$, hence $T = 2\pi\sqrt{\ell/g}$.
\item A seconds pendulum has half-period $1\ \mathrm{s}$, so $T = 2\ \mathrm{s}$:
\[
\ell = \frac{gT^2}{4\pi^2} = \frac{9.78 \times 4}{4\pi^2} = \frac{9.78}{\pi^2} \approx 0.9910\ \mathrm{m}.
\]
\item New period: $T' = 2 \times 1.002 = 2.004\ \mathrm{s}$. Since $T = 2\pi\sqrt{\ell/g}$ with $\ell$ fixed, $g \propto 1/T^2$:
\[
g' = g\left(\frac{T}{T'}\right)^2 = 9.78 \times \left(\frac{1}{1.002}\right)^2 \approx 9.78 \times 0.99602 \approx 9.741\ \mathrm{m\,s^{-2}}.
\]
(The decrease of $g$ with altitude explains the increase of the period.)
\end{enumerate}
\textbf{Examiner's expectations:} the small-angle approximation must be stated explicitly; in (c), using the ratio method avoids recomputing $\ell$ and is the expected elegant route.
\end{corrige}

\begin{corrige}[Exercise 10 — A Damped Oscillator (GCE Style)]
\textbf{Indicative mark scheme:} (a) B1 auxiliary equation, M1 solving, A1 complex roots. (b) B1 form of general solution, M1/M1 applying each initial condition, A1 constants, A1 final $x(t)$. (c) M1 use of $\omega_d$, A1 period. (d) M1 envelope argument, M1 ratio over one period, A1 exponential value, A1 3 s.f.

\begin{enumerate}
\item Auxiliary equation: $m^2 + 4m + 13 = 0$;
\[
m = \frac{-4 \pm \sqrt{16 - 52}}{2} = -2 \pm 3i.
\]
\item Since the roots are complex ($-2 \pm 3i$), the motion is lightly damped and
\[
x(t) = e^{-2t}\left(A\cos 3t + B\sin 3t\right).
\]
$x(0) = A = 0.1$. Differentiating:
\[
\dot{x} = e^{-2t}\left[(-2A + 3B)\cos 3t + (-3A - 2B)\sin 3t\right],
\]
so $\dot{x}(0) = -2A + 3B = 0$, giving $B = \dfrac{0.2}{3} = \tfrac{1}{15} \approx 0.0667$. Hence
\[
x(t) = e^{-2t}\left(0.1\cos 3t + \tfrac{1}{15}\sin 3t\right)\ \mathrm{(m)}.
\]
\item Damped period: $T_d = \dfrac{2\pi}{3} \approx 2.09\ \mathrm{s}$.
\item Successive maxima on the same side are separated by one damped period $T_d$. Since the oscillatory factor repeats after $T_d$ while the envelope is multiplied by $e^{-2T_d}$, the ratio of successive maximum displacements is the constant
\[
\frac{x_{n+1}}{x_n} = e^{-2T_d} = e^{-4\pi/3} \approx 0.0152.
\]
(Equivalently, the logarithmic decrement per cycle is $\dfrac{4\pi}{3} \approx 4.19$.)
\end{enumerate}
\textbf{Examiner's expectations:} candidates must justify that maxima recur with period $T_d$ (the cosine factor is periodic) and that the exponential envelope alone changes; quoting $e^{-kT_d}$ without justification earns only partial credit.
\end{corrige}

\begin{corrige}[Exercise 11 — Floating Log and Car Suspension (GCE Style)]
\textbf{Indicative mark scheme (Part A):} (a) M1 Archimedes, M1 equilibrium equation, A1 depth. (b) M1 net force with extra immersion, A1 SHM equation, M1 identifying $\omega$, M1 period formula, A1 value. (c) B1 damping term, B1 modified equation, B1 qualitative effect.

\textbf{Part A.}
\begin{enumerate}
\item At equilibrium, weight = upthrust. If $d$ is the immersed depth, the displaced volume is $0.05d$, so
\[
40g = 1000 \times 0.05\,d \times g \quad\Longrightarrow\quad d = \frac{40}{50} = 0.8\ \mathrm{m}.
\]
\item Let $x$ be the additional downward displacement from equilibrium. The extra upthrust is $1000 \times 0.05\,x \times g = 490x$ (upward). Newton's second law:
\[
40\ddot{x} = -490x \quad\Longrightarrow\quad \ddot{x} = -12.25x,
\]
which is SHM with $\omega = \sqrt{12.25} = 3.5\ \mathrm{rad\,s^{-1}}$. The period is
\[
T = \frac{2\pi}{3.5} \approx 1.80\ \mathrm{s}.
\]
\item Water resistance opposes the velocity, adding a term $-\lambda\dot{x}$ to the resultant force:
\[
40\ddot{x} + \lambda\dot{x} + 490x = 0.
\]
The oscillations then decay exponentially in amplitude (light damping, quasi-period slightly increased), or the log settles without oscillating if the damping is critical or heavy.
\end{enumerate}

\textbf{Part B.} Indicative marks: (d) B1 equation, M1 critical condition, A1 value of $\lambda$; (e) B1 classification, A1 damped period.
\begin{enumerate}
\item[(d)] The equation of motion is
\[
250\ddot{x} + \lambda\dot{x} + 18\,000x = 0.
\]
Fastest return without oscillation requires \emph{critical damping}: auxiliary equation $250m^2 + \lambda m + 18\,000 = 0$ with zero discriminant:
\[
\lambda^2 = 4 \times 250 \times 18\,000 = 1.8\times 10^{7}
\quad\Longrightarrow\quad \lambda = \sqrt{1.8\times10^{7}} \approx 4243\ \mathrm{N\,s\,m^{-1}}.
\]
\item[(e)] With $\lambda = 0.6 \times 4243 \approx 2546\ \mathrm{N\,s\,m^{-1}}$, the discriminant is negative: the motion is \textbf{lightly damped (oscillatory)}. Writing $\ddot{x} + 2k\dot{x} + \omega^2 x = 0$: $\omega^2 = \dfrac{18\,000}{250} = 72$ and $2k = \dfrac{2546}{250}$, so $k \approx 5.09$. Then
\[
\omega_d = \sqrt{\omega^2 - k^2} = \sqrt{72 - 25.9} \approx 6.79\ \mathrm{rad\,s^{-1}},
\qquad T_d = \frac{2\pi}{6.79} \approx 0.925\ \mathrm{s}.
\]
\end{enumerate}
\textbf{Examiner's expectations:} in Part A the upthrust must be built from Archimedes' principle, not quoted; in Part B the critical-damping condition $\lambda = 2\sqrt{mk}$ (equivalently $\lambda^2 = 4 \times 250 \times 18\,000$) must be justified from the discriminant.
\end{corrige}

\newpage

\coursec{Summary sheet}

\begin{popfigure}
\begin{tikzpicture}[mindmap, grow cyclic, every node/.style={concept, text=white, font=\small},
  root concept/.append style={concept color=popdark, font=\bfseries},
  level 1/.append style={sibling angle=60, level distance=3.4cm},
  level 2/.append style={sibling angle=45, level distance=2.2cm, font=\scriptsize}]
\node[root concept]{SHM \& Damped Harmonic Motion}
  child[concept color=popblue]{ node{Definition \& Equation}
    child[concept color=popblue]{ node{$\ddot{x}=-\omega^2 x$} }
    child[concept color=popblue]{ node{$x=a\cos(\omega t+\phi)$} }
  }
  child[concept color=poporange]{ node{Kinematics}
    child[concept color=poporange]{ node{$v=\pm\omega\sqrt{a^2-x^2}$} }
    child[concept color=poporange]{ node{$T=\dfrac{2\pi}{\omega}$} }
  }
  child[concept color=popgreen]{ node{Energy}
    child[concept color=popgreen]{ node{$E=\tfrac12 m\omega^2 a^2$} }
    child[concept color=popgreen]{ node{KE $\leftrightarrow$ PE} }
  }
  child[concept color=popgold]{ node{Pendulums}
    child[concept color=popgold]{ node{$T=2\pi\sqrt{\ell/g}$} }
    child[concept color=popgold]{ node{Spring: $T=2\pi\sqrt{m/k}$} }
  }
  child[concept color=poppurple]{ node{Damped Motion}
    child[concept color=poppurple]{ node{$m\ddot{x}+c\dot{x}+kx=0$} }
    child[concept color=poppurple]{ node{Light, critical, heavy} }
  }
  child[concept color=poppink]{ node{Applications \& GCE}
    child[concept color=poppink]{ node{Modelling} }
    child[concept color=poppink]{ node{Exam technique} }
  };
\end{tikzpicture}
\legende{Mind map of the chapter: Simple and Damped Harmonic Motion}
\end{popfigure}

\begin{bilan}
\textbf{Key definitions}
\begin{itemize}
\item \cle{Simple harmonic motion} (SHM): motion about a fixed point $O$ in which the acceleration is directed towards $O$ and is proportional to the displacement from $O$: $\ddot{x}=-\omega^2 x$.
\item \cle{Amplitude} $a$: maximum displacement from the centre; \cle{period} $T=\dfrac{2\pi}{\omega}$; \cle{frequency} $f=\dfrac{1}{T}$; $\omega$ is the angular frequency.
\item \cle{Damped harmonic motion}: oscillatory motion with a resisting force proportional to velocity, $-c\dot{x}$, causing the amplitude to decay.
\item \cle{Light (under) damping}: oscillations with exponentially decaying amplitude; \cle{critical damping}: fastest return to equilibrium without oscillation; \cle{heavy (over) damping}: slow return without oscillation.
\end{itemize}

\textbf{Formulas and results to know}
\begin{itemize}
\item SHM equation: $\ddot{x}=-\omega^2 x$, with general solution $x=a\cos(\omega t+\phi)$, or $x=A\cos\omega t+B\sin\omega t$.
\item Velocity: $\dot{x}=\pm\omega\sqrt{a^2-x^2}$; maximum speed $v_{\max}=\omega a$ at the centre; maximum acceleration $\omega^2 a$ at the extremes.
\item Displacement from equilibrium: $x=a\cos\omega t$ (starting at extreme) or $x=a\sin\omega t$ (starting at centre).
\item Energy: kinetic $\tfrac12 m\omega^2(a^2-x^2)$, potential $\tfrac12 m\omega^2 x^2$, total $E=\tfrac12 m\omega^2 a^2$ (constant).
\item Simple pendulum (small oscillations): $T=2\pi\sqrt{\dfrac{\ell}{g}}$.
\item Mass on a spring: $T=2\pi\sqrt{\dfrac{m}{k}}$; for a vertical spring, oscillations are about the \emph{equilibrium} position, and gravity only shifts the centre.
\item Damping equation: $m\ddot{x}+c\dot{x}+kx=0$; auxiliary equation $m\lambda^2+c\lambda+k=0$:
  \begin{itemize}
  \item $c^2<4mk$: complex roots $\Rightarrow$ light damping, $x=e^{-\alpha t}(A\cos\beta t+B\sin\beta t)$;
  \item $c^2=4mk$: repeated root $\Rightarrow$ critical damping, $x=(A+Bt)e^{-\alpha t}$;
  \item $c^2>4mk$: real roots $\Rightarrow$ heavy damping, $x=Ae^{\lambda_1 t}+Be^{\lambda_2 t}$.
  \end{itemize}
\end{itemize}

\textbf{Methods}
\begin{itemize}
\item To prove SHM: choose a displacement $x$ from equilibrium, apply $F=ma$, and reduce to $\ddot{x}=-\omega^2 x$; read off $\omega$ and $T$.
\item To find amplitude and phase: use the initial conditions $x(0)$ and $\dot{x}(0)$.
\item Time between two positions: solve $x(t)=x_0$ using the appropriate trig form, choosing the correct root from the direction of motion.
\item For damped motion: form the auxiliary equation, classify the damping from the discriminant, then fit constants to the initial conditions.
\end{itemize}

\textbf{Common mistakes}
\begin{itemize}
\item Mixing up $x=a\cos\omega t$ and $x=a\sin\omega t$; always check against $x(0)$.
\item Forgetting that $v=\pm\omega\sqrt{a^2-x^2}$ has two signs; the direction of travel selects the sign.
\item Using degrees instead of radians when evaluating $\omega t+\phi$.
\item For a vertical spring, measuring $x$ from the unstretched position instead of equilibrium, and wrongly concluding the motion is not SHM.
\item Writing the damping force with the wrong sign, or classifying damping using $c^2$ against $4k$ instead of $4mk$.
\item Confusing the damped angular frequency $\beta$ with the natural frequency $\omega_0=\sqrt{k/m}$.
\end{itemize}

\textbf{GCE tips}
\begin{itemize}
\item Examiners award a method mark for explicitly reaching the form $\ddot{x}=-\omega^2 x$ before quoting $T=\dfrac{2\pi}{\omega}$.
\item State the centre of oscillation clearly; many marks are lost through a poorly defined variable.
\item In ``show that'' questions, keep every step of the algebra; in numerical parts, give answers to 3 significant figures with units.
\item For pendulum questions, justify the small-angle approximation $\sin\theta\approx\theta$ (in radians).
\item Sketch displacement--time graphs for the three damping regimes: examiners reward correct shape, decay envelope $e^{-\alpha t}$, and correct intercepts.
\end{itemize}
\end{bilan}

\findecours

\end{document}
